\documentclass[11pt,a4paper,openany]{book}

% BEGIN preamble
\usepackage[T1]{fontenc}
\usepackage{lmodern,amsmath,amssymb,amsthm,mathtools}
\usepackage[a4paper,margin=25mm,headheight=15pt]{geometry}
\usepackage{microtype,enumitem,xcolor,longtable,booktabs,array,needspace}
\usepackage[colorlinks=true,linkcolor=blue!40!black,urlcolor=blue!40!black]{hyperref}
\usepackage[nameinlink,noabbrev]{cleveref}
\usepackage{fancyhdr}
\hypersetup{pdftitle={Complex Analysis: Unified Notes},pdfauthor={MathGenome preparation series}}
\pagestyle{fancy}\fancyhf{}
\fancyhead[L]{\small MathGenome: Complex Analysis}\fancyhead[R]{\small\thepage}
\renewcommand{\headrulewidth}{0.3pt}
\setlength{\parindent}{0pt}\setlength{\parskip}{5pt}
\setlength{\emergencystretch}{3em}\allowdisplaybreaks[2]\raggedbottom
\setcounter{tocdepth}{1}
\newtheorem{theorem}{Theorem}[chapter]
\newtheorem{lemma}[theorem]{Lemma}
\newtheorem{proposition}[theorem]{Proposition}
\newtheorem{corollary}[theorem]{Corollary}
\theoremstyle{definition}
\newtheorem{definition}[theorem]{Definition}
\newtheorem{construction}[theorem]{Construction}
\newtheorem{example}[theorem]{Example}
\theoremstyle{remark}\newtheorem{remark}[theorem]{Remark}
\crefname{construction}{construction}{constructions}
\newcommand{\R}{\mathbb R}\newcommand{\C}{\mathbb C}
\newcommand{\N}{\mathbb N}\newcommand{\Z}{\mathbb Z}\newcommand{\Q}{\mathbb Q}
\newcommand{\D}{\mathbb D}\newcommand{\Chat}{\widehat{\mathbb C}}
\newcommand{\eps}{\varepsilon}\newcommand{\ee}{\varepsilon}
\newcommand{\abs}[1]{\lvert#1\rvert}\newcommand{\norm}[1]{\lVert#1\rVert}
\newcommand{\dd}{\,d}\newcommand{\iu}{\mathrm i}
\DeclareMathOperator{\Hol}{Hol}\DeclareMathOperator{\Mer}{Mer}
\DeclareMathOperator{\Res}{Res}\DeclareMathOperator{\ind}{ind}
\DeclareMathOperator{\Ree}{Re}\DeclareMathOperator{\Imm}{Im}
\DeclareMathOperator{\supp}{supp}\DeclareMathOperator{\diam}{diam}
\DeclareMathOperator{\dist}{dist}\DeclareMathOperator{\ord}{ord}
\newcommand{\family}[1]{}
\newcommand{\source}[1]{\par\begingroup\small\sffamily\color{black!65}Source comparison: #1\par\endgroup}

% END preamble

\begin{document}
\frontmatter
% BEGIN frontmatter
\begin{titlepage}\centering
{\Large MathGenome Preparation Series\par}\vspace{25mm}
{\Huge\bfseries Complex Analysis\par}\vspace{8mm}
{\LARGE Unified Notes\par}\vspace{12mm}
{\large Standard and Advanced Undergraduate One-Variable Theory\par}
\vfill{English mathematical manuscript\par November 2026}
\end{titlepage}
\chapter*{Scope and reading contract}
These notes combine one-variable complex analysis into a single notation and
an explicit proof dependency structure. The source comparisons use Conway,
Ahlfors, Stein--Shakarchi, Freitag--Busam, and selected introductory portions of
Forster. This is an original synthesis, not a transcription or a claim to cover
every result in all five books.

The adopted scope includes local and global Cauchy theory, local series,
singularities, residues, zero counting, conformal maps, normal families, the
planar Riemann mapping theorem, harmonic and Perron theory, Runge approximation,
Mittag--Leffler and Weierstrass constructions, growth and factorization,
analytic continuation, introductory surfaces and algebraic branches, elliptic
functions and their concrete projective cubic, and a proved route to Picard.
The stable source-family identifiers are independent of the printed chapter
numbers. They are content identifiers, not prescribed Lean target counts.

The application branches on Gamma/zeta and complex differential equations,
general Jordan boundary extension and prime ends, general surface
uniformization, Riemann--Roch, duality, and full modular-form theory are outside
this manuscript's adopted proof scope. A conditional Schwarz--Christoffel
formula does not assert accessory-parameter existence. A lattice-associated
cubic does not assert general algebraization or a converse uniformization.

Every adopted principal conclusion must retain its precise hypotheses and
conclusion. Equivalent complete proofs and legitimate previously established
foundations are allowed. A particular textbook proof's unused intermediate
steps are not separate deliverables. Definitions and representation bridges
that are actually used still need justification.

\paragraph{Preparation status.}
The adopted mathematical scope has completed author checking and independent
cross-review of every chapter and result family, with the recorded findings
resolved. The verification register describes the evidence and its limits.
Source citations identify comparison material; a citation alone does not close
a proof. This is a reviewed mathematical input, not yet a Lean formalization
or a kernel certificate. No D12 formalization run is started by this manuscript.

% END frontmatter
\tableofcontents
\mainmatter
% BEGIN foundations
\chapter*{Prior foundations and dependency boundaries}
\addcontentsline{toc}{chapter}{Prior foundations and dependency boundaries}
The logical starting point is real and complex arithmetic, real analysis,
finite-dimensional real differential and integral calculus, and point-set
topology. The following prior foundations are allowed only with the stated
scope; they do not include any of the complex-analysis conclusions proved in
later chapters.

\begin{enumerate}[label=F\arabic*.]
\item Real/complex normed-field arithmetic, completeness, compactness of closed
bounded subsets of finite-dimensional real space, and elementary polynomial
algebra including the real intermediate value theorem. The complex fundamental
theorem of algebra is proved here rather than assumed.
\item Series, absolute and compact uniform convergence, the Cauchy criterion,
real one-dimensional differentiation and integration, the real FTC and
integration by parts. Complex-valued integration is the real and imaginary
coordinate integral, with the standard absolute integral estimate.
\item Real total derivatives, chain rule, the real finite-dimensional inverse
and implicit function theorems, and change of variables for injective C1 maps
with the appropriate domain and integrability hypotheses. The complex-linear
differential and its real Jacobian bridge are proved in the manuscript.
The real planar Green circulation/flux formulas are available for an already
specified bounded C1 region, or a finite piecewise regular decomposition with
the stated oriented-boundary cancellation hypotheses, also after removing
interior disks. The fields must be C1 on a neighborhood of the closure;
the first and second Green identities for scalar functions require C2 there.
These are the prior D04 results labelled \texttt{d04:12:green},
\texttt{d04:12:parts}, and \texttt{d04:15:piecewise}, with checked Lean interfaces
\texttt{D04.Green.riemann\_manifold\_green} and
\texttt{D04.GreenIdentities.second\_green}. They assert no interior-existence
or separation theorem for an arbitrary Jordan trace.
\item Metric and general topology: continuity, compactness, connectedness,
countable exhaustion and diagonal subsequences, the real Arzela--Ascoli theorem,
and quotient topology. Any plane separation, winding-number, path-lifting or
surface assertion needed beyond these is proved before use.
\item Lebesgue dominated and monotone convergence, Fubini under its hypotheses,
scalar measure integration, absolute continuity/BV and Stieltjes integrals,
and the real area/change-of-variables theorem. Contour representations and
the holomorphic area factor are separately justified.
\end{enumerate}

MathGenome supplies the already checked D03, D04, D10 and D11 background in a
pinned Lean/Mathlib environment. Its exact interface crosswalk is retained in
the accompanying preparation records. Those libraries are possible later Lean
reuse sources, not a justification to assume an absent complex theorem.
Future directions D22, D23 and D24 are ownership references only; they supply
no unproved theorem to this manuscript. A forward cross-reference is permitted
only when the actual proof dependency remains acyclic.

% END foundations


% BEGIN chapters/01
\chapter{Holomorphicity and real differentiation}\label{ch:d12-01}

Write \(B(a,r)=\{z\in\C:|z-a|<r\}\). Throughout, a domain is a
nonempty connected open subset of \(\C\). Arithmetic, real calculus and
topology are used only within the foundations declared at the beginning.

\begin{definition}[Complex limits and derivatives]\family{D12.01.01}\label{d12:01-01}
For a function defined near \(a\), the assertion \(f(z)\to L\) means:
for every \(\eps>0\), some \(\delta>0\) makes
\(0<|z-a|<\delta\) imply \(|f(z)-L|<\eps\).
Continuity includes the value \(f(a)\). Its complex derivative is
\[
 f'(a)=\lim_{h\to0,\ h\ne0}\frac{f(a+h)-f(a)}h .
\]
A function is holomorphic on an open set when this derivative exists at
each point; no continuity of the derivative is assumed.
\end{definition}
\begin{proposition}[Limits and differentiation rules]\label{d12:01-01-rules}
Complex convergence is equivalent to convergence of both real coordinates,
and, for a punctured limit at \(a\), to convergence along every sequence
\(z_n\ne a\) in the domain tending to \(a\). Differentiability implies
continuity. Sums, products,
quotients where the denominator is nonzero, and compositions have derivatives
\[
 (af+bg)'=af'+bg',\quad (fg)'=f'g+fg',\quad
 (f/g)'=(f'g-fg')/g^2,\quad (g\circ f)'=(g'\circ f)f'.
\]
Every assertion is local at a point where all displayed functions and
derivatives exist; the corresponding holomorphic closure laws hold on the
open sets where they are defined.
\end{proposition}
\begin{proof}
For \(w=x+\iu y\), \(|x|,|y|\leq |w|\leq |x|+|y|\), which proves the
coordinate assertion. A failed limit supplies an \(\eps>0\) and points at
distance less than \(1/n\) violating that tolerance; this proves the
sequential assertion in both directions. Differentiability is precisely
\[
 f(a+h)=f(a)+Ah+h e(h),\qquad e(h)\longrightarrow0.
\]
This identity proves continuity. Adding two such identities proves the sum
rule; multiplying them and collecting the linear term proves the product
rule, since every remaining term divided by \(h\) tends to zero.
If \(g(a)\ne0\), continuity makes \(g\) nonzero near \(a\), and
\[
 \frac{g(a+h)^{-1}-g(a)^{-1}}h
 =-\frac{g(a+h)-g(a)}{h\,g(a+h)g(a)}
 \longrightarrow-\frac{g'(a)}{g(a)^2}.
\]
Combining this with the product rule gives the quotient rule.
For composition put \(b=f(a)\) and write
\(g(b+k)=g(b)+Bk+k\eta(k)\), setting \(\eta(0)=0\).
Then \(k=f(a+h)-f(a)=Ah+he(h)\) tends to zero and \(k/h\to A\).
The composed difference quotient is
\(B(k/h)+(k/h)\eta(k)\), tending to \(BA\); this includes \(k=0\).
\end{proof}
\source{CON III.2, printed 33--35, PDF 47--49; AHL II.1, printed 22--28, PDF 38--44.}

\begin{theorem}[The real differential and CR equations]\family{D12.01.02}\label{d12:01-02}
Identify \(\C\) with \(\R^2\). A function is complex differentiable at \(a\)
if and only if it is real Fr\'echet differentiable there and its real
differential is complex-linear. For \(f=u+\iu v\), assuming real
differentiability at \(a\), this is equivalent to
\[
 u_x=v_y,\qquad u_y=-v_x \quad\hbox{at }a.
\]
In that case \(f'=u_x+\iu v_x=v_y-\iu u_y\), and the real Jacobian is
\(\det D_{\R}f=|f'|^2\).
\end{theorem}
\begin{proof}
The remainder identity in the preceding proof is exactly real Fr\'echet
differentiability with differential \(h\mapsto Ah\).
Conversely a complex-linear map is multiplication by its value \(A\) at
\(1\), and its real remainder \(o(|h|)\) becomes \(o(1)\) after division
by \(h\). A real-linear map commutes with multiplication by \(\iu\)
exactly when its matrix has the form
\(\left(\begin{smallmatrix}p&-q\\q&p\end{smallmatrix}\right)\).
Substitute the real partial derivatives to obtain CR, the derivative
formula and the determinant \(p^2+q^2\).
\end{proof}

\begin{theorem}[Sufficient regularity and its boundary]\family{D12.01.03}\label{d12:01-03}
If \(u,v\) have continuous first real partial derivatives on an open set,
then \(u+\iu v\) is holomorphic exactly when CR holds there.
Existence of the four partial derivatives and CR at one point, even together
with continuity of the function there, does not suffice.
\end{theorem}
\begin{proof}
The real differentiability theorem for \(C^1\) maps and
\cref{d12:01-02} prove sufficiency. Necessity follows from the same matrix
calculation, without needing continuity of the partial derivatives.
For the failure define \(F(0)=0\) and
\(F(z)=z^5/|z|^4\) for \(z\ne0\). Since \(|F(z)|=|z|\), it is continuous
at zero. On both real and imaginary axes it equals \(z\), so at zero
\(u_x=v_y=1\) and \(u_y=v_x=0\). But
\(F(h)/h=h^4/|h|^4\) equals \(1\) for positive real \(h\) and \(-1\)
on the ray \(\arg h=\pi/4\); the derivative does not exist.
\end{proof}
\source{CON III.2, printed 35--44, PDF 49--58; AHL II.1.2, printed 24--28, PDF 40--44.}

\begin{definition}[Three initially distinct notions]\family{D12.01.04}\label{d12:01-04}
A function is \(C^1\)-holomorphic if it is complex differentiable and its
derivative is continuous. It is power-series analytic if near every point
it equals a convergent series \(\sum_{n\ge0}a_n(z-a)^n\).
These definitions initially impose different requirements.
\end{definition}
\begin{proposition}[The eventual equivalence]\label{d12:01-04-equivalence}
On every open set, holomorphic, \(C^1\)-holomorphic and power-series analytic
are equivalent; each implies real smoothness and holomorphic derivatives
of every order.
\end{proposition}
\begin{proof}
The implication from a power series and its derivative regularity is
proved directly in \cref{d12:02-01}. The implication from weak holomorphy
is proved by Goursat and the Cauchy formula in
\cref{d12:04-01,d12:04-02}, with the Taylor identification in
\cref{d12:05-01}. A \(C^1\)-holomorphic function is holomorphic by its
definition. No later conclusion is used in those proofs as an initial
regularity hypothesis.
\end{proof}

\begin{proposition}[Locality and examples]\family{D12.01.05}\label{d12:01-05}
Restrictions of holomorphic functions to open subsets are holomorphic.
Functions holomorphic on members of an open cover and equal on overlaps
glue uniquely to a holomorphic function. Real smoothness does not imply
holomorphy, and differentiability at one point does not imply holomorphy
on any neighborhood of that point.
\end{proposition}
\begin{proof}
A derivative depends only on the values in an arbitrarily small neighborhood;
this proves restriction and the gluing assertion at each point. Equality
on overlaps makes the glued values unambiguous and proves uniqueness.
The real smooth function \(z\mapsto\bar z\) has difference quotient
\(\bar h/h\), equal to \(1\) on real increments and \(-1\) on imaginary
increments, at every point. For \(f(z)=|z|^2\),
\[
 \frac{|a+h|^2-|a|^2}{h}
 =\bar a+a\frac{\bar h}{h}+\bar h.
\]
At \(a=0\) this tends to zero. If \(a\ne0\), real and imaginary increments
give distinct limits, so zero is its only differentiability point.
\end{proof}


% END chapters/01

% BEGIN chapters/02
\chapter{Series, elementary functions and branches}\label{ch:d12-02}

\begin{theorem}[Radius and calculus of a power series]\family{D12.02.01}\label{d12:02-01}
For \(\sum_{n\ge0}a_n(z-a)^n\), put
\[
 L=\limsup_{n\to\infty}|a_n|^{1/n},\qquad R=L^{-1},
 \quad 0^{-1}=\infty,\quad \infty^{-1}=0.
\]
The series converges absolutely for \(|z-a|<R\), diverges for
\(|z-a|>R\), and converges absolutely and uniformly on every closed disk
of radius \(r<R\). Its sum is holomorphic there, its derivatives are the
termwise differentiated series, and each derivative and the termwise
primitive have the same radius \(R\). Thus
\(a_n=f^{(n)}(a)/n!\).
\end{theorem}
\begin{proof}
The root test proves convergence for \(|z-a|L<1\). If \(|z-a|L>1\),
infinitely many terms have modulus bounded below by a number greater than
one, so terms fail to tend to zero. These arguments include \(L=0,\infty\).
For \(r<s<R\), the numbers \(|a_n|s^n\) are bounded by some \(M\), hence
\(|a_n(z-a)^n|\le M(r/s)^n\) on the radius-\(r\) disk.
The geometric tail proves absolute uniform convergence.

For every integer \(k\ge0\), the series with terms bounded by a constant
times \(n^k(r/s)^n\) converges: its successive positive-term ratio tends to
\(r/s<1\). Hence all formal derivative series converge uniformly on smaller
disks. Let \(f_N\) be the polynomial partial sums and \(g\) the uniform
limit of \(f_N'\) on a disk containing the segment from \(z\) to \(z+h\).
Real FTC along that segment gives
\[
 f_N(z+h)-f_N(z)=h\int_0^1f_N'(z+th)\dd t.
\]
Uniform convergence permits passage to the limit. Dividing by \(h\), the
result tends to \(g(z)\), since \(g\) is continuous. Repetition gives every
derivative. As \(n^{1/n}\to1\), multiplying coefficients by the nonzero
polynomial factors \(n(n-1)\cdots(n-k+1)\) does not change the limsup root;
finite initial terms have no effect. Division by \(n+1\) and shifting the
index gives the primitive radius. Polynomial and zero-series cases have
radius infinity and satisfy the same conclusions. Evaluating the
\(n\)-th derivative at the center gives the coefficient formula.
When \(R=0\) there is no open convergence disk and no holomorphy conclusion.
\end{proof}

\begin{theorem}[Series algebra and composition]\family{D12.02.02}\label{d12:02-02}
On a common open convergence disk, sums are coefficientwise and products
are Cauchy products. Power-series coefficients representing a function
near their center are unique. If \(F(a)=b\), \(F\) is represented by a
power series near \(a\), and \(G\) by a power series near \(b\), then
\(G\circ F\) has the convergent formally substituted series on a sufficiently
small disk about \(a\).
\end{theorem}
\begin{proof}
On any smaller disk the sums of absolute values converge. A double series
with finite sum of absolute values can be rearranged: after a finite
rectangle its total omitted absolute mass is arbitrarily small, so every
finite exhaustion has the same limit. Applying this to products proves
the Cauchy product formula, including uniform convergence on that disk.
The derivative formula in \cref{d12:02-01} proves uniqueness.

Write \(F(a+w)-b=\sum_{n\ge1}c_nw^n\) and
\(G(b+\zeta)=\sum_{m\ge0}d_m\zeta^m\).
Choose \(r>0\) so small that
\(H(r)=\sum_{n\ge1}|c_n|r^n\) is strictly inside the convergence radius
of \(G\); this is possible since \(H(r)\to0\).
Every substituted monomial is absolutely dominated on \(|w|\le r\) by
\(|d_m|H(r)^m\), whose sum is finite. The preceding absolute-tail argument
therefore allows expansion of each power and regrouping by total degree.
The resulting power series is absolutely uniformly convergent and equals
the composition. If \(F\) is constant the assertion reduces to a constant
series.
\end{proof}
\source{CON III.1--2, printed 30--37, PDF 44--51; AHL II.2, printed 33--42, PDF 49--58.}

\begin{theorem}[Exponential and trigonometric functions]\family{D12.02.03}\label{d12:02-03}
The entire function \(E(z)=\sum_{n\ge0}z^n/n!\) satisfies
\(E'=E\), \(E(z+w)=E(z)E(w)\), \(E(0)=1\), and never vanishes.
It agrees with the real exponential on \(\R\), and
\[
 E(x+\iu y)=e^x(\cos y+\iu\sin y),\qquad
 E(z)=E(w)\ \Longleftrightarrow\ z-w\in2\pi\iu\Z .
\]
Define
\[
 \cos z=\tfrac12(E(\iu z)+E(-\iu z)),\quad
 \sin z=\tfrac1{2\iu}(E(\iu z)-E(-\iu z)),
\]
and \(\cosh z=(E(z)+E(-z))/2,\ \sinh z=(E(z)-E(-z))/2\).
Their derivatives are \(\sin'=\cos,\ \cos'=-\sin,\
\sinh'=\cosh,\ \cosh'=\sinh\). Sine and hyperbolic sine are odd;
cosine and hyperbolic cosine are even. Their addition identities and
\(\sin^2z+\cos^2z=1,\ \cosh^2z-\sinh^2z=1\) hold for complex
arguments. The period set of sine and cosine is \(2\pi\Z\);
their zeros are respectively \(\pi\Z\) and \(\pi/2+\pi\Z\), all simple.
The hyperbolic functions are obtained by imaginary substitution and
have period set \(2\pi\iu\Z\). On their nonzero-denominator domains,
define \(\tan=\sin/\cos,\ \cot=\cos/\sin,\
\tanh=\sinh/\cosh,\ \coth=\cosh/\sinh\); their derivatives are
\(1/\cos^2,-1/\sin^2,1/\cosh^2,-1/\sinh^2\), respectively.
\end{theorem}
\begin{proof}
The coefficient ratio gives infinite radius and differentiation gives
\(E'=E\). The absolutely convergent double sum and the binomial identity
give \(E(z)E(w)=\sum_k(z+w)^k/k!\).
Consequently \(E(z)E(-z)=1\).
On real arguments the series is the real exponential series; substituting
\(\iu y\) and separating even and odd terms gives the established real
sine and cosine series. Thus the displayed polar formula follows, and
the real exponential's injectivity and the real trigonometric period
criterion prove the kernel assertion. Substitution into the definitions,
multiplication using the addition formula for \(E\), and collecting the
four terms give
\(\sin(z+w)=\sin z\cos w+\cos z\sin w\),
\(\cos(z+w)=\cos z\cos w-\sin z\sin w\), and
\(\sin^2z+\cos^2z=1\).
Differentiating gives \(\sin'=\cos,\ \cos'=-\sin\); replacement by
\(\iu z\) gives \(\sinh'=\cosh,\ \cosh'=\sinh\) and
\(\cosh^2-\sinh^2=1\).
The same substitution in the two addition formulas gives
\(\sinh(z+w)=\sinh z\cosh w+\cosh z\sinh w\) and
\(\cosh(z+w)=\cosh z\cosh w+\sinh z\sinh w\).
The equations \(\sin z=0\) and \(\cos z=0\) become
\(E(2\iu z)=1\) and \(E(2\iu z)=-1\). The kernel criterion gives the
listed zeros. At those points the derivatives are \(\pm1\), so the
zeros are simple. The addition formulas give parity and period \(2\pi\).
If \(p\) is any sine period, substituting \(z=0,\pi/2\) gives
\(\sin p=0,\cos p=1\), so the zero list forces \(p\in2\pi\Z\).
For a cosine period, \(\cos p=1\) gives
\((E(\iu p)-1)^2=0\), hence the same conclusion.
Imaginary substitution gives the hyperbolic period sets.
The quotient rule and the two Pythagorean identities prove all four
quotient derivative formulas on exactly their stated domains.
\end{proof}

\begin{theorem}[Principal logarithm]\family{D12.02.04}\label{d12:02-04}
On \(S=\C\setminus(-\infty,0]\) put
\(\operatorname{Log}z=\log|z|+\iu\operatorname{Arg}z\), where
\(-\pi<\operatorname{Arg}z<\pi\).
It is holomorphic, \((\operatorname{Log})'=1/z\),
\(E(\operatorname{Log}z)=z\), and \(\operatorname{Log}1=0\).
It is the inverse of \(E\) from the strip \(|\Imm z|<\pi\) onto \(S\).
Any two continuous logarithms of the same nonzero function on a connected
set differ by a constant in \(2\pi\iu\Z\).
There is no continuous logarithm of \(z\) on \(\C\setminus\{0\}\).
\end{theorem}
\begin{proof}
For \(z=x+\iu y\in S\), \(r=(x^2+y^2)^{1/2}\) satisfies \(r+x>0\) and
\[
 \operatorname{Arg}z=2\arctan\frac y{r+x}.
\]
Real differentiation gives
\((\log r)_x=x/r^2,\ (\log r)_y=y/r^2\) and
\(\operatorname{Arg}_x=-y/r^2,\ \operatorname{Arg}_y=x/r^2\).
These continuous derivatives satisfy CR; \cref{d12:01-02} gives \(1/z\).
Polar representation proves the exponential identity and strip inverse.
Two logarithms differ pointwise by an element of \(2\pi\iu\Z\), by
\cref{d12:02-03}. A continuous map from a connected set to this discrete
set is constant, since inverse images of individual values are open and
closed. If a logarithm on \(\C^*\) existed, along \(e^{\iu t}\),
\(0\le t\le2\pi\), its difference from \(\iu t\) would be a continuous
discrete-valued function, hence constant. Equal endpoints would then
differ by \(2\pi\iu\), a contradiction.
\end{proof}

\begin{corollary}[Powers with a chosen logarithm]\label{d12:02-04-powers}
On the principal slit plane, for every \(\alpha\in\C\),
\(z^\alpha=\exp(\alpha\operatorname{Log}z)\) is holomorphic with
derivative \(\alpha z^{\alpha-1}\), and
\(z^\alpha z^\beta=z^{\alpha+\beta}\) for this one chosen logarithm.
Changing a logarithm by \(2\pi\iu k\) multiplies the power by
\(e^{2\pi\iu k\alpha}\).
\end{corollary}
\begin{proof}
The chain rule gives \(\alpha e^{\alpha\operatorname{Log}z}/z\).
Since \(e^{\operatorname{Log}z}=z\), this is the stated derivative.
The exponential addition formula gives both multiplication statements.
\end{proof}

\begin{theorem}[Local branches]\family{D12.02.05}\label{d12:02-05}
If \(f\) is holomorphic and \(f(a)\ne0\), every chosen logarithm \(c\) of
\(f(a)\) extends to a holomorphic logarithm \(L\) of \(f\) near \(a\), with
\(L'=f'/f\). For each positive integer \(n\), every chosen \(n\)-th root
of \(f(a)\) extends to a holomorphic \(n\)-th root near \(a\).
Exponential and nonzero power maps have the corresponding local inverse
branches.
\end{theorem}
\begin{proof}
Continuity permits a disk \(U\) about \(a\) on which
\(|f(z)/f(a)-1|<1/2\). Then
\(L(z)=c+\operatorname{Log}(f(z)/f(a))\) has all the stated properties,
by the chain rule and \cref{d12:02-04}.
The function \(E(L/n)\) is a root, and multiplication by a constant
\(n\)-th root of unity adjusts it to any prescribed root at \(a\).
The logarithm is a local inverse of the exponential after adding the
chosen \(2\pi\iu\) multiple; roots are local inverses of powers at
nonzero points.
\end{proof}
\begin{corollary}[The general local inverse branch]\label{d12:02-05-inverse}
A holomorphic function with nonzero derivative at a point has a
holomorphic inverse branch on a neighborhood of its value.
\end{corollary}
\begin{proof}
This is the real inverse theorem and complex differential argument
proved explicitly in \cref{d12:05:inverse-noncritical}. It is recorded
separately so that the local logarithm/root construction just proved
does not depend on the later inverse theorem, which uses those roots
only in the subsequent local mapping normal form.
\end{proof}
\source{CON III.2, printed 37--44, PDF 51--58; AHL II.3, printed 42--48, PDF 58--64.}

\begin{theorem}[Abel's radial limit]\family{D12.02.06}\label{d12:02-06}
If \(\sum_{n\ge0}a_n\zeta^n=S\) converges for a nonzero \(\zeta\), then
\[
 \lim_{t\uparrow1}\sum_{n\ge0}a_n(t\zeta)^n=S,\qquad 0\le t<1.
\]
The interior sums converge absolutely. No general boundary convergence
or boundary extension follows merely from interior holomorphy.
\end{theorem}
\begin{proof}
Put \(b_n=a_n\zeta^n\) and \(S_N=\sum_{n=0}^Nb_n\).
These partial sums are bounded, say by \(M\); hence \(|b_n|\le2M\)
after enlarging \(M\) to include the first term. This proves absolute
convergence for \(t<1\). Finite summation by parts, with the last term
tending to zero, gives
\[
 \sum_{n\ge0}b_nt^n=(1-t)\sum_{n\ge0}S_nt^n.
\]
Given \(\eps>0\), choose \(N\) with \(|S_n-S|<\eps\) for \(n\ge N\).
The tail in \((1-t)\sum(S_n-S)t^n\) is at most \(\eps\), while its
finite head tends to zero. This proves the limit.
The series \(\sum z^n\) is holomorphic for \(|z|<1\), but at every
\(|z|=1\) its terms fail to tend to zero; at \(1\) its sum
\(1/(1-z)\) has no finite radial limit. Thus the boundary hypothesis is
substantive. The theorem asserts a radial limit, not an unrestricted
approach theorem.
\end{proof}
\source{AHL II.2.5, printed 41--42, PDF 57--58; CON III.1, printed 30--33, PDF 44--47.}


% END chapters/02

% BEGIN chapters/03
\chapter{Contours, Stieltjes integrals and primitives}\label{ch:d12-03}

\begin{definition}[Paths and finite chains]\family{D12.03.01}\label{d12:03-01}
A path is a continuous \(\gamma:[a,b]\to\C\), \(a\le b\); its trace is
\(|\gamma|=\gamma([a,b])\). It is piecewise \(C^1\) when a finite
partition makes each restricted path \(C^1\), including its one-sided
endpoint derivatives. For a continuous \(f\) on its trace set
\[
 \int_\gamma f(z)\dd z=\int_a^b f(\gamma(t))\gamma'(t)\dd t,
\]
summing over those pieces. A finite chain is a finite integral linear
combination of paths; its integral is the corresponding linear sum.
A cycle uses closed paths. A degenerate or constant path has integral zero.
\end{definition}
\begin{proposition}[Real coordinate interpretation]\label{d12:03-01-coordinates}
For \(f=u+\iu v,\ \gamma=x+\iu y\),
\[
 \int_\gamma f\dd z
 =\int_\gamma(u\dd x-v\dd y)+\iu\int_\gamma(v\dd x+u\dd y).
\]
Refining the partition in the definition changes no value.
\end{proposition}
\begin{proof}
On each \(C^1\) piece multiply
\((u+\iu v)(x'+\iu y')\) and integrate its coordinates.
Real integral additivity proves refinement invariance, including all
one-sided endpoints, which affect no integral.
\end{proof}

\begin{theorem}[Reversal, concatenation and parameter changes]\family{D12.03.02}\label{d12:03-02}
Integrals add under subdivision and concatenation and change sign under
reversal. For an increasing piecewise \(C^1\) surjection \(\phi\) of
compact parameter intervals, when both paths are piecewise \(C^1\),
\(\int_{\gamma\circ\phi}f\dd z=\int_\gamma f\dd z\).
The same equality holds for every continuous nondecreasing surjection
when the rectifiable integral below is used, with no differentiability
or strict monotonicity of \(\phi\).
\end{theorem}
\begin{proof}
Real integral additivity proves subdivision; affine rescaling of the two
parameter intervals proves concatenation. The substitution
\(t=a+b-s\) gives reversal and its minus sign.
For differentiable monotone \(\phi\), the chain rule and the real
change-of-variable formula on each piece give
\(\int f(\gamma(\phi(s)))\gamma'(\phi(s))\phi'(s)\dd s\).
Its value is the original integral; intervals where \(\phi\) is constant
contribute zero. The asserted general version, which also proves equality
for a nonsmooth reparameterization, is established in
\cref{d12:03-06-reparameterization} directly from sums.
\end{proof}

\begin{theorem}[Length, estimates and integral limits]\family{D12.03.03}\label{d12:03-03}
For a piecewise \(C^1\) path its polygonal length is
\[
 \ell(\gamma)=\sup_{a=t_0<\cdots<t_m=b}
       \sum_j|\gamma(t_j)-\gamma(t_{j-1})|
       =\int_a^b|\gamma'(t)|\dd t.
\]
For continuous \(f\) on its trace,
\[
 \left|\int_\gamma f\dd z\right|
 \le\int_a^b|f(\gamma(t))|\,|\gamma'(t)|\dd t
 \le\ell(\gamma)\sup_{|\gamma|}|f|.
\]
Uniform convergence of continuous \(f_n\) to \(f\) on that fixed trace
implies convergence of their integrals. These estimates and uniform
interchange extend to all rectifiable paths.
\end{theorem}
\begin{proof}
For a \(C^1\) piece, FTC bounds every chord sum by \(\int|\gamma'|\).
Uniform continuity of \(\gamma'\) gives, for a partition of mesh \(h\)
and tags \(\xi_j\),
\[
 \left|\gamma(t_j)-\gamma(t_{j-1})
       -\gamma'(\xi_j)(t_j-t_{j-1})\right|
 \le\omega_{\gamma'}(h)(t_j-t_{j-1}).
\]
Consequently chord sums tend to the Riemann integral of \(|\gamma'|\);
taking suprema gives the reverse length inequality. Insert all the
piece boundaries into any partition and add the equalities to handle
a piecewise path. The integral inequalities are the coordinate complex
integral's norm inequality. Apply the last bound to \(f_n-f\) to obtain
the convergence assertion. The rectifiable extension follows from the
same sum estimate in \cref{d12:03-06}.
\end{proof}

\begin{lemma}[Polygonal connectivity]\label{d12:03:polygonal-connectivity}
Every open connected subset of \(\C\) is polygonally connected.
\end{lemma}
\begin{proof}
Fix \(a\) in the set and let \(E\) be points reachable from \(a\) by a finite
polygonal path in the set. If \(z\) is any point of the set, a small disk
about \(z\) lies in it. If that disk meets \(E\), all its points are reachable
by adding a segment. Thus \(E\) and its relative complement are open.
Connectedness and \(a\in E\) give the assertion.
\end{proof}

\begin{theorem}[Primitives and zero periods]\family{D12.03.04}\label{d12:03-04}
For a continuous \(f\) on a domain \(\Omega\), the following are equivalent:
\(f\) has a holomorphic primitive \(F\) with \(F'=f\);
its integrals over piecewise \(C^1\) paths depend only on endpoints;
its integral over every closed piecewise \(C^1\) path is zero.
For any such primitive and any rectifiable path in \(\Omega\),
\(\int_\gamma f\dd z=F(\gamma(b))-F(\gamma(a))\).
Primitives on a domain are unique up to an additive constant.
\end{theorem}
\begin{proof}
For a piecewise \(C^1\) path the real chain rule and FTC give the endpoint
formula, so a primitive implies path independence and zero periods.
Conversely compare two paths with the same endpoints by concatenating
one with the reverse of the other; zero periods imply independence.
Fix a basepoint and define \(F(z)\) by integration along a polygonal path
to \(z\), available by \cref{d12:03:polygonal-connectivity}.
For a sufficiently small \(h\), independence gives
\[
 \frac{F(z+h)-F(z)}h=\int_0^1f(z+th)\dd t\longrightarrow f(z).
\]
Thus \(F'=f\). A function with zero complex derivative has zero real
differential by \cref{d12:01-02}; real FTC along polygonal paths makes it
constant, proving uniqueness.

For the rectifiable endpoint formula use the Stieltjes definition below.
The trace is compact in \(\Omega\); sufficiently close points of the
trace are joined by a segment in \(\Omega\).
Since \(f\) is uniformly continuous on a fixed compact neighborhood of
the trace, real FTC on those segments gives
\[
 |F(w)-F(z)-f(z)(w-z)|\le\omega(|w-z|)|w-z|,
 \qquad \omega(r)\to0.
\]
Apply this to successive path values of a fine partition. The sum of
the errors is bounded by \(\omega(\max|\Delta\gamma|)\ell(\gamma)\).
Continuity of \(\gamma\) makes this tend to zero. The \(F\)-increments
telescope, while the other sum tends to the Stieltjes integral.
\end{proof}

\begin{theorem}[Triangle criterion for a local primitive]\family{D12.03.05}\label{d12:03-05}
Suppose \(\Omega\) is star-shaped about \(a\), \(f\) is continuous there,
and its integral around every triangle whose closed interior lies in
\(\Omega\) is zero. Then
\(F(z)=\int_{[a,z]}f\dd w\) is a primitive of \(f\).
In particular this holds on a convex open set.
\end{theorem}
\begin{proof}
Fix \(z\). The segment \([a,z]\) is compact in \(\Omega\), so some
\(\delta\)-neighborhood of it lies in \(\Omega\). If \(|h|<\delta\),
the triangle with vertices \(a,z,z+h\) lies in that neighborhood.
Its zero boundary integral gives
\(F(z+h)-F(z)=\int_{[z,z+h]}f\dd w\).
Divide by \(h\) and use the last continuity estimate in the preceding proof.
Degenerate triangles contribute zero by reversal and subdivision.
\end{proof}
\source{CON IV.1, printed 58--68, PDF 72--82; AHL IV.1.1--1.3, printed 101--109, PDF 117--125.}

\begin{theorem}[The rectifiable/Stieltjes representation]\family{D12.03.06}\label{d12:03-06}
A continuous path is rectifiable exactly when its total variation
\(V(\gamma)=\ell(\gamma)\) is finite. For such a rectifiable path and any
continuous \(u:[a,b]\to\C\), the tagged sums
\[
 \sum_j u(\xi_j)\bigl(\gamma(t_j)-\gamma(t_{j-1})\bigr),
 \qquad t_{j-1}\le\xi_j\le t_j,
\]
have a unique limit as the partition mesh tends to zero.
Call it \(\int_a^b u\dd\gamma\), and define
\(\int_\gamma f\dd z=\int_a^b(f\circ\gamma)\dd\gamma\).
It is linear, additive on subintervals, and bounded in modulus by
\(V(\gamma)\sup|u|\). If \(\gamma\) is AC, it agrees with
\(\int_a^b u(t)\gamma'(t)\dd t\), hence with the original
piecewise \(C^1\) integral.
\end{theorem}
\begin{proof}
The first assertion compares identical polygonal suprema.
If \(Q\) refines a partition \(P\), express each \(P\)-increment as the
sum of its \(Q\)-increments. The difference of the tagged sums is at most
\(\omega_u(\operatorname{mesh}P)V(\gamma)\).
For two fine partitions pass to their common refinement; their sums
differ by at most twice that quantity. The complex Cauchy criterion
gives existence and uniqueness independently of tags.
Linearity and the bound hold for every sum and hence their limits.
For additivity insert the desired subdivision points; refinement changes
sums by a quantity tending to zero, and each resulting subinterval sum
has its own limit.

If \(\gamma\) is AC, the already established real/complex AC FTC
(D11, compact AC fundamental theorem) gives
\(\Delta\gamma_j=\int_{t_{j-1}}^{t_j}\gamma'(t)\dd t\).
The sum is therefore \(\int u_P\gamma'\), with \(u_P\) the tagged step
function. Since \(\sup|u_P-u|\le\omega_u(\operatorname{mesh}P)\) except
irrelevant endpoint choices, its difference from \(\int u\gamma'\)
is at most \(\omega_u(\operatorname{mesh}P)\int|\gamma'|\), tending
to zero. A piecewise \(C^1\) path is AC by adding the real FTC estimates
on its finitely many pieces; thus both representations coincide.
\end{proof}

\begin{proposition}[All monotone parameter changes]\label{d12:03-06-reparameterization}
Let \(\phi:[c,d]\to[a,b]\) be a continuous nondecreasing surjection.
Then \(V(\gamma\circ\phi)=V(\gamma)\) and
\(\int_{\gamma\circ\phi}f\dd z=\int_\gamma f\dd z\).
For a nonincreasing surjection the integral has the opposite sign.
\end{proposition}
\begin{proof}
The images of partition points form a nondecreasing partition; deleting
repeated images only deletes zero increments. This bounds one variation
by the other. Conversely, choose preimages of each member of any strictly
increasing partition of \([a,b]\). Their order is strictly increasing;
use \(c,d\) for the endpoints. This proves the reverse bound.
Uniform continuity of \(\phi\) makes the image mesh of fine source
partitions tend to zero. The corresponding sums for \(f\circ\gamma\)
agree exactly after deleting zero increments; mapped tags remain in their
respective image intervals. Passage to the limit proves invariance.
A decreasing map is an increasing one followed by the affine reversal,
whose sums reverse all increments and hence the sign.
\end{proof}

\begin{proposition}[Measure and arclength bridges]\label{d12:03-06-measure-arclength}
For continuous BV \(\gamma\), its normalized constant extension
\(\widetilde\gamma\), zero to the left after subtraction of \(\gamma(a)\),
has a finite complex Stieltjes measure \(\nu_\gamma\) satisfying
\[
 \nu_\gamma((s,t])=\gamma(t)-\gamma(s),\quad
 |\nu_\gamma|([a,b])=V(\gamma),\quad \nu_\gamma(\{a\})=0.
\]
For continuous \(u\) on \([a,b]\),
\(\int u\dd\gamma=\int_{[a,b]}u\dd\nu_\gamma\).
If \(L=V(\gamma)>0\), there is a \(1\)-Lipschitz path
\(\sigma:[0,L]\to\C\) and a continuous nondecreasing surjection
\(s:[a,b]\to[0,L]\) such that
\[
 \gamma=\sigma\circ s,\qquad
 \int_\gamma f\dd z=\int_0^L f(\sigma(v))\sigma'(v)\dd v,\qquad
 |\sigma'|=1\ \hbox{a.e.}
\]
For \(L=0\) the path is constant and every integral is zero.
\end{proposition}
\begin{proof}
The normalized constant extension is continuous globally BV, and is zero
at minus infinity. The exact normalized scalar Stieltjes correspondence
and variation identity of D11 give the displayed measure assertions;
no incomplete-measure representative is substituted for the original path.
A tagged step function on the half-open intervals \((t_{j-1},t_j]\)
has measure integral equal to the Stieltjes sum. Give it any value at
\(a\), whose mass is zero. Uniform convergence of these approximants
to \(u\) and the finite variation estimate prove the equality of integrals.

For completeness the arclength factorization can be proved directly.
Let \(s(t)=V(\gamma;[a,t])\). Inserting a point into partitions gives
variation additivity. The function \(s\) is continuous: for \(x<z\),
choose a partition of \([x,z]\) within \(\eps\) of its variation and let
\(t>x\) be its first point. Inserting partitions of \([x,y]\), \(x<y<t\),
and then \(t\), gives
\[
 V(\gamma;[x,y])\le\eps+
 |\gamma(t)-\gamma(x)|-|\gamma(t)-\gamma(y)|
 \le\eps+|\gamma(y)-\gamma(x)|.
\]
Let \(y\downarrow x\) and then \(\eps\downarrow0\).
Reversing the interval proves the left version, including endpoints.
Thus \(s\) is continuous and onto \([0,L]\).
If \(s(t)=s(u)\), variation between \(u,t\) is zero, so
\(\gamma(t)=\gamma(u)\). Set \(\sigma(s(t))=\gamma(t)\).
For ordered values,
\(|\sigma(s(t))-\sigma(s(u))|\le s(t)-s(u)\); thus \(\sigma\) is
well-defined and \(1\)-Lipschitz. The parameter-change proposition gives
its length \(L\) and the equality of contour integrals.
Lipschitz functions are AC. The D11 exact variation formula for AC
primitives gives \(L=\int_0^L|\sigma'|\); the real Lipschitz derivative
bound is \(|\sigma'|\le1\) a.e. Therefore
\(\int(1-|\sigma'|)=0\), proving the a.e. unit-speed assertion and
the AC integral formula. If \(L=0\), every two path values have distance
at most zero.
\end{proof}
\source{CON IV.1, printed 58--67, PDF 72--81; AHL IV.1.2--1.3, printed 104--109, PDF 120--125; D11 compact AC FTC and normalized scalar Stieltjes/variation correspondence.}

\begin{lemma}[A finite-length trace has empty interior]\label{d12:03:trace-empty-interior}
The trace of a rectifiable path has planar Lebesgue measure zero and hence
empty interior.
\end{lemma}
\begin{proof}
Use its arclength path. Divide \([0,L]\) into at most \(1+L/\delta\)
intervals of length at most \(\delta\). Their images lie in disks of
radius \(\delta\). The trace's outer area is at most
\(\pi\delta^2(1+L/\delta)\), tending to zero.
A nonempty open disk has positive area, so it cannot lie in that trace.
Constant paths have singleton traces and obey the same conclusion.
\end{proof}


% END chapters/03

% BEGIN chapters/04
\chapter{Goursat, winding numbers and Cauchy theory}\label{ch:d12-04}

\begin{theorem}[Weak Goursat theorem]\family{D12.04.01}\label{d12:04-01}
If \(f\) is complex differentiable on a neighborhood of a closed triangle
\(T\), its positively oriented boundary integral is zero. The same holds
for a rectangle by division along a diagonal. Continuity of \(f'\) is
not a hypothesis.
\end{theorem}
\begin{proof}
Divide a nondegenerate triangle into four half-sized triangles by its
side midpoints. Internal integrals cancel, so one child has boundary
integral of modulus at least a quarter of the parent's. Iterate,
obtaining nested closed triangles \(T_n\), diameter \(2^{-n}d\) and
perimeter \(2^{-n}p\), with
\[
 \left|\int_{\partial T}f\dd z\right|
 \le4^n\left|\int_{\partial T_n}f\dd z\right|.
\]
Compactness and the shrinking diameter give a unique common point \(a\).
Differentiability there gives
\(f(z)=f(a)+f'(a)(z-a)+r(z)\), with
\(|r(z)|\le\eps|z-a|\) in a sufficiently small disk.
Constants and linear functions have primitives and zero boundary integrals.
For sufficiently large \(n\), therefore,
\[
 \left|\int_{\partial T_n}f\dd z\right|
 \le\eps\,\diam(T_n)\operatorname{length}(\partial T_n)
 =\eps\,4^{-n}dp .
\]
Thus the original integral is at most \(\eps dp\) for every \(\eps>0\).
Degenerate triangles cancel by oriented segment additivity. Adding the
two triangle formulas proves the rectangle assertion.
\end{proof}

\begin{lemma}[A continuous exceptional point]\label{d12:04:exceptional-point}
The same zero-triangle conclusion holds if \(f\) is continuous near the
triangle and holomorphic there except possibly at one point.
\end{lemma}
\begin{proof}
Subdivide the triangle by the \(N\)-fold affine triangular grid into
\(N^2\) small triangles of perimeter \(p/N\).
At most six of their closures contain the exceptional point. Every other
closed small triangle has a holomorphic neighborhood, so its integral
is zero by Goursat. On a compact neighborhood let \(|f|\le M\).
The remaining integrals have total modulus at most \(6Mp/N\).
All interior sides cancel. Let \(N\to\infty\).
If the exceptional point lies outside the triangle there are no remaining
terms; if it lies on a side or vertex the same bound applies.
\end{proof}

\begin{theorem}[Disk Cauchy theorem, formula and all derivatives]\family{D12.04.02}\label{d12:04-02}
A holomorphic function on a disk has a primitive and zero integral on every
closed rectifiable path there. If \(f\) is holomorphic on a neighborhood
of \(\overline{B(a,r)}\), and \(C\) is its positively oriented circle, then
for \(|z-a|<r\),
\[
 f(z)=\frac1{2\pi\iu}\int_C\frac{f(w)}{w-z}\dd w,\qquad
 f^{(n)}(z)=\frac{n!}{2\pi\iu}
               \int_C\frac{f(w)}{(w-z)^{n+1}}\dd w .
\]
All derivatives exist and are holomorphic. At the center, if
\(M=\max_C|f|\), then \(|f^{(n)}(a)|\le n!M/r^n\).
\end{theorem}
\begin{proof}
Goursat and \cref{d12:03-05} construct the primitive on a disk;
\cref{d12:03-04} gives zero closed integrals, including rectifiable ones.
For a fixed \(z\) in the circle define
\[
 g_z(w)=\frac{f(w)-f(z)}{w-z}\quad(w\ne z),\qquad g_z(z)=f'(z).
\]
It is continuous and holomorphic off \(z\).
The exceptional-point lemma and the triangle primitive criterion show
that \(g_z\) has a primitive on a slightly larger disk containing \(C\).
Its integral on \(C\) is zero. On \(C\), the expansion
\[
 \frac1{w-z}=\sum_{k\ge0}\frac{(z-a)^k}{(w-a)^{k+1}}
\]
is uniformly convergent. The first term has integral \(2\pi\iu\) by
\(w=a+re^{\iu t}\); every other term has a primitive on the punctured
plane and zero circle integral. Thus
\(\int_C(w-z)^{-1}\dd w=2\pi\iu\), proving the first formula.

For \(J_m(z)=\int_C f(w)(w-z)^{-m}\dd w\), on a neighborhood of any
interior \(z_0\) the denominators are uniformly separated from zero.
The rational kernel has difference quotient
\[
 \frac{(w-z-h)^{-m}-(w-z)^{-m}}h
   =m(w-z)^{-m-1}+O(|h|),
\]
uniformly on \(C\); the bound follows by factoring the difference of
the \(m\)-th powers of the denominators, or by its finite geometric
expansion, with both denominators bounded below.
The ML estimate therefore gives \(J_m'=mJ_{m+1}\).
The same uniform denominator estimate proves continuity of every \(J_m\).
Differentiate the first formula repeatedly to get the displayed formulas;
they show each derivative is again differentiable. The circle length
\(2\pi r\) gives the center estimate.
\end{proof}

\begin{lemma}[A preliminary Liouville theorem]\label{d12:04:liouville}
A bounded entire function is constant.
\end{lemma}
\begin{proof}
If \(|f|\le M\), the derivative estimate at any \(a\), on every radius
\(r>0\), gives \(|f'(a)|\le M/r\). Thus \(f'(a)=0\). The zero-derivative
conclusion of \cref{d12:03-04} makes \(f\) constant.
This lemma uses only the local disk formula, not the global formula below.
\end{proof}
\source{CON IV.2, printed 68--76, PDF 82--90; IV.8, printed 100--102, PDF 114--116; AHL IV.1.4--2.3, printed 109--124, PDF 125--140.}

\begin{definition}[Index of a cycle]\family{D12.04.03}\label{d12:04-03}
For a cycle \(\Gamma=\sum_jm_j\gamma_j\), with
\(|\Gamma|=\bigcup_{m_j\ne0}|\gamma_j|\), and \(a\notin|\Gamma|\), set
\[
 \ind(\Gamma,a)=\frac1{2\pi\iu}\int_\Gamma\frac{\dd z}{z-a}.
\]
A zero chain has empty trace and zero index.
\end{definition}
\begin{theorem}[Integer index and its support]\label{d12:04-03-properties}
The index is integer-valued, continuous and locally constant off the
trace, and zero on its unbounded complementary component.
\end{theorem}
\begin{proof}
First consider a closed path \(\gamma\) avoiding \(a\).
Uniform continuity provides a partition such that each subpath
\(\gamma(t)-a\) lies in a small disk avoiding zero on which a logarithm
is available from \cref{d12:02-05}. Choose its logarithm on the first
piece, and add a \(2\pi\iu\) multiple on each following piece to make
endpoint values agree. The resulting continuous \(L(t)\), piecewise
\(C^1\), satisfies
\(E(L(t))=\gamma(t)-a\) and \(L'=\gamma'/(\gamma-a)\).
Thus the integral is \(L(b)-L(a)\); equal path endpoints force that
difference to be \(2\pi\iu k\), \(k\in\Z\).
Integer linear combinations preserve this property.

On a small disk in the complement of the trace the reciprocal kernels
vary uniformly continuously, since their denominators have a fixed
positive lower bound. The ML estimate proves continuity in \(a\).
A continuous integer-valued function is locally constant.
For \(|a|\to\infty\), the estimate
\[
 |\ind(\Gamma,a)|\le
 \frac{\sum_j|m_j|\ell(\gamma_j)}
 {2\pi\,\dist(a,|\Gamma|)}
\]
makes the index tend to zero, so its integer value is eventually zero.
Local constancy then makes it zero throughout the connected unbounded
component.
\end{proof}

\begin{proposition}[Index operations and homotopy]\family{D12.04.04}\label{d12:04-04}
Index adds under chain addition and concatenation and changes sign under
reversal. If two closed piecewise \(C^1\) paths are joined by a continuous
homotopy of closed paths avoiding \(a\), their indices about \(a\) agree.
The homotopy need not be differentiable and its intermediate paths need
not be rectifiable.
\end{proposition}
\begin{proof}
The algebraic assertions follow from integral additivity.
We prove a more general homotopy integral assertion without differentiating
the homotopy. Let \(f\) be holomorphic on an open set \(\Omega\), and let
\(H:[0,1]^2\to\Omega\) be a continuous homotopy with piecewise \(C^1\)
bottom and top paths. Suppose either its endpoints are fixed or its
left and right sides agree, as in a homotopy of closed loops.
Its compact image has positive distance from the complement of \(\Omega\)
(if the complement is empty choose any positive scale).
Uniform continuity permits a rectangular grid so fine that the image
of every cell lies in a disk contained in \(\Omega\).
Join images of adjacent grid vertices by straight segments.
Every cell polygon lies in such a disk, so its integral is zero by
the disk primitive theorem. Summing cells cancels every interior edge.
For each bottom or top subpath, both it and its endpoint chord lie in
one of those disks, and their integrals agree by that disk's primitive.
The side chord integrals cancel, because their vertex images agree on
the two sides or are constant for fixed endpoints.
Thus bottom and top integrals agree.
Apply this assertion to \(f(z)=1/(z-a)\) on \(\C\setminus\{a\}\).
\end{proof}

\begin{lemma}[Practical index computations]\label{d12:04:crossing-index}
A positively traversed circle has index one at each interior point and
zero at each exterior point. More generally, suppose a closed piecewise
\(C^1\) path avoids \(a\), starts off a ray \(a+[0,\infty)\eta\)
(\(|\eta|=1\)), and has finitely many transverse crossings of that ray,
none at a corner. Then its index is
\[
 \sum_{\gamma(t)-a\in(0,\infty)\eta}
        \operatorname{sign}\Imm\!\left(\gamma'(t)/\eta\right).
\]
In particular positive rectangle and half-disk boundaries have index
one in their indicated interiors and zero in their exteriors.
These statements concern these explicit regions, not separation by an
arbitrary Jordan curve.
\end{lemma}
\begin{proof}
The circle's interior value was computed by the geometric reciprocal
expansion in the disk formula proof. For an exterior point expand instead
in powers of \((w-a_0)/(a-a_0)\), where \(a_0\) is the circle's center;
all resulting circle monomials integrate to zero.
For the crossing formula rotate and translate the ray to the positive
axis. Between crossings choose the argument in \((0,2\pi)\).
At an upward crossing it jumps from \(2\pi\) to zero; its continuous
lift must be increased by \(2\pi\). At a downward crossing it must be
decreased by \(2\pi\). Away from these finitely many neighborhoods the
same argument branch is continuous. At the endpoint the branch agrees
with its starting value; hence the total lifted argument increment is
exactly the signed crossing sum times \(2\pi\).
The index proof above identifies this increment with the index.
For a point inside a rectangle or a half-disk, a horizontal ray to its
right exits across precisely one positively crossed boundary piece;
choose a slightly tilted ray if it would meet a corner. It gives one.
For an exterior point, the convex closed region contracts to a point
within itself, avoiding that exterior point; homotopy invariance gives
zero.
\end{proof}

\begin{corollary}[Index from a specified region and Green's theorem]\label{d12:04:green-index}
Let \(D\) be an already specified bounded planar region whose finite
piecewise \(C^1\) oriented boundary cycle \(\Gamma\) satisfies the
hypotheses of the real Green theorem, also after removing a sufficiently
small closed interior disk. Then
\(\ind(\Gamma,a)=1\) for \(a\in D\) and \(0\) for
\(a\notin\overline D\), or for any other complementary component on
which the same Green calculation has no pole.
No existence or separation theorem for the interior of an arbitrary
Jordan trace is asserted here.
\end{corollary}
\begin{proof}
For a \(C^1\) holomorphic \(q=u+\iu v\), the coordinate contour formula
and the real Green theorem give
\[
 \int_{\partial D}q\dd z
 =\iint_D(-v_x-u_y)\dd x\dd y+
       \iu\iint_D(u_x-v_y)\dd x\dd y=0
\]
by CR, whenever \(q\) is holomorphic on a neighborhood of the closed
integration region. Take \(q(z)=1/(z-a)\).
For \(a\notin\overline D\) it has no pole and the integral is zero.
For \(a\in D\), remove a small disk centered at \(a\). Green on the
punctured region gives the integral over \(\Gamma\) equal to the
positive small-circle integral, namely \(2\pi\iu\).
The same zero calculation applies if \(a\) lies in a hole of the
specified region. Only the declared, already available real Green
theorem for such regions is used.
\end{proof}

\begin{theorem}[Full finite-cycle Cauchy theorem and formula]\family{D12.04.05}\label{d12:04-05}
Let \(f\) be holomorphic on an open set \(\Omega\), and let a finite
cycle \(\Gamma\) have its trace in \(\Omega\). If
\(\ind(\Gamma,w)=0\) for every \(w\notin\Omega\), then
\[
 \int_\Gamma f(z)\dd z=0,\qquad
 \ind(\Gamma,a)f(a)=\frac1{2\pi\iu}
       \int_\Gamma\frac{f(z)}{z-a}\dd z\quad(a\in\Omega\setminus|\Gamma|).
\]
The derivative formula is
\[
 \ind(\Gamma,a)f^{(n)}(a)
 =\frac{n!}{2\pi\iu}\int_\Gamma
            \frac{f(z)}{(z-a)^{n+1}}\dd z.
\]
Neither simple connectedness nor a simple trace is required.
\end{theorem}
\begin{proof}
We give the global gluing argument, including its diagonal regularity.
For \(z,w\in\Omega\) let
\(q(z,w)=(f(w)-f(z))/(w-z)\), with \(q(z,z)=f'(z)\).
It is jointly continuous and holomorphic in \(z\).
Off the diagonal this follows by quotient rules.
Near a diagonal point choose a circle containing both \(z,w\) in its
disk and contained in \(\Omega\). The local Cauchy formula gives
\[
 q(z,w)=\frac1{2\pi\iu}\int_C
        \frac{f(\zeta)}{(\zeta-z)(\zeta-w)}\dd\zeta .
\]
It holds for distinct \(z,w\) by subtraction of the two formulas and
at equality by the derivative formula. Uniform denominator bounds
prove the claimed continuity and complex differentiation in \(z\),
with continuous derivative, on this neighborhood.

Let \(H=\{z\notin|\Gamma|:\ind(\Gamma,z)=0\}\).
It is open, and \(\Omega\cup H=\C\).
Define
\[
 Q(z)=
 \begin{cases}
 \displaystyle\int_\Gamma q(z,w)\dd w,&z\in\Omega,\\[3pt]
 \displaystyle\int_\Gamma f(w)(w-z)^{-1}\dd w,&z\in H .
 \end{cases}
\]
On the overlap the two values differ by
\(f(z)\int_\Gamma(w-z)^{-1}\dd w=0\), so they agree.
The second formula is holomorphic by the uniform rational-kernel
argument used in the disk theorem.
For the first, fix a small disk about \(z_0\) compact in \(\Omega\).
The joint regularity just established, compactness of \(|\Gamma|\)
and a finite subcover give uniform convergence of the \(z\)-difference
quotients of \(q(z,w)\) for \(w\) on that trace.
The ML estimate permits differentiation of the contour integral.
Thus \(Q\) is entire.

Outside a large circle every point is in \(H\), and the second formula
tends to zero by the ML estimate, since \(f\) is bounded on the trace.
It is therefore bounded on the exterior and, by continuity, on the
remaining compact disk. The preliminary Liouville lemma makes \(Q=0\).
Its first definition now yields the Cauchy formula.

If \(\Gamma\) is nonzero, its finite union of finite-length traces has
empty interior by \cref{d12:03:trace-empty-interior}. Choose
\(a\in\Omega\setminus|\Gamma|\).
Apply the established formula to \((z-a)f(z)\); its value at \(a\)
is zero and its kernel integral is \(\int_\Gamma f\), proving the
theorem. The zero cycle is immediate.
On a neighborhood avoiding the trace, the index is constant. Differentiate
the formula there using the already justified kernel derivatives;
this gives the \(n\)-th derivative formula.
\end{proof}
\source{CON IV.4--5, printed 80--87, PDF 94--101, especially the global gluing proof at printed 84--85; AHL IV.4, printed 137--148, PDF 153--164.}

\begin{theorem}[Homotopic Cauchy theorem and simple connectivity]\family{D12.04.06}\label{d12:04-06}
Integrals of a holomorphic function agree along piecewise \(C^1\) paths
that are fixed-endpoint homotopic in its open domain. They also agree
for freely homotopic closed loops. A null-homotopic closed path has zero
integral. Consequently every holomorphic function on a simply connected
domain has a primitive, and every zero-free holomorphic function there
has a holomorphic logarithm and an \(n\)-th root for every \(n\ge1\).
\end{theorem}
\begin{proof}
The grid proof in \cref{d12:04-04} establishes exactly the two homotopy
assertions, including continuous intermediate paths.
Compare a null-homotopic loop with the constant loop.
For a simply connected domain every closed loop has such a contraction,
so the primitive criterion \cref{d12:03-04} applies.
The logarithm and root consequences follow from the period criterion
proved next, applied to the holomorphic function \(f'/f\).
The holomorphy of \(f'\) is already a local Cauchy conclusion.
\end{proof}

\begin{theorem}[Global logarithm and root criteria]\family{D12.04.07}\label{d12:04-07}
Let \(f\) be zero-free and holomorphic on a domain \(\Omega\).
It has a holomorphic logarithm if and only if
\(\int_\gamma f'/f=0\) on every closed piecewise \(C^1\) path.
It has a holomorphic \(n\)-th root, \(n\ge1\), if and only if
\[
 \frac1{2\pi\iu}\int_\gamma\frac{f'}f\dd z\in n\Z
 \quad\hbox{for every such closed path}.
\]
All logarithms differ by a constant \(2\pi\iu k\); all \(n\)-th roots
differ by a constant \(n\)-th root of unity.
\end{theorem}
\begin{proof}
For a logarithm \(L\), differentiation gives \(L'=f'/f\), so periods
vanish. Conversely zero periods give a primitive \(L_0\) of \(f'/f\).
The derivative of \(fE(-L_0)\) is zero, so it is a nonzero constant.
Choose a logarithm of that constant and add it to \(L_0\).

For the root assertion fix \(a\in\Omega\), choose \(c^n=f(a)\), and put
\[
 g(z)=c\exp\left(\frac1n\int_{a}^{z}\frac{f'}f\dd w\right),
\]
using any polygonal path. Differences of path integrals are periods;
the stated condition and the exponential kernel make \(g\) unambiguous.
On a small disk it has a local primitive representation, so it is
holomorphic with \(g'/g=f'/(nf)\). Thus \(f/g^n\) is constant, and its
value at \(a\) is one. If \(f=g^n\) already, the path \(g\circ\gamma\)
is piecewise \(C^1\) and zero-free. The integer-index argument gives
\(\int_\gamma g'/g\in2\pi\iu\Z\), proving necessity.
Logarithm differences and root ratios are continuous with values in
the respective discrete sets. Connectedness makes them constant.
\end{proof}

\begin{theorem}[Morera and continuous gluing]\family{D12.04.08}\label{d12:04-08}
A continuous function on an open set is holomorphic if its integral
around every triangle compactly contained there is zero.
A continuous function holomorphic off one point, or off a locally
finite union of compact rectifiable arcs, is holomorphic everywhere.
The latter assertion is local: in each relatively compact small disk
only finitely many of the arcs are required.
\end{theorem}
\begin{proof}
On a disk the triangle criterion gives a holomorphic \(F\) with \(F'=f\).
The disk Cauchy derivative theorem proves \(F'\) holomorphic, hence \(f\).
For a continuous exceptional point use the earlier exceptional-point
lemma and this criterion.

For the arc assertion fix a closed triangle in a small disk where the
exceptional set is contained in finitely many rectifiable arcs of total
length \(L\). The arcs are compact, and their locally finite union is
relatively closed; consequently a closed triangle missing them has a
holomorphic neighborhood. Subdivide by a triangular grid of mesh \(h\).
The arclength covering in \cref{d12:03:trace-empty-interior} covers those
arcs by at most \(C+L/h\) disks of radius \(h\), where \(C\) is the
finite number of arcs. A disk of radius \(h\) meets a bounded number
of grid triangles, independently of \(h\); this follows by scaling the
fixed nondegenerate triangular grid, whose vertices have a fixed
positive minimum separation.
Thus the total perimeter of triangles meeting the exceptional set is
at most \(C_1(L+Ch)\). All other small triangles have zero integrals
by Goursat. On a bad triangle subtract \(f\) at one of its vertices;
the constant's boundary integral is zero, and the remaining integral
is at most its perimeter times \(\omega_f(C_2h)\).
Interior sides cancel in the sum over every triangle, so the original
boundary integral is at most
\(C_1(L+Ch)\omega_f(C_2h)\), tending to zero.
Morera applies. Local disks covering the open set prove the full claim.
This includes straight segments and regular analytic arcs without a
separate topological separation assumption.
\end{proof}
\source{CON IV.5--6, printed 86--97, PDF 100--111; AHL IV.6.5, printed 172--174, PDF 188--190. The finite-length gluing argument is the explicit subdivision proof above.}


% END chapters/04


% BEGIN chapters/05
\chapter{Taylor series, zeros and local maps}\label{ch:d12-05}

\begin{theorem}[Taylor representation and smoothness]\family{D12.05.01}\label{d12:05-01}
If \(f\) is holomorphic on \(B(a,R)\), \(0<R\le\infty\), then
\[
 f(z)=\sum_{n\ge0}\frac{f^{(n)}(a)}{n!}(z-a)^n,\qquad |z-a|<R.
\]
The series and all its termwise derivatives converge absolutely uniformly
on compact subdisks. Holomorphic functions are real \(C^\infty\);
\(\partial_x f=f'\), \(\partial_y f=\iu f'\), with the corresponding
iterated identities. Thus the equivalences stated in
\cref{d12:01-04-equivalence} are now proved.
\end{theorem}
\begin{proof}
Choose \(|z-a|<r<R\) and use the circle Cauchy formula.
On that circle the geometric expansion
\[
 (w-z)^{-1}=\sum_{n\ge0}(z-a)^n(w-a)^{-n-1}
\]
is uniformly convergent, uniformly also for \(z\) in any smaller disk.
Its integral coefficients are \(f^{(n)}(a)/n!\) by
\cref{d12:04-02}. This proves the equality. The same estimate gives
absolute normal convergence, or it follows from \cref{d12:02-01}.
As \(r\) can approach \(R\), the asserted disk is obtained, including
\(R=\infty\). The power-series calculus proves all derivative assertions.
The differential matrix in \cref{d12:01-02} gives the two first real
partial identities; applying it to every holomorphic derivative gives
every higher mixed derivative and its continuity.
\end{proof}

\begin{theorem}[Coefficient and remainder bounds]\family{D12.05.02}\label{d12:05-02}
If \(M_r=\max_{|w-a|=r}|f(w)|\), where \(r<R\), then
\[
 |a_n|\le M_r r^{-n},\qquad
 \left|f(z)-\sum_{n=0}^Na_n(z-a)^n\right|
 \le M_r\frac{(|z-a|/r)^{N+1}}{1-|z-a|/r}
 \quad(|z-a|<r).
\]
For every \(k\ge0\), center derivatives obey
\(|f^{(k)}(a)|\le k!M_r/r^k\).
\end{theorem}
\begin{proof}
The coefficient formula is the Cauchy integral divided by \(2\pi\iu\).
Its ML bound is \(M_r/r^n\). Sum the geometric majorant over
\(n\ge N+1\) for the remainder. Multiplication by \(k!\) gives the
derivative estimate.
\end{proof}
\source{CON IV.2, printed 68--76, PDF 82--90; AHL V.1.2, printed 179--184, PDF 195--200.}

\begin{theorem}[Local zero factorization]\family{D12.05.03}\label{d12:05-03}
At each point \(a\), a holomorphic function either vanishes on a
neighborhood or has a unique integer \(m\ge0\) and a holomorphic \(g\)
near \(a\) with \(g(a)\ne0\), such that
\(f(z)=(z-a)^m g(z)\). If \(f(a)=0\), then \(m\ge1\) is its zero
multiplicity. A finite-order zero is isolated. Multiplicities add under
products of functions with finite orders.
\end{theorem}
\begin{proof}
If every Taylor coefficient is zero the Taylor representation vanishes
near \(a\). Otherwise choose the least nonzero coefficient \(a_m\)
and divide the series by \((z-a)^m\).
The remaining series has constant coefficient \(a_m\ne0\), hence is
nonzero on a smaller disk by continuity. The coefficient formula
makes \(m\) unique. Multiplication of the two factorizations gives
addition of the orders.
\end{proof}

\begin{theorem}[Identity theorem]\family{D12.05.04}\label{d12:05-04}
On a domain, a holomorphic function with zeros accumulating at a point
of the domain is identically zero. Two holomorphic functions equal on
a set with an interior accumulation point are therefore equal everywhere.
\end{theorem}
\begin{proof}
A zero accumulation at \(a\) excludes a finite-order factorization there,
so the function vanishes on a neighborhood of \(a\).
Let \(E\) consist of points with a zero neighborhood. It is nonempty
and open. If \(b\) is in its closure within the domain, either it is
already in \(E\) or distinct zero points approach \(b\). In the latter
case a finite-order factorization is again impossible, so \(b\in E\).
Thus \(E\) is also relatively closed; connectedness gives the whole
domain. Apply this result to the difference of the two functions.
\end{proof}

\begin{corollary}[Discrete and countable zero sets]\label{d12:05-03-discrete}
A holomorphic function on a domain, not identically zero, has isolated
zeros, finitely many in each compact subset, and at most countably many
in total. These claims are componentwise on a disconnected open set,
and do not apply to components on which the function vanishes identically.
\end{corollary}
\begin{proof}
The identity theorem excludes the zero-neighborhood alternative at every
point, leaving finite-order isolated zeros.
An infinite subset of a compact planar set has an accumulation point,
which would contradict that theorem. To prove countability, for each zero
choose a disk with rational center coordinates and rational radius which
contains that zero and no other zero. Among an enumeration of all such
disks assign the first permissible disk. Two distinct zeros cannot receive
the same disk, giving an injection into a countable set.
\end{proof}
\source{CON IV.3, printed 76--80, PDF 90--94; AHL IV.3.2, printed 126--130, PDF 142--146.}

\begin{lemma}[Inverse at a noncritical point]\label{d12:05:inverse-noncritical}
If \(f'(a)\ne0\), \(f\) restricts to a biholomorphism between
neighborhoods of \(a\) and \(f(a)\), and its inverse \(h\) satisfies
\(h'(w)=1/f'(h(w))\).
\end{lemma}
\begin{proof}
Taylor smoothness and \(\det D_{\R}f(a)=|f'(a)|^2>0\) permit application
of the real \(C^1\) inverse theorem. It supplies a real \(C^1\)
local inverse. Its differential is the inverse of the complex-linear
differential of \(f\), hence is multiplication by \(1/f'\).
The real/complex derivative bridge \cref{d12:01-02} makes that inverse
holomorphic, with the displayed derivative.
\end{proof}

\begin{theorem}[Local mapping normal form]\family{D12.05.05}\label{d12:05-05}
If \(f\) is holomorphic near \(a\) and not constant on any neighborhood
of \(a\), there are \(m\ge1\), \(r>0\), a neighborhood \(U\) of \(a\),
and a biholomorphism \(\phi:U\to B(0,r)\) with \(\phi(a)=0\), such that
\[
 f(z)-f(a)=\phi(z)^m .
\]
Its image is \(B(f(a),r^m)\); every value there other than \(f(a)\)
has exactly \(m\) distinct preimages in \(U\).
\end{theorem}
\begin{proof}
Factor \(f(z)-f(a)=(z-a)^m g(z)\), \(g(a)\ne0\).
A local \(m\)-th root \(h\) of \(g\) exists by
\cref{d12:02-05}. Set \(\phi(z)=(z-a)h(z)\).
Its derivative at \(a\) is \(h(a)\ne0\).
The inverse lemma makes it a biholomorphism on a smaller neighborhood;
shrink its target to a disk and restrict its source accordingly.
For a nonzero \(w\) with \(|w|<r^m\), its \(m\) roots are
\[
 \exp\bigl((\log|w|+\iu\theta+2\pi\iu k)/m\bigr),
 \quad 0\le k<m,\quad w=|w|e^{\iu\theta}.
\]
The exponential kernel proves these are all distinct and all the roots.
Their moduli are \(|w|^{1/m}<r\). Taking their inverse images under
\(\phi\) proves the exact multiplicity assertion, including the image.
\end{proof}

\begin{theorem}[Open maps and holomorphic inverses]\family{D12.05.06}\label{d12:05-06}
A nonconstant holomorphic function on a domain is open.
An injective holomorphic function on a domain has nonzero derivative
everywhere, is a homeomorphism onto its open image, and has a
holomorphic inverse. A nonzero derivative alone guarantees only local
invertibility.
\end{theorem}
\begin{proof}
For any \(a\) and any source neighborhood, the identity theorem makes
\(f-f(a)\) nonzero on every sufficiently small neighborhood.
Apply the normal form inside that neighborhood: its image contains a
disk about \(f(a)\). This proves openness.
If an injective function had local multiplicity \(m>1\), the last
theorem would give distinct preimages of nearby values, a contradiction.
Thus \(m=1\) and \(f'(a)\ne0\). A continuous injective open map has a
continuous inverse on its image. The local inverse lemma shows that
inverse is holomorphic near every point, hence everywhere by locality.
For the final boundary, \(E(z)\) has derivative \(E(z)\ne0\) but takes
the same values at \(z\) and \(z+2\pi\iu\).
\end{proof}

\begin{proposition}[Polynomial and analytic multiplicity]\family{D12.05.07}\label{d12:05-07}
For a nonzero polynomial \(P\), the analytic multiplicity of a root
\(a\) equals the largest integer \(m\) such that \((z-a)^m\) divides
\(P\) algebraically.
\end{proposition}
\begin{proof}
Repeated Euclidean division by \(z-a\) gives
\(P=(z-a)^mQ\), \(Q(a)\ne0\), and stops after at most \(\deg P\) steps.
This is the local holomorphic factorization with a nonzero remaining
factor, whose order is unique. No fundamental theorem of algebra is
assumed in this argument.
\end{proof}
\source{CON IV.7, printed 97--100, PDF 111--114; AHL IV.3.3, printed 130--133, PDF 146--149.}


% END chapters/05

% BEGIN chapters/06
\chapter{Maximum principles and growth barriers}\label{ch:d12-06}

\begin{theorem}[Maximum and minimum modulus]\family{D12.06.01}\label{d12:06-01}
On a domain, an interior local maximum of the modulus of a holomorphic
function forces the function to be constant. If it is zero-free, the
same holds for an interior local minimum of its modulus.
On a bounded domain, a function continuous on the closure and holomorphic
inside has modulus bounded by its boundary maximum.
More generally, the boundary hypothesis
\[
 \limsup_{\Omega\ni z\to\xi}|f(z)|\le M
 \quad(\xi\in\partial\Omega),\qquad M\ge0,
\]
already implies \(|f|\le M\) in a bounded domain.
\end{theorem}
\begin{proof}
If a function on a domain is nonconstant, \cref{d12:05-06} says that
the image of every neighborhood is open. Any open neighborhood of a
complex number contains a value of strictly larger modulus, even when
that number is zero. This excludes a local maximum. Apply it to \(1/f\)
for the minimum assertion.

For the limsup form fix \(\eps>0\). Each boundary point has a neighborhood
where interior values are at most \(M+\eps\).
A finite collection covers the compact boundary.
The portion of the bounded domain outside their union is a compact
subset of the domain. If some value exceeded \(M+\eps\), the global
maximum would therefore be attained on that interior compact set.
The strong principle would make \(f\) constant, and its boundary limsup
would contradict that strict excess. Let \(\eps\downarrow0\).
The continuous-closure version follows from this result, or by applying
it with \(M=\max_{\partial\Omega}|f|\).
\end{proof}

\begin{theorem}[Liouville and algebraic closure]\family{D12.06.02}\label{d12:06-02}
Every bounded entire function is constant. Every nonconstant complex
polynomial has a root, and every polynomial of degree \(n\ge1\) is a
product of \(n\) linear factors, counted with their analytic/algebraic
multiplicities.
\end{theorem}
\begin{proof}
Liouville was proved from the disk derivative estimate in
\cref{d12:04:liouville}; that is the proof used here.
Suppose a polynomial \(P\) of positive degree had no root.
Its leading term gives \(|P(z)|\to\infty\) as \(|z|\to\infty\):
divide by the leading monomial, and the remaining finite inverse-power
sum tends to one. Thus \(1/P\) is entire and bounded outside a disk.
It is bounded on the remaining compact disk because its denominator
never vanishes. Liouville makes \(1/P\), hence \(P\), constant, a
contradiction. Divide out one linear factor and repeat until the degree
is zero. The multiplicity comparison \cref{d12:05-07} identifies the
repeated factors with analytic multiplicity.
\end{proof}
\source{CON VI.1, printed 128--130, PDF 142--144; AHL IV.2.3 and IV.3.4, printed 120--124, 133--137, PDF 136--140, 149--153.}

\begin{theorem}[Schwarz lemma and all equality cases]\family{D12.06.03}\label{d12:06-03}
If \(f:\D\to\D\) is holomorphic and \(f(0)=0\), then
\(|f(z)|\le|z|\) and \(|f'(0)|\le1\).
If equality holds at a nonzero point, or \(|f'(0)|=1\), then
\(f(z)=cz\), \(|c|=1\). Conversely these rotations give equality.
\end{theorem}
\begin{proof}
The Taylor series removes the singularity of \(g(z)=f(z)/z\), setting
\(g(0)=f'(0)\). On \(|z|=r<1\), \(|g(z)|\le1/r\).
The boundary maximum theorem gives this bound on the radius-\(r\)
disk. For a fixed interior \(z\), let \(r\uparrow1\), obtaining
\(|g(z)|\le1\), also at zero. Equality at any of the stated points
is an interior maximum of \(|g|\), so \(g\) is constant of modulus one.
The converse is a direct calculation.
\end{proof}

\begin{theorem}[Disk automorphisms and Schwarz--Pick]\family{D12.06.04}\label{d12:06-04}
For \(a\in\D\), the map
\(\phi_a(z)=(z-a)/(1-\bar a z)\) is a disk automorphism with inverse
\((w+a)/(1+\bar a w)\).
Every holomorphic \(f:\D\to\D\) satisfies
\[
 \left|\frac{f(z)-f(a)}{1-\overline{f(a)}f(z)}\right|
 \le\left|\frac{z-a}{1-\bar a z}\right|,\qquad
 \frac{|f'(a)|}{1-|f(a)|^2}\le\frac1{1-|a|^2}.
\]
Equality in the first inequality at \(z\ne a\), or in the derivative
inequality at any point, holds exactly for disk automorphisms.
All such automorphisms are \(c\phi_a\), \(|c|=1\).
\end{theorem}
\begin{proof}
The denominators are nonzero in the disk. Expansion gives
\[
 1-|\phi_a(z)|^2=
 \frac{(1-|a|^2)(1-|z|^2)}{|1-\bar a z|^2}>0.
\]
Solving for \(z\) gives the stated holomorphic inverse, which obeys
the same disk inequality. Apply Schwarz to
\(h=\phi_{f(a)}\circ f\circ\phi_a^{-1}\).
It vanishes at zero. The Schwarz value bound is exactly the first
inequality; its derivative bound, using
\(\phi_a'(a)=1/(1-|a|^2)\) and
\((\phi_a^{-1})'(0)=1-|a|^2\), is the second.
Either equality makes \(h\) a rotation, hence \(f\) an automorphism.
Conversely an automorphism and its inverse give opposite inequalities,
so both are equalities.
For an arbitrary automorphism choose its unique preimage \(a\) of zero.
The automorphism \(f\circ\phi_a^{-1}\) fixes zero; applying Schwarz to
it and its inverse gives derivative equality, hence it is a rotation.
\end{proof}
\source{CON VI.2, printed 130--133, PDF 144--147; AHL IV.3.4, printed 133--137, PDF 149--153.}

\begin{theorem}[Hadamard three circles]\family{D12.06.05}\label{d12:06-05}
Let \(f\) be holomorphic on \(r_0<|z|<r_3\) and not identically zero.
For \(r_0<r_1<r<r_2<r_3\), put
\(M(t)=\max_{|z|=t}|f(z)|>0\).
Then
\[
 \log M(r)\le
 \frac{\log(r_2/r)}{\log(r_2/r_1)}\log M(r_1)
 +\frac{\log(r/r_1)}{\log(r_2/r_1)}\log M(r_2).
\]
Thus \(\log M(t)\) is convex as a function of \(\log t\).
The identically zero function satisfies the corresponding multiplicative
inequality with zero values.
\end{theorem}
\begin{proof}
Set
\(A=(\log M(r_2)-\log M(r_1))/\log(r_2/r_1)\).
For positive integers \(n\), choose integers \(k_n\) with
\(k_n/n\to A\). On the closed annulus \(r_1\le|z|\le r_2\),
the holomorphic function \(f(z)^n z^{-k_n}\) is continuous.
Its maximum is bounded by the larger boundary maximum, so
\[
 \log M(r)-\frac{k_n}{n}\log r
 \le\max_{j=1,2}\left(\log M(r_j)-\frac{k_n}{n}\log r_j\right).
\]
Both entries on the right tend to the same value
\(\log M(r_1)-A\log r_1\). Taking limits and rearranging gives the formula.
If some \(M(t)\) were zero, all points on that circle would be zeros,
forcing identically zero by the identity theorem; this proves the strict
positivity used in the logarithms. The zero function is handled before
taking logarithms.
\end{proof}

\begin{theorem}[A sector Phragm\'en--Lindel\"of theorem]\family{D12.06.06}\label{d12:06-06}
Let \(\alpha\ge1/2\) and
\[
 S_\alpha=\{re^{\iu\theta}:r>0,\ |\theta|<\pi/(2\alpha)\}.
\]
For \(\alpha=1/2\), use \(-\pi<\theta<\pi\), so the negative ray is
the cut. Suppose \(f\) is holomorphic there,
\(\limsup_{S_\alpha\ni z\to\xi}|f(z)|\le M\) at every finite
boundary point, including zero, and \(M\ge0\).
If for some \(A>0,B\ge0,0\le\beta<\alpha\),
\[
 |f(z)|\le A\exp(B|z|^\beta)
 \quad(z\in S_\alpha,\ |z|\hbox{ sufficiently large}),
\]
then \(|f|\le M\) throughout the sector.
The strict growth threshold cannot be replaced by \(\beta=\alpha\).
\end{theorem}
\begin{proof}
Choose \(\gamma\) with \(\beta<\gamma<\alpha\).
The principal logarithm on the sector defines \(z^\gamma\), and
\[
 \Ree z^\gamma\ge c|z|^\gamma,\qquad
 c=\cos(\gamma\pi/(2\alpha))>0 .
\]
For \(\eps>0\), let \(F_\eps(z)=f(z)\exp(-\eps z^\gamma)\).
Along each boundary ray its limsup is at most \(M\), since the multiplier
has modulus at most one. The boundary limsup at zero has the same bound.
On a large outer arc of radius \(R\),
\[
 |F_\eps|\le A\exp(BR^\beta-\eps cR^\gamma)\longrightarrow0.
\]
Given \(\delta>0\), choose such an outer radius with bound \(M+\delta\);
choose an inner radius so small that \(|F_\eps|\le M+\delta\)
on its arc, by the limsup at zero. The bounded-domain boundary-limsup
maximum principle on the truncated sector gives
\(|F_\eps(z)|\le M+\delta\) at every fixed point between these radii.
The radii can enclose any specified point, so let \(\delta\downarrow0\).
Then let \(\eps\downarrow0\) to obtain \(|f(z)|\le M\).
This argument also covers \(M=0\).

For sharpness, \(f(z)=\exp(z^\alpha)\), using the sector branch,
has modulus one on both rays and tends to one at the vertex.
It obeys growth \(\exp(|z|^\alpha)\), but is unbounded on the positive
axis. Thus the critical exponent does not give the same conclusion.
Rotated or translated sectors obey the identical result after their
explicit affine coordinate change.
\end{proof}
\source{CON VI.3--4, printed 133--140, PDF 147--154, especially the sector barrier at printed 139--140.}


% END chapters/06

% BEGIN chapters/07
\chapter{Laurent series, singularities and the sphere}\label{ch:d12-07}

\begin{theorem}[Laurent expansion]\family{D12.07.01}\label{d12:07-01}
Let \(0\le r_0<R_0\le\infty\), and let \(f\) be holomorphic on
\(A=\{r_0<|z-a|<R_0\}\). Then uniquely
\[
 f(z)=\sum_{n\in\Z}c_n(z-a)^n,\qquad
 c_n=\frac1{2\pi\iu}\int_{|w-a|=\rho}
                    \frac{f(w)}{(w-a)^{n+1}}\dd w
 \quad(r_0<\rho<R_0).
\]
The two one-sided series converge absolutely and uniformly on every
closed subannulus; both may be differentiated termwise any number of times.
The coefficient integral is independent of \(\rho\).
\end{theorem}
\begin{proof}
For two admissible radii the difference of their positive circles has
index zero outside \(A\), by the circle computation.
Apply full Cauchy to \(f(w)/(w-a)^{n+1}\), which is holomorphic on \(A\),
to prove coefficient independence.

For a fixed \(z\in A\) choose
\(r_0<\rho_1<|z-a|<\rho_2<R_0\).
The outer positive circle minus the inner positive circle has index
one at \(z\) and zero on \(\C\setminus A\).
The full Cauchy formula gives
\[
 f(z)=\frac1{2\pi\iu}\int_{|w-a|=\rho_2}\frac{f(w)}{w-z}\dd w
       -\frac1{2\pi\iu}\int_{|w-a|=\rho_1}\frac{f(w)}{w-z}\dd w.
\]
In the first integral expand in powers of \((z-a)/(w-a)\).
In the second use
\((w-z)^{-1}=-\sum_{k\ge0}(w-a)^k/(z-a)^{k+1}\).
Both geometric expansions are uniformly convergent on their respective
circles, and coefficient independence identifies all terms with the
displayed \(c_n\).

For a fixed subannulus \(s\le|z-a|\le t\), choose
\(\rho_1<s\le t<\rho_2\).
If \(M_j=\max_{|w-a|=\rho_j}|f(w)|\), then
\[
 |c_n|\le M_2\rho_2^{-n}\ (n\ge0),\qquad
 |c_{-k}|\le M_1\rho_1^k\ (k\ge1).
\]
Geometric majorants with ratios \(t/\rho_2\) and \(\rho_1/s\) give
absolute uniform convergence. Any fixed number of derivatives introduces
only polynomial factors in \(n,k\) and fixed powers of \(s^{-1}\);
the same convergent geometric majorants with these polynomial factors
prove uniform convergence of the derivative series. The argument in
\cref{d12:02-01} justifies differentiation.
Finally, integrate a proposed normally convergent Laurent series against
\((w-a)^{-m-1}\) on one circle. Each term vanishes except \(n=m\),
whose integral is \(2\pi\iu\), proving uniqueness.
\end{proof}
\source{CON V.1, printed 106--109, PDF 120--123; AHL V.1.3, printed 184--187, PDF 200--203.}

\begin{definition}[Isolated singularities]\family{D12.07.02}\label{d12:07-02}
An isolated singularity at \(a\) means that \(f\) is holomorphic on
\(0<|z-a|<R\). It is removable if \(f\) extends holomorphically to \(a\);
a pole of order \(m\ge1\) if \(f=(z-a)^{-m}g\), with \(g\) holomorphic
and \(g(a)\ne0\); and essential if neither alternative holds.
A general domain boundary point need not have a punctured neighborhood
and is not covered by this terminology.
\end{definition}
\begin{theorem}[Coefficient and growth classification]\label{d12:07-02-classification}
In the punctured-disk Laurent expansion:
removability means all negative coefficients vanish;
a pole of order \(m\) means \(c_{-m}\ne0\) and \(c_n=0\) for \(n<-m\);
essentiality means infinitely many negative coefficients are nonzero.
For an integer \(m\ge0\), the bound \(f(z)=O(|z-a|^{-m})\) means at
most a pole of order \(m\), including removability.
For \(m\ge1\), \(f(z)=o(|z-a|^{-m})\) means at most order \(m-1\).
A finite limit at \(a\) is equivalent to removability, and
\(|f(z)|\to\infty\) is equivalent to a pole.
\end{theorem}
\begin{proof}
If there are no negative coefficients, the positive series is a
holomorphic extension. Conversely the extension's Taylor series and
Laurent uniqueness exclude negative coefficients.
A finite nonzero negative principal part factors as
\((z-a)^{-m}g(z)\), \(g(a)=c_{-m}\ne0\); its converse follows by
expanding \(g\). The remaining case is precisely an infinite negative
part.

On the circle of radius \(\rho\), the coefficient bound is
\(|c_{-k}|\le\rho^k\max|f|\).
An \(O(\rho^{-m})\) bound makes \(c_{-k}=0\) for every \(k>m\) by
letting \(\rho\downarrow0\). An \(o(\rho^{-m})\) bound also makes
\(c_{-m}=0\). The corresponding factorizations give the converse bounds.

A finite limit gives a bounded punctured function and hence only
nonnegative coefficients by the case \(m=0\); its extension value is
the limit. A pole plainly tends to infinity in modulus.
If \(|f|\to\infty\), then on a smaller punctured disk \(1/f\) is
holomorphic, bounded and tends to zero. It extends holomorphically
with value zero, and cannot vanish identically because it is nonzero
off the center. Its zero has some finite order by the identity theorem,
so its reciprocal has a pole.
\end{proof}

\begin{theorem}[Riemann removability and Casorati--Weierstrass]\family{D12.07.03}\label{d12:07-03}
A bounded isolated singularity is removable; the weaker condition
\((z-a)f(z)\to0\) also suffices.
At an essential isolated singularity the image of every punctured
neighborhood is dense in \(\C\).
\end{theorem}
\begin{proof}
Boundedness is the \(m=0\) estimate above. Under the weaker condition,
\(\rho\max_{|z-a|=\rho}|f(z)|\to0\); the coefficient bound gives
\(c_{-k}=0\) for every \(k\ge1\), so the singularity is removable.
If an essential singularity had a punctured neighborhood whose image
missed \(B(w,\eps)\), then \(h=1/(f-w)\) would be bounded and
holomorphic there. It extends to \(a\). If \(h(a)\ne0\), then
\(f=w+1/h\) extends. If \(h(a)=0\), its finite-order zero gives a
pole of \(f\); \(h\) is not identically zero since it never vanishes
off \(a\). Both alternatives contradict essentiality.
\end{proof}
\source{CON V.1, printed 103--112, PDF 117--126; AHL IV.3.1--3.2, printed 124--130, PDF 140--146.}

\begin{definition}[Scalar meromorphic functions]\family{D12.07.04}\label{d12:07-04}
A scalar meromorphic function on an open set is holomorphic except at
a relatively closed discrete set of poles, with value \(\infty\) at those poles.
It is never identically infinity on a connected component.
At a point where it is neither identically zero locally nor a pole,
its signed order is its zero multiplicity; a pole of order \(m\) has
signed order \(-m\).
\end{definition}
\begin{proposition}[Local quotients and arithmetic]\label{d12:07-04-quotients}
Scalar meromorphic functions are exactly the local quotients \(p/q\)
of holomorphic functions with \(q\) not identically zero on the local
disk, after filling the cancellable singularities.
They are closed under addition, subtraction and multiplication;
the reciprocal exists meromorphically for a function not identically
zero on any component. On a domain these functions form a field.
For nonzero factors, signed orders add under multiplication and subtract
under division.
\end{proposition}
\begin{proof}
At a finite regular point take \(q=1\); at a pole take
\(q=(z-a)^m,\ p=(z-a)^mf\). Conversely, factor the finite-order zeros
of \(p,q\). If \(p\) vanishes identically, the quotient extends as zero.
Otherwise their quotient is
\((z-a)^{\ord_a p-\ord_a q}\) times a nonzero holomorphic factor,
so it has only a removable singularity or a pole at each zero of \(q\).
Those zeros are discrete. Local common-denominator formulas
\[
 p_1/q_1+p_2/q_2=(p_1q_2+p_2q_1)/(q_1q_2),
 \qquad (p_1/q_1)(p_2/q_2)=p_1p_2/(q_1q_2)
\]
prove closure. Denominators do not become identically zero, since a
nonzero holomorphic function has only isolated zeros on a disk.
For reciprocation exchange \(p,q\); a function not identically zero
has a nonzero local numerator by the identity theorem.
The same factorizations prove the order rules.
These formulas define the meromorphic operations even at apparent
expressions such as infinity minus infinity; they are not undefined
pointwise extended-number arithmetic.
\end{proof}

\begin{definition}[Sphere and the coordinate at infinity]\family{D12.07.05}\label{d12:07-05}
Let \(\Chat=\C\cup\{\infty\}\), with neighborhoods of infinity the sets
containing \(\{\infty\}\cup\{|z|>R\}\).
Its finite coordinate is \(z\); near infinity use \(w=1/z\), with
\(w(\infty)=0\). A scalar function is holomorphic or meromorphic at
infinity if its expression \(f(1/w)\) has that property at zero.
An isolated singularity at infinity requires holomorphy on some entire
exterior \(|z|>R\).
\end{definition}
\begin{proposition}[Compact sphere and entire functions at infinity]\label{d12:07-05-classification}
The sphere is compact, Hausdorff and connected, and its two coordinates
have holomorphic transition \(w=1/z\).
An entire function has a removable singularity at infinity precisely
when it is constant; it has a pole there of order \(m\) precisely when
it is a polynomial of degree \(m\ge1\). A nonpolynomial entire function
has an essential singularity at infinity.
For reduced \(P/Q\) with degrees \(p,q\), the behavior at infinity is
a pole of order \(p-q\) if \(p>q\), a zero of order \(q-p\) if \(p<q\),
and a finite nonzero value if \(p=q\).
\end{proposition}
\begin{proof}
An open cover has a member containing infinity and the exterior of a
large disk; finitely many other members cover the remaining compact
disk. This proves compactness. Two finite points can be separated by
disjoint small disks; a finite point and infinity are separated by a
bounded neighborhood and a sufficiently large exterior.
The connected plane is dense in \(\Chat\), so its closure is connected.
The two displayed coordinates and their transition are continuous
inverse homeomorphisms, and \(1/z\) is holomorphic off zero.

An entire function removable at infinity is bounded on a large exterior
and on the remaining compact disk, hence constant by Liouville.
If it has a pole of order \(m\) at infinity, \(|f(z)|\le C|z|^m\)
for large \(|z|\). Cauchy's coefficient estimate on radius \(R\) makes
every Taylor coefficient of degree \(n>m\) at most \(CR^{m-n}\),
hence zero as \(R\to\infty\). The resulting polynomial has degree
exactly \(m\), by the pole's exact order; the converse is obtained by
substitution \(z=1/w\).
The isolated-singularity trichotomy gives the nonpolynomial case.
Finally divide \(P/Q\) by their leading monomials:
\(P(z)/Q(z)=(a_p/b_q)z^{p-q}(1+o(1))\), or substitute \(1/w\)
and factor the exact powers of \(w\); this gives all three orders.
\end{proof}
\source{CON I.6, printed 8--11, PDF 22--25; V.1, printed 111, PDF 125; AHL I.2.4, printed 18--20, PDF 34--36.}

\begin{theorem}[Rationality and degree]\family{D12.07.06}\label{d12:07-06}
Every scalar meromorphic function on the whole compact sphere is rational.
For a nonconstant reduced rational function \(P/Q\),
\(d=\max(\deg P,\deg Q)\) is both the total number of its zeros and the
total number of its poles on the sphere, counted with multiplicity.
More generally every sphere value has \(d\) preimages counted with its
local multiplicities. A nonzero constant has no zeros or poles;
the identically zero function is excluded from finite zero-count claims.
\end{theorem}
\begin{proof}
Poles are isolated and have no accumulation on the sphere: at an ordinary
point continuity gives a pole-free neighborhood, and at a pole its
local factorization excludes other poles nearby. Compactness therefore
makes their number finite. Subtract the finite Laurent principal parts
at all finite poles. The difference \(h\) is entire and still
meromorphic at infinity, since local quotient arithmetic preserves
meromorphy. The preceding proposition makes \(h\) a polynomial or
constant. Adding back the finitely many rational principal parts proves
rationality.

For reduced \(P/Q\), finite zeros and poles number \(p=\deg P\) and
\(q=\deg Q\), by the FTA, with no cancellation. The infinity orders in
the preceding proposition make each total equal to \(\max(p,q)\).
For a finite target \(c\), preimages solve the nonzero polynomial
\(P-cQ=0\), of degree \(e\le d\).
If \(q<d\), then \(p=d\), \(e=d\), and infinity maps to infinity,
giving all \(d\) preimages finitely.
If \(q=d\), multiply numerator and denominator of
\(P(1/w)/Q(1/w)-c\) by \(w^d\).
The denominator has a nonzero value at zero and the numerator has a
zero of order \(d-e\). The \(e\) finite roots plus that infinity order
give \(d\), with \(d-e=0\) meaning infinity is not a preimage.
For the target infinity the total pole count gives \(d\).
\end{proof}

\begin{definition}[Holomorphic maps into the sphere]\family{D12.07.07}\label{d12:07-07}
A continuous map into \(\Chat\) is holomorphic when its expression in
a target coordinate is holomorphic near every source point; at an
infinite value this expression is \(1/F\), with value zero.
The source may be a plane open set, or a sphere open set using its
stated coordinates. Transition \(z\mapsto1/z\) proves independence
of these coordinates.
\end{definition}
\begin{theorem}[The scalar/sphere-valued bridge]\label{d12:07-07-bridge}
On a connected plane or sphere domain, the holomorphic sphere-valued
maps are exactly the scalar meromorphic functions with their pole values,
together with the constant infinity map.
A finite-valued holomorphic function on the whole sphere is constant.
\end{theorem}
\begin{proof}
For a scalar meromorphic function, an ordinary point uses its usual
holomorphic expression. At a pole \(a\), reciprocal factorization is
holomorphic and zero; this proves continuity into infinity and
holomorphy in its infinity coordinate.
Conversely near any point mapped to infinity the reciprocal coordinate
is holomorphic. Unless it vanishes identically there, its zero is of
finite order and gives an isolated pole of the original scalar expression.
If \(F\) is infinity on an open neighborhood, that property extends
through every boundary point of the set of such neighborhoods:
continuity permits an infinity chart near that point, and its
holomorphic reciprocal vanishes on a nonempty open subset, hence on
the entire small disk by the identity theorem.
The set is therefore both open and relatively closed; connectedness
makes \(F\) identically infinity. In all other cases its infinite
values are isolated poles and its finite values are holomorphic in
the ordinary chart, proving the bridge. The coordinate transition
\(1/z\), holomorphic with holomorphic inverse off zero, justifies each
source or target coordinate replacement.

For a finite-valued holomorphic function on the compact sphere, its
continuous modulus attains a maximum. In a source coordinate this is an
interior local maximum. The strong maximum principle makes the function
constant on a neighborhood; the identity theorem, continued through
overlapping coordinates on the connected sphere, makes it constant
everywhere.
\end{proof}
\source{CON V.3, printed 123--128, PDF 137--142, and VII.3, printed 155--157, PDF 169--171; AHL II.1.4, printed 30--33, PDF 46--49; the coordinate and compactness bridges are proved above.}


% END chapters/07

% BEGIN chapters/08
\chapter{Residues and controlled integral applications}\label{ch:d12-08}

\begin{definition}[Residue]\family{D12.08.01}\label{d12:08-01}
For an isolated singularity with Laurent expansion
\(f(z)=\sum_{n\in\Z}c_n(z-a)^n\), set
\(\Res_a(f(z)\dd z)=c_{-1}\), also written \(\Res(f;a)\).
Equivalently it is \((2\pi\iu)^{-1}\int_C f\dd z\) on any sufficiently
small positive circle about \(a\).
\end{definition}
\begin{proposition}[Computation and coordinate invariance]\label{d12:08-01-formulas}
If \(f\) has a pole of order at most \(m\ge1\), then
\[
 \Res(f;a)=\frac1{(m-1)!}
  \left.\frac{\dd^{\,m-1}}{\dd z^{m-1}}
             ((z-a)^m f(z))\right|_{z=a}.
\]
For \(f=p/q\), with \(p,q\) holomorphic,
\(q(a)=0,\ q'(a)\ne0\), the residue is \(p(a)/q'(a)\);
if \(p(a)=0\) this may be a removable singularity.
If \(z=\phi(w)\), \(\phi(b)=a,\ \phi'(b)\ne0\), is a local
holomorphic coordinate, then
\[
 \Res_b\bigl(f(\phi(w))\phi'(w)\dd w\bigr)
       =\Res_a(f(z)\dd z).
\]
\end{proposition}
\begin{proof}
The coefficient \(c_{-1}\) is the coefficient of degree \(m-1\)
in the Taylor series of \((z-a)^mf\), giving the derivative formula.
Factor \(q(z)=(z-a)q_1(z)\), \(q_1(a)=q'(a)\), for the simple formula.

For invariance take a small circle \(C_b\) on which \(\phi\) is a local
biholomorphism and its image lies in the punctured Laurent disk.
Uniform convergence of the Laurent series on that compact image permits
termwise integration. For \(n\ne-1\), the composed term
\((\phi-a)^n\phi'\) has primitive
\((\phi-a)^{n+1}/(n+1)\); only the \(n=-1\) term remains.
Write \(\phi(w)-a=(w-b)h(w)\), where \(h\) is holomorphic and nonzero
on a sufficiently small disk. Its local logarithm gives
\[
 \frac{\phi'}{\phi-a}=\frac1{w-b}+\frac{h'}h,
\]
whose circle integral is \(2\pi\iu\). Thus the transformed differential's
circle integral is \(2\pi\iu c_{-1}\), proving invariance even for an
essential isolated singularity.
The integral characterization itself follows by integrating the normally
convergent Laurent series on its circle.
\end{proof}

\begin{theorem}[Residue theorem for cycles]\family{D12.08.02}\label{d12:08-02}
Let \(\Omega\) be open, \(A\subset\Omega\) a relatively closed discrete
set, and \(f\) holomorphic on \(\Omega\setminus A\).
Let a finite cycle \(\Gamma\) avoid \(A\), with trace in \(\Omega\),
and suppose \(\ind(\Gamma,w)=0\) for \(w\notin\Omega\).
Then
\[
 \int_\Gamma f\dd z
   =2\pi\iu\sum_{a\in A}\ind(\Gamma,a)\Res(f;a).
\]
Only finitely many terms in the sum can be nonzero.
\end{theorem}
\begin{proof}
First suppose \(A=\{a_1,\ldots,a_m\}\) is finite.
Choose pairwise disjoint small disks around these points, with closures
in \(\Omega\) and missing the trace. Let \(C_j\) be their positive
circles and \(n_j=\ind(\Gamma,a_j)\).
The cycle \(\Gamma-\sum_jn_jC_j\) has index zero outside
\(\Omega\setminus A\): outside \(\Omega\) every small circle has
index zero, and at each \(a_j\) its correction cancels \(n_j\).
Full Cauchy on that punctured open set gives the result, since
\(\int_{C_j}f=2\pi\iu\Res(f;a_j)\).

For the discrete case let
\[
 K=|\Gamma|\cup
   \{z\notin|\Gamma|:\ind(\Gamma,z)\ne0\}.
\]
Its complement is open by local constancy. The index vanishes outside
a large disk, so \(K\) is compact; the complement-index hypothesis
also makes \(K\subset\Omega\).
Choose a small \(\delta>0\) with the closed \(\delta\)-neighborhood
of \(K\) compact in \(\Omega\).
Only finitely many members of the relatively closed discrete set \(A\)
lie in that compact neighborhood: otherwise an accumulation point
there would contradict discreteness and relative closedness.
Let \(U=\{z:\dist(z,K)<\delta\}\).
The cycle has index zero outside \(U\), and \(f\) has only finitely
many singularities in \(U\). Apply the finite case on \(U\).
Every other \(a\in A\) has zero index by the definition of \(K\).
\end{proof}
\source{CON V.2, printed 112--114, PDF 126--128; AHL IV.5.1, printed 148--152, PDF 164--168.}

\begin{definition}[Residue at infinity]\family{D12.08.03}\label{d12:08-03}
For a differential \(f(z)\dd z\) holomorphic on a large punctured
neighborhood of infinity, define
\[
 \Res_\infty(f(z)\dd z)
   =\Res_{w=0}\bigl(-f(1/w)w^{-2}\dd w\bigr).
\]
It is minus the coefficient of \(z^{-1}\) in the exterior Laurent
expansion of \(f\).
\end{definition}
\begin{theorem}[Sum on the sphere]\label{d12:08-03-sum}
A meromorphic differential on the sphere, expressed as \(f(z)\dd z\)
in the finite coordinate and by the transformed expression at infinity,
has total residue zero.
\end{theorem}
\begin{proof}
The Jacobian of \(z=1/w\) is \(-w^{-2}\), which gives the definition
and the coefficient statement directly.
Its finitely many finite poles are enclosed by a large positive circle.
The residue theorem makes that circle's integral \(2\pi\iu\) times
their sum. In the \(w\)-coordinate the circle runs clockwise about zero.
Changing parameter therefore identifies its negative with the positive
small-circle integral of \(-f(1/w)w^{-2}\dd w\).
This gives minus the residue at infinity as the sum of finite residues.
A meromorphic differential has finitely many poles by compactness
and its local Laurent descriptions; hence no infinite sum was used.
\end{proof}

\begin{theorem}[Rational and trigonometric integrals]\family{D12.08.04}\label{d12:08-04}
Let \(R=P/Q\) be a rational function in reduced form, with no real
poles and \(\deg Q\ge\deg P+2\). Then its ordinary improper integral
is absolutely convergent and
\[
 \int_{-\infty}^{\infty}R(x)\dd x
    =2\pi\iu\sum_{\Imm a>0}\Res(R;a).
\]
If a rational expression \(S(\cos t,\sin t)\) has no poles for real
\(t\), then
\[
 \int_0^{2\pi}S(\cos t,\sin t)\dd t
 =2\pi\iu\sum_{|a|<1}\Res(T;a),\quad
 T(z)=\frac1{\iu z}
 S\!\left(\frac{z+z^{-1}}2,\frac{z-z^{-1}}{2\iu}\right),
\]
after meromorphic cancellation.
In particular, for \(a>|b|\),
\[
 \int_0^{2\pi}\frac{\dd t}{a+b\cos t}
       =\frac{2\pi}{\sqrt{a^2-b^2}},\qquad
 \int_{-\infty}^{\infty}\frac{\dd x}{1+x^2}=\pi .
\]
\end{theorem}
\begin{proof}
The degree gap gives \(|R(z)|\le C|z|^{-2}\) uniformly outside
a large disk. On the real line this proves absolute convergence of
both tails; compact intervals have no poles.
Integrate along \([-r,r]\) followed by the upper semicircle.
Its index is one at the enclosed upper poles and zero at all other
poles, by \cref{d12:04:crossing-index}.
The arc integral is at most \(\pi r\,Cr^{-2}\to0\).
The residue theorem proves the first formula.

For the second formula use \(z=e^{\iu t}\),
\(\dd t=\dd z/(\iu z)\) and the two displayed expressions for sine
and cosine. The resulting rational function has no unit-circle pole;
its possible pole at zero is included. The circle residue theorem
gives the identity.
For the stated example \(b=0\) gives \(2\pi/a\) directly.
If \(b\ne0\), the transformed differential is
\(2\,\dd z/(\iu(bz^2+2az+b))\).
Its roots are
\(\alpha=(-a+\sqrt{a^2-b^2})/b\) and
\(\beta=(-a-\sqrt{a^2-b^2})/b\).
Their product is one and
\(|\alpha|=|b|/(a+\sqrt{a^2-b^2})<1<|\beta|\).
The residue at \(\alpha\) is
\(2/(\iu b(\alpha-\beta))=1/(\iu\sqrt{a^2-b^2})\).
This yields the value. For \(1/(1+z^2)\) the only upper pole is
\(\iu\), with residue \(1/(2\iu)\).
\end{proof}
\source{AHL IV.5.3, printed 154--157, PDF 170--173; CON V.2, printed 114--123, PDF 128--137.}

\begin{theorem}[Jordan estimate and oscillatory rational integrals]\family{D12.08.05}\label{d12:08-05}
Let \(c>0\), and let \(q\) be continuous on the upper semicircle \(C_r\).
Then
\[
 \left|\int_{C_r}e^{\iu cz}q(z)\dd z\right|
 \le\frac{\pi}{c}\sup_{C_r}|q|.
\]
Consequently this arc tends to zero if its supremum tends to zero.
For a proper rational function \(R\), without real poles, and \(c>0\),
the two one-sided ordinary improper tails exist and
\[
 \int_{-\infty}^{\infty}e^{\iu cx}R(x)\dd x
   =2\pi\iu\sum_{\Imm a>0}\Res(e^{\iu cz}R(z);a).
\]
For \(c<0\) use the lower poles and the factor \(-2\pi\iu\).
If the degree gap is at least two, the integrals are also absolutely
convergent. For \(b>0,\ c\in\R\),
\[
 \int_{-\infty}^{\infty}\frac{\cos(cx)}{x^2+b^2}\dd x
       =\frac{\pi}{b}e^{-b|c|}.
\]
\end{theorem}
\begin{proof}
Parameterize \(z=re^{\iu t}\). Since
\(\sin t\ge 2t/\pi\) for \(0\le t\le\pi/2\), symmetry gives
\[
 r\int_0^\pi e^{-cr\sin t}\dd t
 \le2r\int_0^{\pi/2}e^{-2crt/\pi}\dd t
 \le\pi/c,
\]
proving the estimate. A proper rational function satisfies
\(R(z)=O(1/z)\), uniformly outside a large disk, so the arc vanishes.
To justify ordinary tails rather than only symmetric cutoffs, note that
\(R'(x)=O(x^{-2})\), from its convergent inverse-power expansion at
infinity (or by polynomial division of its derivative).
Integration by parts on \([A,B]\) gives
\[
 \int_A^B e^{\iu cx}R(x)\dd x
 =\left[\frac{e^{\iu cx}R(x)}{\iu c}\right]_A^B
   -\frac1{\iu c}\int_A^B e^{\iu cx}R'(x)\dd x.
\]
The endpoint terms have limits and the derivative integral is absolutely
convergent; the negative tail is proved with the same formula.
Residues on the upper half-disk and the vanishing arc now give the
identity. Replacing \(z\) by \(-z\) proves the lower-plane case and its
orientation sign. The stronger degree gap gives absolute convergence
by the preceding theorem.
For \(c>0\), the residue of
\(e^{\iu cz}/(z^2+b^2)\) at \(\iu b\) is
\(e^{-cb}/(2\iu b)\). The imaginary real-axis integrand is odd and
integrable, leaving the cosine formula. Evenness in \(c\) gives
\(c<0\), and the rational formula at \(c=0\) gives the remaining case.
\end{proof}

\begin{theorem}[Indentations and principal values]\family{D12.08.06}\label{d12:08-06}
Let \(R\) be proper rational with only simple real poles, and \(c>0\).
At each real pole take symmetric deleted intervals, independently
shrinking to zero. The tails at plus and minus infinity are taken as
ordinary one-sided limits. Then
\[
 \operatorname{PV}\int_{-\infty}^{\infty}e^{\iu cx}R(x)\dd x
 =2\pi\iu\sum_{\Imm a>0}\Res(e^{\iu cz}R(z);a)
   +\pi\iu\sum_{a\in\R}\Res(e^{\iu cz}R(z);a).
\]
In particular
\[
 \operatorname{PV}\int_{\R}\frac{e^{\iu cx}}x\dd x=\pi\iu,\qquad
 \int_0^\infty\frac{\sin(cx)}x\dd x=\frac\pi2 .
\]
The latter integral is ordinary improper but not absolutely convergent;
the former is not a licensed absolutely integrable scalar integral.
\end{theorem}
\begin{proof}
Near a simple real pole \(a\), write
\(e^{\iu cz}R(z)=A_a/(z-a)+h_a(z)\), with \(h_a\) holomorphic.
A small upper semicircle, followed clockwise from \(a-\eps\) to
\(a+\eps\), contributes \(-\pi\iu A_a+o(1)\):
the principal term integrates exactly and the remaining term is bounded
by its length times a fixed local bound.
The upper half-disk boundary with these semicircular indentations
has index one about upper poles and zero about the removed real poles.
This follows also by subtracting the positively oriented small
half-disk boundary cycles from the large half-disk cycle; the deleted
real segments cancel. For a point in a deleted open upper half-disk
the two indices are both one, giving zero for the indented cycle.
The real pole is in that same complementary component of the indented
trace, so its index is zero by local constancy; no individual unindented
index is evaluated at a point on its trace. Apply the residue theorem, then the Jordan
estimate and the small-arc limits. Rearrangement gives the displayed
sign and half-residue factors.

Existence of the independent local principal values follows directly
from the same decomposition: the symmetric integrals of
\(A_a/(x-a)\) cancel, and \(h_a\) is integrable locally.
The two far tails exist by the integration-by-parts argument above.
For \(R(z)=1/z\), there are no upper poles and the real residue is one.
The imaginary part has removable value \(c\) at zero.
Its tails converge by integration by parts, and it is even, giving
the half-line value \(\pi/2\).
Absolute convergence fails: on each interval
\([(k+1/6)\pi/c,(k+5/6)\pi/c]\), \(|\sin(cx)|\ge1/2\);
the integrals of \(|\sin(cx)|/x\) on those disjoint intervals dominate
a positive multiple of \(\sum_{k\ge1}1/k\).
The complex integrand also has a nonintegrable modulus at zero.
\end{proof}
\source{AHL IV.5.3, printed 156--159, PDF 172--175; CON V.2, printed 114--123, PDF 128--137.}

\begin{theorem}[A branch-cut integral]\family{D12.08.07}\label{d12:08-07}
For real \(0<s<1\), the ordinary integral is absolutely convergent and
\[
 \int_0^\infty\frac{x^{s-1}}{1+x}\dd x=\frac{\pi}{\sin\pi s}.
\]
\end{theorem}
\begin{proof}
The integrand is bounded by \(x^{s-1}\) near zero and by \(x^{s-2}\)
near infinity, proving convergence.
Use the branch \(z^{s-1}=\exp((s-1)(\log|z|+\iu\arg z))\),
\(0<\arg z<2\pi\), with cut along the positive ray.
Integrate \(z^{s-1}/(1+z)\) around the positively oriented boundary
of \(\eps<|z|<R,\ \delta<\arg z<2\pi-\delta\), where
\(0<\eps<1<R\) and \(0<\delta<\pi/2\).
This is a piecewise smooth cycle in the branch domain. Its index at
\(-1\) is one: the negative ray starting at \(-1\) crosses the outer
circle once positively and no other boundary piece, by the crossing
formula. The only enclosed pole has residue \(e^{\iu\pi(s-1)}\).

Let \(\delta\downarrow0\) with \(\eps,R\) fixed. The radial integrals
converge uniformly on their compact radial intervals; the upper bank
runs outwards, the lower bank inwards and acquires multiplier
\(e^{2\pi\iu(s-1)}=e^{2\pi\iu s}\).
The inner arc has modulus at most
\(2\pi\eps^s/(1-\eps)\), tending to zero.
The outer arc is at most \(2\pi R^s/(R-1)\), tending to zero since
\(s<1\). Thus the residue equation becomes
\[
 (1-e^{2\pi\iu s})\int_0^\infty\frac{x^{s-1}}{1+x}\dd x
       =2\pi\iu e^{\iu\pi(s-1)}.
\]
Since \(1-e^{2\pi\iu s}=-2\iu e^{\iu\pi s}\sin\pi s\),
division gives the result. Both strict endpoint conditions have been
used; the real integrals diverge at \(s=0\) and \(s=1\).
\end{proof}

\begin{theorem}[A residue summation formula]\label{d12:08-07-summation}
For real \(a>0\),
\[
 \sum_{n\in\Z}\frac1{n^2+a^2}=\frac{\pi}{a}\coth(\pi a).
\]
The bilateral sum is absolutely convergent.
\end{theorem}
\begin{proof}
Integrate
\(F(z)=\pi\cot(\pi z)/(z^2+a^2)\)
around the square with sides
\(\Ree z,\Imm z=\pm(N+1/2)\), choosing \(N+1/2>a\).
At each integer \(n\) inside it, \(\pi\cot(\pi z)\) has residue one:
\(\sin\pi z\) has derivative \(\pi(-1)^n\) and numerator
\(\pi\cos\pi z\) has value \(\pi(-1)^n\).
The two additional residues at \(\pm\iu a\) are each
\(-\pi\coth(\pi a)/(2a)\), since
\(\cot(\iu t)=-\iu\coth t\).

Let \(r=N+1/2\). On the vertical sides the exponential expression
\(\cot w=\iu(e^{2\iu w}+1)/(e^{2\iu w}-1)\) gives
\(|\cot\pi z|\le1\). On the horizontal sides it gives
\(|\cot\pi z|\le(1+e^{-2\pi r})/(1-e^{-2\pi r})
 \le\coth(\pi/2)\).
Moreover \(|z^2+a^2|\ge |z|^2-a^2\ge r^2-a^2\).
The contour integral is therefore bounded in modulus by
\(8r\pi\coth(\pi/2)/(r^2-a^2)\), tending to zero.
Its index at each interior pole is one by the explicit rectangle
computation. The residue theorem yields
\[
 \lim_{N\to\infty}
 \left(\sum_{n=-N}^N\frac1{n^2+a^2}
            -\frac{\pi}{a}\coth(\pi a)\right)=0.
\]
Comparison with \(\sum_{n\ge1}n^{-2}\) proves absolute convergence,
so this symmetric limit is the ordinary bilateral sum and may be
reordered. No unproved cotangent partial-fraction expansion is used.
\end{proof}
\source{AHL IV.5.3, printed 159--160, PDF 175--176; CON V.2, printed 112--123, PDF 126--137. The keyhole bank limits and the square-contour summation estimates are supplied explicitly above.}


% END chapters/08


% BEGIN chapters/09
\chapter{Argument principle, Rouch\'e and stability of zeros}\label{ch:d12-09}

Orders of zeros and poles are the signed Laurent orders of Chapter~\ref{ch:d12-07}. A cycle below is finite, avoids every zero and pole under discussion, and has index zero on the complement of the stated open set. Sums with index weights are finite: the closure of the nonzero-index set together with the trace is compact in that open set. Indeed the index vanishes outside a large disk, and a fixed neighborhood of the complement misses the trace and has index zero by local constancy.

\begin{proposition}[Logarithmic derivative]\family{D12.09.01}\label{d12:09-01}
If a nonzero meromorphic function has signed order $m$ at $a$, then $f'/f$ has residue $m$ there. Its only poles are zeros and poles of $f$, and they are simple unless their residue is zero.
\end{proposition}
\begin{proof}
Write $f(z)=(z-a)^m u(z)$ with $u$ holomorphic and nonzero near $a$. Differentiation on the punctured disk gives $f'/f=m/(z-a)+u'/u$. The last term is holomorphic throughout the disk. This proves every assertion, including negative $m$ and $m=0$.
\end{proof}
\source{CON V.3 pp.123--124 (PDF137--138); AHL IV.5.2 pp.152--153.}

\begin{theorem}[Argument principle]\family{D12.09.02}\label{d12:09-02}
For a meromorphic $f$ on $\Omega$, nonzero on every component relevant to a cycle $\Gamma$, with no zero or pole on its trace,
\[
 \frac1{2\pi\iu}\int_\Gamma\frac{f'}f\,dz
 =\sum_{a}\ind(\Gamma,a)\ord_a f.
\]
The left side is $\ind(f\circ\Gamma,0)$, where images retain their parametrizations and integer coefficients. For an already specified bounded region $D$ satisfying the real Green theorem hypotheses (a $C^1$ boundary or the stated finite piecewise regular decomposition), its full positive boundary gives the ordinary number of zeros minus poles in $D$. This includes a Jordan-domain case when that given domain and boundary satisfy those hypotheses; it does not assert an interior for an arbitrary Jordan trace.
\end{theorem}
\begin{proof}
Apply the finite-cycle residue theorem to $f'/f$, using \cref{d12:09-01}. On every smooth piece the chain rule gives $d(f\circ\gamma)/(f\circ\gamma)=(f'/f)(\gamma)\gamma' dt$; add the pieces. For the last assertion use \cref{d12:04:green-index}: for $a\in D$ apply real Green to $D$ with a small disk about $a$ removed, using the closed form $dz/(z-a)$. Its inner-circle term is $-2\pi\iu$, so the full boundary integral is $2\pi\iu$. For $a\notin\overline D$ apply Green directly to get zero. Thus its boundary index is one on the given region and zero outside its closure. No general Jordan separation theorem is invoked.
\end{proof}
\source{CON V.3 pp.123--128 (PDF137--142); AHL IV.5.2 pp.152--154.}

\begin{theorem}[Rouch\'e]\family{D12.09.03}\label{d12:09-03}
Let $f,g$ be holomorphic on a neighborhood of the trace and all nonzero-index points of a finite cycle $\Gamma$. Suppose $|g-f|<|f|$ on its trace. Their index-weighted numbers of zeros agree. In particular their ordinary zero counts agree for the positive boundary of a given Green-regular region, or for a positive circle, under the preceding hypotheses.
\end{theorem}
\begin{proof}
Set $f_t=f+t(g-f)$, $0\le t\le1$. On the trace $|f_t|\ge |f|-|g-f|>0$, with a common positive lower bound by compactness. Thus $t\mapsto(2\pi\iu)^{-1}\int_\Gamma f_t'/f_t$ is continuous: the integrands are jointly continuous and uniformly bounded on each compact parameter interval and contour piece. It is integer-valued by the argument principle, hence constant on $[0,1]$. Its endpoint values give the assertion. The strict inequality also excludes the identically zero endpoint on any component having nonzero index.
\end{proof}
\source{CON V.3 pp.125--128 (PDF139--142); SS3.4 pp.89--93 (PDF108--112).}

\begin{corollary}[Perturbations and polynomial roots]\family{D12.09.04}\label{d12:09-04}
If $f$ has no zero on a compact counting contour, every holomorphic perturbation within $\min|f|$ there has the same count. A monic polynomial of degree $n$ has exactly $n$ complex roots with multiplicity. Simple roots of a polynomial remain in any chosen disjoint sufficiently small disks about them under sufficiently small changes of its coefficients; there is exactly one root in each such disk. More generally a zero of multiplicity $m$ produces exactly $m$ perturbed zeros in its isolating disk.
\end{corollary}
\begin{proof}
The first assertion is \cref{d12:09-03}. For $P(z)=z^n+\sum_{j<n}c_jz^j$, choose $R>1+\sum|c_j|$. On $|z|=R$, the lower terms have modulus less than $R^n$, so the count equals that of $z^n$, namely $n$. No zero lies outside this disk, since the same inequality holds for $|z|\ge R$. Polynomial and analytic multiplicity agree by local factorization, giving the FTA bridge without changing the meaning of multiplicity. Choose each isolating circle with no other zero. Coefficient convergence is uniform on the finitely many circles, since each monomial is bounded there; apply the first assertion on each circle. For a general holomorphic perturbation the same argument uses its uniform boundary error directly.
\end{proof}
\source{CON V.3; AHL IV.5.2; polynomial/analytic multiplicity bridge \cref{d12:05-07}.}

\begin{theorem}[Hurwitz]\family{D12.09.05}\label{d12:09-05}
Let holomorphic $f_n$ converge locally uniformly on a domain to a holomorphic $f\not\equiv0$. On any isolating circle with closed disk in the domain and $f\ne0$ on the circle, $f_n$ and $f$ eventually have the same number of zeros. If every $f_n$ is zero-free, the limit is zero-free or identically zero.
\end{theorem}
\begin{proof}
The boundary minimum of $|f|$ is positive, so local uniform convergence and Rouch\'e give the count. If a nonzero limit had a zero, isolated-zero factorization would supply a small circle containing a positive number of zeros and no boundary zero, contradicting zero-freeness of $f_n$.
\end{proof}
\source{CON VII.2 pp.152--153 (PDF166--167); SS8.3.2--3.3 pp.227--228 (PDF246--247).}

\begin{theorem}[Limits of univalent functions]\family{D12.09.06}\label{d12:09-06}
A locally uniform holomorphic limit of injective holomorphic functions on a domain is either injective or constant.
\end{theorem}
\begin{proof}
Suppose $f$ is nonconstant and $f(a)=f(b)$ with $a\ne b$. The functions $h_n(z)=f_n(z)-f_n(a)$ converge locally uniformly to $h=f-f(a)$. Choose a small closed disk about $b$ missing $a$, with $h$ nonzero on its boundary. Hurwitz gives a zero of $h_n$ inside that disk for large $n$. Injectivity says its only zero in the domain is $a$, a contradiction. The constant case cannot be removed: $f_n(z)=z/n$ on a disk gives such a limit.
\end{proof}
\source{CON VII.2 Exercise10 p.154 (PDF168), proved here; SS8 Proposition3.5 pp.227--228 (PDF246--247).}

% END chapters/09

% BEGIN chapters/10
\chapter{M\"obius geometry and explicit conformal maps}\label{ch:d12-10}

\begin{proposition}[Angles and orientation]\family{D12.10.01}\label{d12:10-01}
At a point where $f'(a)=c\ne0$, the real differential of a holomorphic $f$ is multiplication by $c$: a rotation followed by dilation $|c|$, with determinant $|c|^2>0$. It preserves oriented angles of nonzero tangent vectors. The map is a conformal equivalence locally; a nonzero derivative everywhere does not assert global injectivity.
\end{proposition}
\begin{proof}
The real differential is $h\mapsto ch$ by \cref{d12:01-02}. In real coordinates its matrix is $\left(\begin{smallmatrix}\Ree c&-\Imm c\\\Imm c&\Ree c\end{smallmatrix}\right)$. For tangents $v,w\ne0$, their image ratio is $cw/(cv)=w/v$, proving angle preservation, and the determinant is $|c|^2$. The inverse theorem \cref{d12:05-06} gives the local assertion. The exponential has nonzero derivative and repeats with period $2\pi\iu$, proving the limitation.
\end{proof}

\begin{theorem}[M\"obius transformations]\family{D12.10.02}\label{d12:10-02}
For $ad-bc\ne0$, $T(z)=(az+b)/(cz+d)$ extends to a biholomorphism of $\Chat$, with $T(-d/c)=\infty$ and $T(\infty)=a/c$ when $c\ne0$, and $T(\infty)=\infty$ when $c=0$. These maps form a group under composition; matrices represent the same map exactly when they differ by a nonzero scalar. Given two ordered triples of distinct sphere points there is a unique such map taking the first triple to the second. Cross ratios are preserved, including their limits when an entry is infinity.
\end{theorem}
\begin{proof}
The inverse matrix $\left(\begin{smallmatrix}d&-b\\-c&a\end{smallmatrix}\right)$ gives the inverse map. In finite charts the derivative is $(ad-bc)/(cz+d)^2$; reciprocal charts at a pole and at infinity show the extensions and nonzero chart derivatives. Matrix multiplication gives composition. If two matrices give the same map, their quotient gives the identity, and $(az+b)/(cz+d)=z$ on an open set forces $c=b=0$ and $a=d\ne0$.
For distinct finite $p,q,r$, the map
\[
 A(z)=\frac{(z-p)(q-r)}{(z-r)(q-p)}
\]
takes them to $0,1,\infty$; its limits give the formula if an entry is infinity. Compose the map for the first triple with the inverse for the second. A map fixing $0,1,\infty$ has $b=c=0$ and $a=d$, proving uniqueness. Define the cross ratio $[z,p;q,r]=A(z)$. Composing with any M\"obius map gives two normalized three-point maps, so uniqueness proves invariance, with the same chart limits in infinity cases.
\end{proof}
\source{CON III.3 pp.44--58 (PDF58--72); AHL III.3 pp.76--86.}

\begin{proposition}[Circles, reflections and automorphisms]\family{D12.10.03}\label{d12:10-03}
M\"obius maps take generalized circles (circles or lines with infinity) to generalized circles. Reflection in one is the conjugate of $z\mapsto\bar z$ by a map taking it to the extended real line; it is independent of that map. Disk automorphisms are
$e^{\iu\theta}(z-a)/(1-\bar a z)$, $|a|<1$. Upper-half-plane automorphisms are real-coefficient fractions with positive determinant, modulo nonzero real scalar.
\end{proposition}
\begin{proof}
A generalized circle has equation $\alpha|z|^2+\beta z+\bar\beta\bar z+\gamma=0$, with real $\alpha,\gamma$ and $|\beta|^2-\alpha\gamma>0$. Substitution of the inverse fraction and multiplication by its denominator modulus squared gives another equation of this form. The discriminant is multiplied by a positive determinant modulus square; the case of a line is $\alpha=0$, interpreted at infinity. A circle can be taken to the real line by first sending one of its points to infinity and then rotating/translating its image line. Two choices differ by a map preserving the real line. Such a map has real coefficients up to scalar: normalize the images of three distinct real points by the real versions of the preceding formula and use uniqueness. It therefore commutes with conjugation, proving independence of reflection.
The disk classification is the equality case of Schwarz's lemma after sending $a=T^{-1}(0)$ to zero, as proved in \cref{d12:06-04}. Conjugate by $C(z)=(z-\iu)/(z+\iu)$ for the half-plane. The resulting fractions preserve the extended real line, hence have real coefficients up to scalar, and
$\Imm T(z)=(ad-bc)\Imm z/|cz+d|^2$ fixes precisely the positive sign. Conversely this formula and the inverse prove that every such real fraction is a half-plane automorphism.
\end{proof}

\begin{proposition}[Explicit injective maps]\family{D12.10.04}\label{d12:10-04}
The Cayley map $C$ above maps $\Imm z>0$ biholomorphically onto $\D$. The exponential maps $\{0<\Imm z<\theta\}$ onto $\{r e^{\iu t}:r>0,0<t<\theta\}$ for $0<\theta\le2\pi$, with the endpoint rays excluded, and inverse the logarithm with argument in $(0,\theta)$. On a sector of opening $\theta<2\pi$, that logarithm gives a power map of exponent $\pi/\theta$ onto a half-plane after rotation. Finally
$J(z)=(z+z^{-1})/2$ is a biholomorphism $\{|z|>1\}\to\C\setminus[-1,1]$.
\end{proposition}
\begin{proof}
The identity $1-|C(z)|^2=4\Imm z/|z+\iu|^2$ gives the Cayley image; its inverse is $\iu(1+w)/(1-w)$. Exponential modulus and argument give surjectivity on the stated sector, and equality of exponential values means the real parts agree and imaginary parts differ by $2\pi k$. The strict strip width at most $2\pi$ forces $k=0$. The logarithm is holomorphic with derivative $1/w$, so these are biholomorphisms. Multiplying its imaginary part interval by $\pi/\theta$ proves the power-map assertion with its chosen branch.
For $J(z)=w$ the roots satisfy $z^2-2wz+1=0$, with product one. A unit-modulus root would give $w=\cos t\in[-1,1]$. If $w$ is off that segment, neither root has modulus one, so exactly one has modulus greater than one. The quadratic thus proves both surjectivity and injectivity. Its derivative $(1-z^{-2})/2$ has no zero in the exterior, so the inverse is holomorphic. Equivalently the inverse is $w+\sqrt{w^2-1}$ with the branch asymptotic to $w$ at infinity; this branch is obtained from the proved inverse, not assumed globally beforehand.
\end{proof}
\source{SS8.1 pp.206--212 (PDF225--231); AHL III.4 pp.86--100.}

\begin{theorem}[Length and area]\family{D12.10.05}\label{d12:10-05}
For an absolutely continuous path $\gamma$ with compact trace in the domain of a holomorphic $f$,
$L(f\circ\gamma)=\int |f'(\gamma(t))|\,|\gamma'(t)|\,dt$.
For an injective holomorphic $f:\Omega\to f(\Omega)$ and a Lebesgue measurable $E\subset\Omega$,
\[
 m_2(f(E))=\int_E |f'|^2\dd m_2.
\]
More generally, for nonnegative measurable $h$ on $f(\Omega)$ the integral over $f(E)$ equals $\int_E h(f(z))|f'(z)|^2\dd m_2(z)$; the scalar identity holds when the absolute integral is finite.
\end{theorem}
\begin{proof}
On a neighborhood of the compact trace, $f$ is $C^1$ with bounded differential, and is Lipschitz on finitely many small convex neighborhoods. The AC chain rule applies on a finite subdivision of the path; the AC length formula and $|(f\circ\gamma)'|=|f'(\gamma)||\gamma'|$ give the first assertion. Injectivity implies $f'\ne0$ by the local multiplicity theorem. Hence $f$ is a real $C^1$ diffeomorphism between the two open plane sets, and its absolute real Jacobian is $|f'|^2$ by \cref{d12:10-01}. The D11 $C^1$ diffeomorphism change-of-variables theorem on completed Lebesgue spaces applies with these exact hypotheses, proving the nonnegative identity and image measurability; apply it to $h=1$ for area and to absolute values and real/imaginary parts for the integrable scalar version. For a merely rectifiable parametrization first use its arc-length/AC representation from \cref{d12:03-06}; the derivative formula is not assigned to an arbitrary non-AC parametrization.
\end{proof}
\source{SS8.4.3 p.238 (PDF257); D11 Chapter9 $C^1$ Jacobian/change-of-variables theorem and Chapter12 AC chain/length bridge.}

\begin{corollary}[Specified classical domains]\family{D12.10.06}\label{d12:10-06}
Every disk and every half-plane is conformally equivalent to $\D$. A horizontal strip of finite positive width, every sector of opening in $(0,2\pi)$, and a plane slit along one ray are conformally equivalent to $\D$. The exterior disk and the complement of a closed segment are equivalent to one another, not asserted equivalent to $\D$.
\end{corollary}
\begin{proof}
Affine maps reduce disks and half-planes to their unit models. Translate/rescale a strip to $0<\Imm z<\pi$, apply exponential onto the upper half-plane, then Cayley. Rotate a sector, use its selected power to the half-plane, then Cayley. A slit plane has argument interval of length $2\pi$; its selected square root gives a half-plane, with inverse squaring on that half-plane. The last equivalence is $J$ and affine changes of its segment endpoints from \cref{d12:10-04}. Every composition has the explicitly proved inverse on its stated image.
\end{proof}

% END chapters/10

% BEGIN chapters/11
\chapter{Compact convergence and normal families}\label{ch:d12-11}

\begin{definition}[The compact-open metric]\family{D12.11.01}\label{d12:11-01}
For an open $\Omega\subset\C$ choose
$K_j=\{z\in\Omega:|z|\le j,\ \dist(z,\C\setminus\Omega)\ge1/j\}$, with distance to the empty set interpreted as infinity. Empty suprema are zero. On continuous complex-valued functions put
\[
 d(f,g)=\sum_{j\ge1}2^{-j}\min(1,\sup_{K_j}|f-g|).
\]
Its convergence is exactly uniform convergence on compact subsets. The space $\Hol(\Omega)$ is complete for this metric.
\end{definition}
\begin{proof}
The $K_j$ are compact, $K_j\subset\operatorname{int}K_{j+1}$, and every compact subset of $\Omega$ lies in one $K_j$, by boundedness and positive distance from the complement. The inequality $\min(1,a+b)\le\min(1,a)+\min(1,b)$ proves the triangle inequality. Vanishing distance implies equality at every point because the $K_j$ exhaust $\Omega$. For convergence in the metric, each fixed nonnegative summand tends to zero, giving uniform convergence on $K_j$; conversely split the sum at a large index and use the bound on its geometric tail. A metric-Cauchy sequence is uniformly Cauchy on each $K_j$. Completeness of $\C$ gives compatible uniform limits there, and the interior nesting makes the assembled limit continuous. It is holomorphic by \cref{d12:11-02}. This also treats $\Omega=\varnothing$, whose function space has a single member.
\end{proof}

\begin{theorem}[Holomorphic limits and all derivatives]\family{D12.11.02}\label{d12:11-02}
A locally uniform limit of holomorphic functions is holomorphic, and every derivative converges locally uniformly to the corresponding derivative of the limit.
\end{theorem}
\begin{proof}
The limit is continuous. Its integral about each triangle in a disk compactly contained in $\Omega$ is the limit of zero integrals, so Morera proves holomorphy there and thus throughout $\Omega$. On $\overline{B(a,r)}\subset B(a,R)\Subset\Omega$, Cauchy's formula gives
\[
 \sup_{|z-a|\le r}|f_n^{(k)}(z)-f^{(k)}(z)|
 \le \frac{k!R}{(R-r)^{k+1}}\sup_{|w-a|=R}|f_n(w)-f(w)|.
\]
The right side tends to zero. Finitely many such smaller disks cover any given compact set, proving the assertion for every integer $k\ge0$.
\end{proof}
\source{CON VII.2 pp.151--152 (PDF165--166); SS2.5.2 pp.53--55 (PDF72--74).}

\begin{lemma}[The compactness mechanism]\family{D12.11.04}\label{d12:11-04}
A locally bounded family of holomorphic functions is equicontinuous on compact sets and every sequence has a subsequence converging uniformly on every compact subset. For continuous maps into a compact metric space, local equicontinuity alone gives such a subsequence.
\end{lemma}
\begin{proof}
If $|f|\le M$ on $B(a,2r)$, the Cauchy estimate gives $|f'|\le M/r$ on $B(a,r)$. Integrating over a segment in this smaller disk gives a common Lipschitz estimate. Finitely many such neighborhoods cover a compact set; a Lebesgue number for that finite cover gives uniform equicontinuity there. Choose a countable dense subset of $\Omega$ (rational-coordinate points suffice). Successive subsequence extraction at its first, second, and later points, followed by a diagonal, gives convergence at every point of this set: use boundedness in $\C$, or compactness of the target. On any compact $K$, equicontinuity on a compact neighborhood of $K$ and a finite net from the dense set make this subsequence uniformly Cauchy on $K$: compare two values at $z$ through a net point, with errors $\eps,\eps,\eps$. Completeness of the target gives the uniform limit. The same diagonal works on every $K_j$, hence on all compact sets. For complex-valued holomorphic maps the limit is holomorphic by \cref{d12:11-02}. This proves the precise Arzel\`a--Ascoli bridge used here rather than assuming compactness of the function space.
\end{proof}

\begin{theorem}[Finite-valued Montel]\family{D12.11.03}\label{d12:11-03}
A family $\mathcal F\subset\Hol(\Omega)$ is locally bounded if and only if every sequence in it has a locally uniformly convergent subsequence with finite holomorphic limit. Equivalently its closure in $\Hol(\Omega)$ is compact.
\end{theorem}
\begin{proof}
The forward implication was proved in \cref{d12:11-04}. If local boundedness fails, there are a compact $K$, $f_n\in\mathcal F$ and $z_n\in K$ with $|f_n(z_n)|>n$. No subsequence can converge uniformly on $K$ to a finite continuous function, which proves the reverse implication. Sequential compactness and compactness agree in metric spaces: failure of total boundedness produces a separated sequence, whereas a finite $2^{-j}$-net at each scale gives Cauchy subsequences by nested extraction; completeness then supplies limits. Apply this metric argument to the closure, approximating its $n$th chosen point by a family member within $1/n$.
\end{proof}
\source{CON VII.2 pp.153--154 (PDF167--168); SS8.3.2 pp.225--227 (PDF244--246).}

\begin{theorem}[Vitali convergence]\family{D12.11.05}\label{d12:11-05}
Let $\Omega$ be a domain and $(f_n)$ locally bounded in $\Hol(\Omega)$. If $f_n$ converges at every point of a set having an accumulation point in $\Omega$, it converges locally uniformly on $\Omega$ to a holomorphic function.
\end{theorem}
\begin{proof}
Montel supplies at least one subsequential holomorphic limit $f$. Every other subsequential holomorphic limit agrees with it on the specified set, and hence everywhere by the identity theorem. If the original sequence did not converge to $f$ uniformly on a compact $K$, a subsequence would have error at least some $\eps>0$ there. Montel would give a further subsequence with holomorphic limit, necessarily $f$, contradicting that error. Connectedness is used exactly in the identity theorem.
\end{proof}

\begin{theorem}[Spherical limits]\family{D12.11.06}\label{d12:11-06}
Use the chordal metric
\[
 \chi(z,w)=\frac{|z-w|}{\sqrt{1+|z|^2}\sqrt{1+|w|^2}},\qquad
 \chi(z,\infty)=\frac1{\sqrt{1+|z|^2}}
\]
on $\Chat$. It is a compact complete metric with the sphere topology. A locally uniform spherical limit of meromorphic maps on a domain is meromorphic or identically infinity. If the original functions are holomorphic and finite, the limit is finite holomorphic or identically infinity. The space of meromorphic maps together with the constant infinity is complete in the compact-open chordal metric. Spherical normality does not imply local ordinary boundedness.
\end{theorem}
\begin{proof}
Stereographic projection identifies $\chi$ with half the Euclidean chord distance on the unit sphere: its coordinate formula is
$(2\Ree z,2\Imm z,|z|^2-1)/(1+|z|^2)$, with infinity the north pole. The metric and compactness assertions follow from this formula and compactness of the sphere. The maps
$T_a(w)=(w-a)/(1+\bar a w)$, together with reciprocal charts for $a=\infty$, are chordal isometries, as direct substitution shows.
Let $f_n\to f$ spherically. Continuity of $f$ and uniform convergence give a small disk about each point on which $T_{f(a)}f_n$ are eventually bounded by one fixed finite constant: the image of $f$ is contained in a chordal ball whose closure misses the pole of this chart. There the chart transforms are finite holomorphic and converge ordinarily, because on a bounded set $|u-v|\le(1+M^2)\chi(u,v)$. Thus the limit is a holomorphic sphere-valued map in every chart. Unless it is identically infinity, its infinity fibers are isolated by the identity theorem in a reciprocal chart, and reciprocal zeros are precisely scalar poles. If $f_n$ are finite holomorphic and $f(a)=\infty$, on a small disk $1/f_n$ are holomorphic, zero-free and converge ordinarily to $1/f$. Hurwitz makes this limit zero-free or zero throughout the disk. Its zero at $a$ forces the latter, and chart identity propagates $f\equiv\infty$ on the domain. Completeness follows by applying compact metric completeness uniformly on $K_j$, as in \cref{d12:11-01}, and the limit classification. Finally the constants $f_n=n$ converge spherically to infinity uniformly everywhere, while their ordinary values are unbounded.
\end{proof}
\source{CON VII.3 pp.155--157 (PDF169--171), including the constant-infinity convention.}

\begin{theorem}[Marty's criterion]\family{D12.11.07}\label{d12:11-07}
For a meromorphic map put $f^\#=|f'|/(1+|f|^2)$ off its poles, and use $(1/f)^\#$ at a pole. This is continuous and well defined. A family of meromorphic maps on a domain is spherically normal if and only if its spherical derivatives are locally bounded.
\end{theorem}
\begin{proof}
Where $f$ is finite and nonzero the two formulas agree by differentiation. A reciprocal coordinate is holomorphic at a pole, so it also proves continuity there. Direct substitution and differentiation show $(T_a\circ f)^\#=f^\#$ for the isometries above. The spherical line element is $2|dw|/(1+|w|^2)$; differentiating the stereographic formula proves this identity. If $f^\#\le M$ on a disk, the length of the spherical image of a line segment is at most $2M$ times its length. Poles cause no difficulty: on the compact segment there are finitely many, and the reciprocal chart supplies the same estimate across each. Chord distance is bounded by spherical length, so the family is equicontinuous there. Compactness of the target and \cref{d12:11-04,d12:11-06} give spherical normality.
Conversely, suppose $f_n^\#(a_n)\to\infty$ for points in a fixed compact set. Pass to $a_n\to a$ and to a spherically locally uniform subsequence $f_n\to f$. A fixed isometry $T$ taking $f(a)$ to zero makes $T f_n$ uniformly bounded on a fixed disk about $a$, by the chart argument in \cref{d12:11-06}. Cauchy's estimate bounds $(T f_n)'$ uniformly on its half-disk. For large $n$ this bounds $f_n^\#(a_n)=(T f_n)^\#(a_n)\le|(T f_n)'(a_n)|$, a contradiction. The criterion therefore includes the converse left as an exercise in the cited source.
\end{proof}
\source{CON VII.3 pp.157--160 (PDF171--174); forward proof and exercise-converse checked and completed here.}

% END chapters/11

% BEGIN chapters/12
\chapter{Riemann mapping, reflection and boundary mapping}\label{ch:d12-12}

All square roots used in the mapping proof are supplied by the homotopy-Cauchy logarithm theorem \cref{d12:04-06}. No connected-complement characterization from Chapter~\ref{ch:d12-15} is used.

\begin{lemma}[A bounded univalent class]\family{D12.12.02}\label{d12:12-02}
If a proper domain $\Omega$ admits a holomorphic square root of every zero-free holomorphic function, then, for $z_0\in\Omega$, there exists an injective holomorphic $h:\Omega\to\D$ with $h(z_0)=0$ and $h'(z_0)>0$.
\end{lemma}
\begin{proof}
Choose $a\notin\Omega$ and $g^2=z-a$. Then $g$ is injective, nonzero and has nonzero derivative. Its image is disjoint from its negative: $g(z)=-g(w)$ would imply $z=w$ and $g(z)=0$. Put $b=g(z_0)$ and choose $r>0$ with $B(b,r)\subset g(\Omega)$. Thus $B(-b,r)$ misses $g(\Omega)$, and $q=1/(g+b)$ is injective and bounded by $1/r$. The function $h=(r/4)(q-q(z_0))$ has modulus at most $1/2$, vanishes at $z_0$, and has nonzero derivative there. Multiply by a unimodular constant to make that derivative positive real. This proves nonemptiness, not merely boundedness of a hypothetical class.
\end{proof}

\begin{lemma}[The extremal map fills the disk]\family{D12.12.03}\label{d12:12-03}
Under the same square-root hypothesis there is a biholomorphism $F:\Omega\to\D$, $F(z_0)=0$, with positive real derivative.
\end{lemma}
\begin{proof}
Let $\mathcal F$ be all injective holomorphic maps into $\D$ fixing $z_0$. The preceding lemma shows its derivative supremum $s=\sup_{f\in\mathcal F}|f'(z_0)|$ is positive. A closed disk about $z_0$ contained in $\Omega$ and the Cauchy estimate show $s<\infty$. Select $f_n$ whose derivatives approach $s$, and rotate their target values to make those derivatives positive. Montel gives a locally uniform subsequence with holomorphic limit $F$ and $F'(z_0)=s>0$. It is injective by \cref{d12:09-06}. Its values lie in the closed disk, and the maximum principle and nonconstancy put them in the open disk.
Suppose $\alpha\in\D\setminus F(\Omega)$; then $\alpha\ne0$. Put $A(w)=(w-\alpha)/(1-\bar\alpha w)$ and choose $g^2=A\circ F$ on $\Omega$. It is an injective map into $\D$ because equality of its values implies equality of their squares. If $b=g(z_0)$, then $0<|b|<1$, and $H=(g-b)/(1-\bar b g)$ belongs to $\mathcal F$. There is a disk self-map
$\psi=A^{-1}\circ(w\mapsto w^2)\circ((w-b)/(1-\bar b w))^{-1}$
with $F=\psi\circ H$ and $\psi(0)=0$. This map is not injective, since the square map on $\D$ is not injective. Schwarz rigidity gives $|\psi'(0)|<1$. Differentiation gives $s=|F'(z_0)|<|H'(z_0)|\le s$, a contradiction. Thus $F$ is onto; its holomorphic inverse follows from \cref{d12:05-06}.
\end{proof}

\begin{theorem}[Normalized planar Riemann mapping]\family{D12.12.01}\label{d12:12-01}
For a nonempty simply connected proper plane domain and $z_0$ in it, there is a unique biholomorphism $F:\Omega\to\D$ such that $F(z_0)=0$ and $F'(z_0)>0$.
\end{theorem}
\begin{proof}
Homotopy-Cauchy gives a logarithm, and hence a square root, for every zero-free holomorphic function on this domain. The preceding two lemmas give existence. If $F,G$ have the stated normalization, $F\circ G^{-1}$ is a disk automorphism fixing zero, so it is $w\mapsto cw$, $|c|=1$. Positive real derivatives force $c=1$, proving uniqueness. The exclusions are necessary: the empty set has no basepoint, and a holomorphic map from $\C$ into $\D$ is constant by Liouville. No boundary extension is claimed.
\end{proof}
\source{CON VII.4 pp.160--164 (PDF174--178); AHL VI.1 pp.229--232; SS8.3 pp.224--231 (PDF243--250).}

\begin{theorem}[Line and circle reflection]\family{D12.12.04}\label{d12:12-04}
Let a disk centered on a real interval be split into upper and lower halves. If $f$ is holomorphic in the upper half, continuous on it together with the interval, and real on the interval, its extension $F(z)=\overline{f(\bar z)}$ in the lower half is holomorphic on the whole disk. The conjugated assertion holds for any source and target generalized circles, locally away from a chosen chart pole.
\end{theorem}
\begin{proof}
Conjugating the local power series proves holomorphy in the lower half. The real boundary values make the two functions continuous across the interval. To verify Morera, split a triangle by the horizontal lines $\Imm z=\pm\eps$. The integrals on its finitely many pieces vanish in each open half by Cauchy. Internal edges cancel; the two remaining parallel cut integrals cancel in the limit by uniform continuity on the compact triangle, and the short joining edges have length tending to zero and bounded integrand. Thus the full triangle integral is zero. Morera proves the assertion. M\"obius charts taking the two circles to the real line reduce the local circle version to this formula; their inverses are holomorphic on the selected neighborhoods. At a chart pole use the reciprocal sphere chart, obtaining a meromorphic extension if appropriate.
\end{proof}
\source{SS2.5.4 pp.57--60 (PDF76--79); CON IX.1 pp.210--213 (PDF224--227).}

\begin{theorem}[Analytic free arcs]\family{D12.12.05}\label{d12:12-05}
Suppose source and target boundary arcs have parametrizations extending to injective holomorphic maps on neighborhoods of real intervals, with nonzero derivative, and the adjacent domains lie locally on one side. A holomorphic map on the source side, continuous through the arc and mapping the arc into the target arc, extends holomorphically across every interior arc point. If it is a conformal equivalence, is one-to-one on the arc, and carries the given sides to one another, its extension has nonzero derivative there.
\end{theorem}
\begin{proof}
Use the parametrizations and their local holomorphic inverses to make both arcs real. Apply \cref{d12:12-04} and transform back. For the last assertion the straightened extension has a real Taylor expansion at the boundary point and first nonzero term $c(z-a)^m$, $c\ne0$. If $m\ge2$, its imaginary part on a sufficiently small upper semicircle changes sign: choose two angles in $(0,\pi)$ at which $\Imm(c e^{\iu m\theta})$ has opposite nonzero signs; the higher terms cannot change these signs at small radius. This contradicts the prescribed one-sided image. Thus $m=1$. The extension cannot be locally constant, since it was injective on an open side. This proves the nonvanishing derivative without assuming a boundary inverse theorem.
\end{proof}

\begin{theorem}[Conditional Schwarz--Christoffel formula]\family{D12.12.06}\label{d12:12-06}
Let $P$ be a bounded simple polygonal domain with vertices $v_1,\ldots,v_n$ and interior angles $\pi\alpha_j$, $0<\alpha_j<2$. Suppose a biholomorphism $F:\{\Imm z>0\}\to P$ is already given with a homeomorphic extension of the compactified boundaries. If the vertices correspond to finite real points $a_1<\cdots<a_n$ and infinity corresponds to an interior side point, then
\[
 F'(z)=C\prod_{j=1}^n(z-a_j)^{\alpha_j-1},\qquad C\ne0,
\]
where each factor uses its upper-half-plane logarithm. Hence $F(z)=A+C\int_{z_*}^z\prod(\zeta-a_j)^{\alpha_j-1}\dd\zeta$ for $z_* $ in the half-plane. If infinity corresponds to the last vertex, the product has just the finite $n-1$ prevertices. The finite-vertex local behavior is $F(z)-v_j=(z-a_j)^{\alpha_j}u_j(z)$, with a nonvanishing holomorphic factor after selecting the indicated branch.
\end{theorem}
\begin{proof}
Near a finite prevertex, the polygon is exactly a wedge. The given boundary homeomorphism ensures that nearby points map into that wedge. Straighten it by the branch $(F-v_j)^{1/\alpha_j}$ and rotate its two edges onto a straight interval. Reflection extends this function across the real axis. It has a simple zero at $a_j$ by the one-sided Taylor argument in \cref{d12:12-05}; write it as $(z-a_j)b_j(z)$ with $b_j$ nonzero. A local logarithm of $b_j$ proves the asserted angle behavior. Differentiating gives
$F''/F'=(\alpha_j-1)/(z-a_j)+h_j$, with $h_j$ holomorphic near $a_j$.
Across every open side reflection extends $F$ with nonzero derivative. On its real preimage interval $F'$ has a constant complex argument, so $F''/F'$ is real. Its lower-half-plane reflection is therefore the single function $\overline{(F''/F')(\bar z)}$, independent of the target side's reflecting line. These extensions glue globally; the vertex expansions remove their stated simple principal parts. At infinity, use $\zeta=-1/z$ as source half-plane chart. Its boundary image is an open side, so the same reflection and simple-zero argument give $F(z)=v+c/z+O(z^{-2})$, $c\ne0$. Therefore $F''/F'=-2/z+O(z^{-2})$ there. It follows that
$F''/F'-\sum_j(\alpha_j-1)/(z-a_j)$ extends to an entire function tending to zero at infinity. Liouville makes it zero. The logarithmic derivative of $Q(z)=\prod_j(z-a_j)^{\alpha_j-1}$ is precisely that sum. Thus $(F'/Q)'=0$ on the half-plane, and $F'=CQ$ with $C\ne0$. Taking a primitive gives the integral formula. Comparing coefficients at infinity also gives $\sum_j(1-\alpha_j)=2$.
If infinity is a prevertex, precompose $F$ by a real M\"obius automorphism moving it to a finite point and a nonvertex to infinity. Apply the proved formula, then the chain rule. Each factor transforms as a nonzero constant times $(z-a_j)/(cz+d)$; the derivative contributes $(cz+d)^{-2}$. The identity $\sum(\alpha_j-1)=-2$ cancels the denominator powers. The factor whose prevertex has moved to infinity becomes constant, leaving exactly the claimed finite product. No accessory-parameter existence or general boundary extension is used or asserted.
\end{proof}
\source{SS8.4.4 pp.241--245 (PDF260--264); AHL VI.2 pp.233--241. Boundary extension is an explicit premise; general Jordan extension and SC existence are deferred.}

% END chapters/12


% BEGIN chapters/13
\chapter{Harmonic functions, conjugates and Poisson theory}\label{ch:d12-13}

\begin{proposition}[The differential definition]\family{D12.13.01}\label{d12:13-01}
A real harmonic function is initially a $C^2$ function $u$ with $u_{xx}+u_{yy}=0$. Real and imaginary parts of a holomorphic function are harmonic. If $F:V\to U$ is holomorphic and $u$ harmonic on $U$, then $u\circ F$ is harmonic on $V$, including at critical points of $F$.
\end{proposition}
\begin{proof}
Holomorphic functions are $C^\infty$ by Taylor theory. Differentiating the Cauchy--Riemann equations gives $u_{xx}=-v_{yx}$ and $u_{yy}=v_{xy}$, proving the assertion for the real part; interchange the parts for the imaginary part. If $F=p+\iu q$, the real chain rule, $\Delta p=\Delta q=0$, and $p_x=q_y$, $p_y=-q_x$ give
$\Delta(u\circ F)=|F'|^2(\Delta u)\circ F$. The identity also holds when $F'=0$ and proves invariance without an injectivity assumption.
\end{proof}

\begin{theorem}[Conjugates and periods]\family{D12.13.02}\label{d12:13-02}
Every harmonic function has a harmonic conjugate locally. On a domain it has a global single-valued conjugate if and only if the real one-form $-u_y\,dx+u_x\,dy$ has zero integral around every closed piecewise $C^1$ path. The conjugate is unique up to a real constant. In particular it exists on simply connected domains. Harmonic functions are consequently $C^\infty$ and real analytic locally.
\end{theorem}
\begin{proof}
In a rectangle about $(x_0,y_0)$ define
\[
 v(x,y)=-\int_{x_0}^x u_y(t,y_0)\dd t+\int_{y_0}^y u_x(x,s)\dd s.
\]
Then $v_y=u_x$ and $v_x=-u_y(x,y_0)+\int_{y_0}^y u_{xx}(x,s)\dd s=-u_y(x,y)$, using $\Delta u=0$. Thus $u+\iu v$ is $C^1$ and satisfies CR, hence is holomorphic. Globally the form is closed because $\partial_xu_x=\partial_y(-u_y)$. Zero periods make its path integral from a fixed basepoint independent of path; on small rectangles it differs from the preceding primitive by a constant, so it gives a global $v$ with these derivatives. Conversely $dv$ has zero closed-path integrals by FTC. Two such $v$ have zero gradient and are constant on a connected domain. For a fixed-endpoint homotopy subdivide its compact parameter square into finitely many pieces whose images lie in primitive rectangles. The primitive increments cancel on internal edges, so closed-form integrals are homotopy invariant; a contraction therefore gives zero periods. Finally local holomorphic Taylor series prove the claimed regularity. For $u=\log|z|$ on $\C^*$ the conjugate form is $d\arg z$ with circle integral $2\pi$, exhibiting the obstruction.
\end{proof}

\begin{proposition}[Mean values and maximum principles]\family{D12.13.03}\label{d12:13-03}
A harmonic function equals its average on every circle whose closed disk lies in its domain. On a domain an attained maximum or minimum forces constancy. On a bounded open set, $\limsup_{z\to\xi}u(z)\le M$ at every boundary point implies $u\le M$. Conversely a continuous real function with the circle mean-value property on every sufficiently small disk is harmonic.
\end{proposition}
\begin{proof}
On a disk choose a holomorphic function with real part $u$ by \cref{d12:13-02}, and take real parts of its center Cauchy formula. If $u\le u(a)$ everywhere, its average on every small circle about $a$ equals $u(a)$, so continuity forces equality at every circle point and throughout a disk. The maximum set is thus open and closed; connectedness proves constancy. Apply this to $-u$ for a minimum. For the bounded-domain boundary version, a positive excess would give a nonempty compact interior superlevel set: the limsup bounds and finite covering of the compact boundary put it away from the boundary. A maximum on that set is an interior global maximum on its component, contradicting the boundary bound. The same proof works for a continuous function with the mean-value property. To prove the converse, on a small closed disk let $h$ be the Poisson solution with boundary values $u$, constructed in \cref{d12:13-04}. The continuous function $u-h$ has mean values and zero boundary values, so the just-proved maximum and minimum argument gives $u=h$ there. This proves local harmonicity. The Poisson construction below uses only the forward maximum principle, so this reference is acyclic.
\end{proof}

\begin{theorem}[Poisson and the disk Dirichlet problem]\family{D12.13.04}\label{d12:13-04}
For continuous real $\phi$ on $|\zeta|=R$, $R>0$, put
\[
 P\phi(z)=\frac1{2\pi}\int_0^{2\pi}
 \frac{R^2-|z|^2}{|Re^{\iu t}-z|^2}\phi(Re^{\iu t})\dd t,
 \qquad |z|<R.
\]
It is harmonic, extends continuously to the circle as $\phi$, and is the unique continuous closed-disk harmonic solution. Every harmonic $u$ satisfies this formula with its own values on each interior circle.
\end{theorem}
\begin{proof}
The kernel is positive and equals $\Ree((Re^{\iu t}+z)/(Re^{\iu t}-z))$. The geometric expansion of this fraction, uniform on compact subdisks, has constant term one and all other angular averages zero; thus the kernel's normalized integral is one. Integrating that holomorphic fraction against $\phi$ gives a holomorphic function: its derivatives converge uniformly on compact subdisks, or follow from the same uniform geometric expansion. Its real part is $P\phi$, proving harmonicity.
At a boundary point $\xi$, write $P\phi(z)-\phi(\xi)$ as the kernel integral of $\phi(\zeta)-\phi(\xi)$. On $|\zeta-\xi|<\delta$ this difference is at most $\eps$ by uniform continuity, contributing at most $\eps$. For $|z-\xi|<\delta/2$, the remaining contribution is bounded by $2\|\phi\|_\infty(R^2-|z|^2)/(\delta/2)^2$, which tends to zero as $z\to\xi$ from any direction inside the disk. Compactness of the boundary gives a continuous extension. Maximum/minimum principles give uniqueness. Subtract this solution from a harmonic function on a neighborhood of the closed disk to obtain the last assertion. No continuous-mean-value converse was used in this construction.
\end{proof}
\source{FOR22.3--22.5 printed176--178 (PDF185--187); AHL VI.3 pp.241--244.}

\begin{theorem}[Schwarz integral]\family{D12.13.05}\label{d12:13-05}
If $f$ is holomorphic in $|z|<R$ and its real part extends continuously to the closed disk with boundary values $\phi$, then
\[
 f(z)=\frac1{2\pi}\int_0^{2\pi}
 \frac{Re^{\iu t}+z}{Re^{\iu t}-z}\phi(Re^{\iu t})\dd t
 +\iu\Imm f(0).
\]
\end{theorem}
\begin{proof}
The integral is holomorphic and its real part is the Poisson solution, hence equals $\Ree f$ by uniqueness. A holomorphic function with zero real part has zero derivative by CR and is a purely imaginary constant. At zero the integral is real, so that constant is exactly $\iu\Imm f(0)$.
\end{proof}

\begin{theorem}[Harnack inequality and convergence]\family{D12.13.06}\label{d12:13-06}
For $u\ge0$ harmonic on $B(a,R)$ and $|z-a|=r<R$,
\[
 \frac{R-r}{R+r}u(a)\le u(z)\le\frac{R+r}{R-r}u(a).
\]
An increasing sequence of real harmonic functions on a domain either converges locally uniformly to a finite harmonic function or tends to $+\infty$ uniformly on every compact subset.
\end{theorem}
\begin{proof}
Apply the Poisson formula on a radius $s$ with $r<s<R$. Its kernel lies between $(s-r)/(s+r)$ and its reciprocal, and the boundary average is $u(a)$. Let $s\uparrow R$. Chains of overlapping disks propagate these two-sided estimates from one fixed point to any compact set, with a finite positive comparison constant: join each point to the basepoint by a polygonal path, cover its compact trace by a finite chain of interior disks, and then use a finite subcover for the compact set.
Subtract $u_1$ from an increasing sequence. If its values at the basepoint are bounded, they converge, and Harnack applied to the nonnegative differences $u_n-u_m$ bounds their suprema on a compact set by a fixed constant times their basepoint differences. Thus they are locally uniformly Cauchy. Their limit is harmonic: on a smaller disk pass to the limit in the Poisson formula using uniform convergence on an interior boundary circle, whose integral is harmonic by \cref{d12:13-04}. If basepoint values are unbounded, the lower Harnack-chain bound forces uniform compact divergence. The finite limit at one point and divergence at another are therefore impossible.
\end{proof}
\source{FOR22.6 printed178--179 (PDF187--188); AHL VI.3 pp.242--244.}

\begin{theorem}[Harmonic reflection and removable punctures]\family{D12.13.07}\label{d12:13-07}
A harmonic function on a half-disk, continuous and zero on its diameter, extends harmonically by odd reflection. A bounded real harmonic function on a punctured disk extends harmonically across the puncture.
\end{theorem}
\begin{proof}
For reflection take a smaller concentric disk. Give its upper semicircle the original boundary values and its lower semicircle their negative reflected values. These fit continuously at the endpoints because the diameter data are zero. The full-disk Poisson solution is odd by uniqueness, hence vanishes on the diameter. On the upper half-disk it agrees with the original function by the bounded-domain maximum/minimum principle. The smaller disks establish the claimed extension.
For removability, solve the full-disk Poisson problem on a circle of radius $R$ about the puncture, using the given harmonic boundary values. The difference $w$ is bounded by some $M$, harmonic on $0<|z|<R$, and zero on the outer circle. On $\eps<|z|<R$, comparison gives
$|w(z)|\le M\log(R/|z|)/\log(R/\eps)$, since this harmonic comparison is $M$ on the inner circle and zero on the outer one. For fixed $z\ne0$ let $\eps\downarrow0$. The bound tends to zero, so the original function equals the full-disk solution. Boundedness is essential: $\log|z|$ is harmonic on a punctured disk and has no finite extension.
\end{proof}

% END chapters/13

% BEGIN chapters/14
\chapter{Subharmonic comparison, Perron and Green functions}\label{ch:d12-14}

\begin{definition}[Continuous subharmonic functions]\family{D12.14.01}\label{d12:14-01}
A continuous real $u$ on an open set is subharmonic if on each closed interior disk it is bounded above by the harmonic Poisson extension of its boundary values. Equivalently it has the submean property on every interior circle; equivalently it obeys comparison with every harmonic function continuous on such a closed disk.
\end{definition}
\begin{proof}
Poisson comparison implies the center submean inequality. Subtracting a harmonic function preserves the submean inequality. To deduce comparison, on a disk of radius $R$ centered at $a$ with $u-h\le0$ on its boundary, consider $u-h+\eps(|z-a|^2-R^2)$. A positive maximum would be interior. The submean inequality there is strict because the average of $|z-a|^2$ on a circle about $z$ exceeds its center value by the circle radius squared. This contradicts a maximum. Let $\eps\downarrow0$ to get $u\le h$. Applying this to the Poisson extension proves equivalence. The same strict-perturbation argument gives the boundary maximum principle on a bounded domain for subharmonic functions with prescribed limsup bounds. Its strong version follows when a global maximum is attained: equality in the submean inequality forces constancy on each small circle and hence on a neighborhood; connectedness then propagates it.
\end{proof}

\begin{lemma}[Operations and harmonic replacement]\family{D12.14.02}\label{d12:14-02}
Nonnegative linear combinations and finite maxima of continuous subharmonic functions are subharmonic. Replacing $u$ on an interior disk $B$ by its harmonic Poisson extension $H$ gives a continuous subharmonic function $u^B\ge u$. The replacement is monotone in $u$. A local subharmonic function may also be pasted by taking its maximum with another subharmonic function, when its boundary limsup is at most the latter and the resulting function is continuous.
\end{lemma}
\begin{proof}
Linear combinations follow from submeans; maxima follow from comparison, applying it to each function. For replacement $H\ge u$ by definition and equality on the circle ensures continuity. To verify comparison across that circle, let a test disk $D$ meet it and let $h$ be a harmonic majorant of $u^B$ on $\partial D$. Since $u\le u^B$, comparison gives $u\le h$ in $D$. On $B\cap D$, the harmonic function $H-h$ has boundary limsup at most zero: on $\partial B$ this is $u-h$, and on $\partial D$ it is the given bound. The bounded-domain maximum principle therefore gives $H\le h$ there. This proves comparison for $u^B$. Boundary data monotonicity of Poisson gives monotonicity of replacement. For pasting, the same test-disk argument first bounds the unchanged function by $h$, then bounds the pasted function on its portion using its boundary limsup inequality; the maximum is therefore also bounded by $h$. This proves the stated gluing rule, including at an interface.
\end{proof}

\begin{theorem}[Perron's harmonic envelope]\family{D12.14.03}\label{d12:14-03}
Let $\Omega$ be a bounded domain and $\phi:\partial\Omega\to\R$ continuous. Let $\mathcal P_\phi$ consist of continuous subharmonic $u$ on $\Omega$ with $\limsup_{z\to\xi}u(z)\le\phi(\xi)$ at every boundary point. Then
$H_\phi(z)=\sup_{u\in\mathcal P_\phi}u(z)$ is finite and harmonic on $\Omega$.
\end{theorem}
\begin{proof}
Put $m=\min\phi$, $M=\max\phi$. The constant $m$ belongs to the class, and the subharmonic boundary maximum principle gives $u\le M$ for each member. The class is closed under finite maxima and interior harmonic replacement, because replacements change nothing near the boundary.
Fix a closed disk $\overline B\subset\Omega$ and its center $a$. Choose class members whose values at $a$ tend to $H_\phi(a)$. Take successive maxima with $m$, then harmonic replacements on $B$, producing increasing class members $v_n$ harmonic on $B$, bounded between $m$ and $M$, and converging at $a$ to $H_\phi(a)$. Harnack gives a harmonic local uniform limit $v$ on $B$, with $v\le H_\phi$ and $v(a)=H_\phi(a)$. For any fixed $w\in\mathcal P_\phi$, replace $\max(v_n,w)$ on $B$. These replacements increase to a harmonic $V\ge v$, bounded by $H_\phi$, and $V(a)=v(a)$. The strong minimum principle applied to $V-v$ gives $V=v$. Since $w\le V$ on $B$, taking the supremum over $w$ gives $H_\phi\le v$. Thus the envelope is harmonic on $B$. Every point is such a center, proving the theorem. No continuity of a general supremum was assumed.
\end{proof}
\source{FOR22.8--22.13 printed179--182 (PDF188--191); CON X.3--4 pp.263--270 (PDF277--284).}

\begin{theorem}[Barrier squeeze]\family{D12.14.04}\label{d12:14-04}
A positive local barrier at $\xi\in\partial\Omega$ is a continuous superharmonic function $b$ on $\Omega\cap B(\xi,r)$, extending continuously to its relative closure, with $b(\xi)=0$ and $b>0$ at all other points there. Its existence implies $H_\phi(z)\to\phi(\xi)$ for every continuous boundary datum. Such a boundary point is called regular.
\end{theorem}
\begin{proof}
Choose $0<\rho<r$ so that $|\phi(\eta)-\phi(\xi)|<\eps$ on boundary points within $\rho$. On the compact artificial boundary $\overline\Omega\cap\partial B(\xi,\rho)$ the barrier has a positive minimum $\beta$, unless that set is empty, in which case the comparisons below have no artificial-boundary condition. Choose $C$ large enough that $\phi(\xi)+\eps+Cb\ge M$ there. Subharmonic comparison applied to each class member gives $u\le\phi(\xi)+\eps+Cb$ in the cap, hence the same upper bound for $H_\phi$.
Choose $C$ also large enough that $\phi(\xi)-\eps-Cb\le m$ on the artificial boundary. Inside the cap take the maximum of this subharmonic function and $m$, and set it equal to $m$ outside. The gluing lemma makes a class member: on the actual boundary its upper values are at most $\phi$, and on the artificial boundary it fits $m$. It gives $H_\phi\ge\phi(\xi)-\eps-Cb$ near $\xi$. Let $z\to\xi$ and then $\eps\downarrow0$. The squeeze proves attainment.
\end{proof}

\begin{proposition}[Explicit regular points and a failure]\family{D12.14.05}\label{d12:14-05}
Disk boundary points and boundary points with an exterior cone are regular. An exterior tangent disk also gives a barrier. The puncture of a punctured disk is not regular for arbitrary continuous data on that domain's boundary.
\end{proposition}
\begin{proof}
For $B(c,R)$ and $\xi=c+Re^{\iu\theta}$ use the positive harmonic barrier $R-\Ree(e^{-\iu\theta}(z-c))$, zero only at $\xi$ on the closed disk. For an exterior cone, shrink its opening slightly. The remaining allowed wedge has opening $\omega<2\pi$, and directions of $\overline\Omega$ near $\xi$ stay a positive angle away from its two edges. In coordinates $z-\xi=re^{\iu t}$ on that wedge,
$b=r^{\pi/\omega}\sin(\pi t/\omega)$ is positive harmonic, tends to zero at $\xi$, and is strictly positive at every other point of the cap closure. Its positivity follows from $0<t<\omega$, and harmonicity from the imaginary part of the selected holomorphic power. If an exterior disk is tangent at $\xi$, take the disk of half its radius centered halfway from its center to $\xi$. Its circle touches the original disk only at $\xi$. The logarithm of distance to this new center divided by its radius is a strict positive harmonic barrier locally on $\overline\Omega\setminus\{\xi\}$.
For $\D\setminus\{0\}$ assign boundary datum zero on the unit circle and one at the isolated boundary point zero. This datum is continuous on the boundary. A continuous closed-domain harmonic solution would be bounded near zero and thus removable by \cref{d12:13-07}; the outer zero data would make it identically zero, contradicting its value at the puncture. In fact the Perron envelope is zero: if $u$ is in its class, $u+\eps\log|z|$ has limsup at most zero on the outer circle and tends to minus infinity at zero, so its maximum principle gives $u\le-\eps\log|z|$. Let $\eps\downarrow0$; the constant zero is itself in the class.
\end{proof}
\source{FOR22.14--22.18 printed182--184 (PDF191--193); CON X.4 pp.269--275 (PDF283--289).}

\begin{theorem}[The regular Dirichlet problem]\family{D12.14.06}\label{d12:14-06}
On a bounded domain every boundary point of which has a barrier, every continuous real boundary datum has a unique harmonic solution continuous on the closure. The solution is $H_\phi$, lies between the minimum and maximum of its boundary data, and depends linearly and positively on those data.
\end{theorem}
\begin{proof}
The preceding envelope and barrier theorems give harmonicity and every boundary limit. These limits give continuity of the extension, including along boundary sequences because $\phi$ is continuous there. Uniqueness is the maximum/minimum principle for the difference. The sum or real scalar multiple of two solutions is a solution for the corresponding boundary data, so uniqueness proves linearity. Nonnegative boundary values give a nonnegative solution by the minimum principle.
\end{proof}

\begin{lemma}[Circular exhaustions and the Green identity]\label{d12:14:circular-exhaustion}
Every plane domain has an exhaustion $D_j$ by connected finite unions of open disks, with $\overline D_j\subset D_{j+1}\Subset\Omega$, such that circle boundaries have no tangencies or triple intersections. These are regular domains. Their boundaries consist of finitely many analytic arcs and corners of interior angle in $(0,2\pi)$. If a harmonic function on $D_j$ is continuous and zero on a boundary arc through a corner of angle $\theta$, then near that corner
$u=O(r^{\alpha})$, $\nabla u=O(r^{\alpha-1})$, $\alpha=\pi/\theta>0$.
Consequently the second Green identity applies to pairs of their Green functions after removing poles and corners, and its corner boundary terms tend to zero.
\end{lemma}
\begin{proof}
Start with a compact exhaustion of $\Omega$ and a basepoint. At each stage cover that compact and the previous $\overline D_j$ by finitely many closed disks lying in $\Omega$, and join their centers to the basepoint by finitely many polygonal paths in $\Omega$. Their union is a connected compact set $K$. Cover $K$ by finitely many smaller open disks centered on $K$, retaining an open interval of possible larger radii whose closed disks still lie in $\Omega$. Choose these radii successively. For a new circle exclude its finitely many tangency radii with earlier circles and its finitely many radii through intersections of two earlier circles. Distinct centers can be chosen by arbitrarily small perturbation before fixing the covering margins. The union of the resulting disks contains $K$ and is connected because each disk meets $K$. This gives nesting and exhaustion. The finite arrangement of transverse circles has a finite boundary graph; every boundary point has either an analytic arc neighborhood or a sector neighborhood with angle strictly between zero and $2\pi$. Its complement contains an exterior cone, proving regularity by \cref{d12:14-05}.
At an intersection, a M\"obius map taking the other intersection of the two full circles to infinity makes both arcs straight, with nonzero local derivative at the corner in question. The interior is a wedge of angle $\theta$, possibly reentrant. A selected power of exponent $\alpha=\pi/\theta$ takes it locally to a half-plane. The transformed harmonic function is continuous and zero on the diameter, so odd harmonic reflection extends it across the diameter and the corner. It is $O(|w|)$ there with bounded gradient, by its convergent harmonic series. Transforming back proves the two estimates. On a smooth arc, the same reflection without a power gives smooth extension.
Remove small disks about the corners and, when needed, about the poles. Choose their radii sufficiently small that the disks are disjoint, each corner disk meets only its two incident boundary arcs, and all new circle intersections are transverse with no triple intersection. The finitely many forbidden tangency and intersection radii can be avoided. On every retained original arc the harmonic functions extend $C^\infty$ by reflection. The new joins with the truncating circles still require justification before applying a $C^1$-boundary Green theorem; we give that justification explicitly.

Fix these truncation radii first. At a new join the two intersecting circles are straightened by a local M\"obius map taking their other intersection to infinity. The cut domain becomes a wedge $0<\arg w<\theta$, $0<\theta<2\pi$. In a small coordinate neighborhood replace its corner tip by the arc
\[
 \gamma_\eps(t)=\eps\left(1-\tfrac12\sin(\pi t)\right)
       \exp\!\left(\iu\theta(3t^2-2t^3)\right),\qquad 0\le t\le1.
\]
Its radius is between $\eps/2$ and $\eps$, and its angle increases strictly in the interior, so it is injective and stays inside the wedge. Its derivative is nonzero: the radial derivative is nonzero at both endpoints, while the angular derivative is positive at every interior point. Its endpoint tangents are radial and match the two retained rays. Its length is $O(\eps)$. Removing the tip bounded by this arc and the two short radial segments therefore gives a $C^1$ boundary at this join; continuous unit tangents give a regular $C^1$ parametrization across the joins after arclength reparametrization. Pull back these arcs in finitely many disjoint coordinate neighborhoods. This produces a domain with a genuine $C^1$ boundary, without assuming a finite-cell decomposition of the original cut domain.

For fixed truncation radii, $u,v$ extend $C^2$ to a neighborhood of every new join by reflection across its single original analytic boundary arc; the truncation circle itself lies in the original domain on its interior side. Thus the functions and their first derivatives are bounded on all the rounding neighborhoods. The ordinary D04 Green identity now applies on the rounded domain with its actual near-closure hypotheses. The added and removed boundary arcs have length $O(\eps)$, so their flux terms tend to zero as $\eps\downarrow0$; any bounded area terms have area $O(\eps^2)$. This first limit proves the identity on the fixed cut domain. Finally, for two zero-boundary functions the integrand $u\partial_\nu v-v\partial_\nu u$ on an original corner circle is $O(r^{2\alpha-1})$ and its length is $O(r)$, so its integral tends to zero when that corner radius tends to zero. The rounding limit was taken first. This does not assume differentiability of a Green function on an untrimmed reentrant corner.
\end{proof}
\source{Prior D04 notes: d04:12:parts (second Green identity, $C^2$ near the closure), d04:12:green and d04:15:piecewise. The exhaustion and corner estimates above supply the additional hypotheses; no boundary PDE regularity theorem is assumed.}

\begin{theorem}[Green functions, symmetry and invariance]\family{D12.14.07}\label{d12:14-07}
Every bounded domain has a canonical positive Green function $g_\Omega(z,a)$: it is harmonic for $z\ne a$, and $g_\Omega(z,a)+\log|z-a|$ is harmonic at $a$. It is the least positive function with these properties. It is symmetric in $z,a$ and invariant under biholomorphisms between bounded domains. At every regular boundary point it tends to zero. On a bounded regular domain it is the unique Green function continuous with value zero at all boundary points, and is obtained by solving the boundary correction $\log|\xi-a|$.
\end{theorem}
\begin{proof}
On a bounded regular domain $D$ solve that correction by \cref{d12:14-06}: if $h$ has boundary data $\log|\xi-a|$, then $g=h-\log|z-a|$ has the required pole and zero boundary. It is positive by the minimum principle on $D\setminus\overline{B(a,\eps)}$, where the inner boundary tends uniformly to $+\infty$, and the strong principle. Any two such zero-boundary functions differ harmonically across the pole and vanish on the boundary, so are equal.
Use the circular exhaustion containing $a$. Its Green functions $g_j$ increase: the difference $g_{j+1}-g_j$ is harmonic across $a$, and its boundary values on $\partial D_j$ are nonnegative. Enclose $\Omega$ in a larger disk. Its explicit zero-boundary Green function, obtained from the disk automorphism formula $-\log|T_a(z)|$, bounds every $g_j$ above by the same comparison. Harnack therefore gives a finite harmonic limit off $a$. A small disk centered at $a$ contained in $D_1$ bounds this limit below by its own Green function. Subtracting $-\log|z-a|$, the upper and lower disk bounds give a bounded harmonic correction near $a$; \cref{d12:13-07} removes its puncture.
If $v$ is any positive Green function with this normalized pole on $\Omega$, $v-g_j$ is harmonic across $a$ and nonnegative on $\partial D_j$, so $v\ge g_j$. The limit is therefore least and is independent of exhaustion. For a regular outer domain its Dirichlet-corrected Green function bounds the limit above, so the limit tends to zero on the boundary and equals that correction by uniqueness. For an individual regular point $\xi$, choose a barrier cap avoiding $a$. The enclosing-disk Green function bounds all $g_j$ by one constant $M$ on the artificial cap circle. The barrier has positive minimum $\beta$ on that circle. Choose $C\beta\ge M$. On the artificial circle $g_j\le Cb$, and on every portion of $\partial D_j$ inside the cap, $g_j=0\le Cb$. The maximum principle on $D_j$ intersected with the cap gives $g_j\le Cb$ before any limit is taken. Pass to the increasing limit, then let $z\to\xi$; $0\le g_\Omega(z,a)\le Cb(z)\to0$. This imposes no unproved boundary data on the limiting function.
For symmetry first work on $D_j$ with distinct poles $a,b$. Apply the second Green identity to $u=g_j(\cdot,a)$ and $v=g_j(\cdot,b)$ after removing small disks about the poles and corners. The Laplacians are zero there. The original boundary arcs contribute zero; corner contributions vanish by \cref{d12:14:circular-exhaustion}. The inner boundary about $a$, whose outward normal for the punctured region points toward $a$, contributes $-2\pi v(a)$ in the limit, since $u=-\log r+O(1)$ and $\partial_\nu u=1/r+O(1)$. The boundary about $b$ contributes $2\pi u(b)$. The terms involving $r\log r$ tend to zero. Hence $g_j(a,b)=g_j(b,a)$. Pass to the increasing exhaustion limit for both poles to obtain symmetry on $\Omega$.
Finally if $F:\Omega\to\Omega'$ is a biholomorphism of bounded domains, $g_{\Omega'}(F(z),F(a))$ is a positive Green function at $a$. Indeed $F(z)-F(a)=(z-a)q(z)$ with $q(a)\ne0$, and $\log|q|$ is harmonic locally. Minimality gives one inequality; apply the same argument to $F^{-1}$ for the reverse inequality. No boundary extension of $F$ is used. At irregular points the canonical Green function has not been asserted to vanish; in the punctured disk that assertion would be false.
\end{proof}
\source{CON X.5 pp.275--278 (PDF289--292) proves bounded regular existence and invariance; its Green identity exercise is expanded here with an actual exhaustion and corner control.}

\begin{corollary}[Harmonic measures of boundary components]\family{D12.14.08}\label{d12:14-08}
Suppose a bounded regular domain has finitely many disjoint compact boundary components $C_1,\ldots,C_m$. There are unique continuous closed-domain harmonic functions $\omega_j$ equal to one on $C_j$ and zero on the others. They satisfy $0\le\omega_j\le1$, $\sum_j\omega_j=1$, and the solution with constant component data $c_j$ is $\sum_jc_j\omega_j$. If $m>1$, each $\omega_j$ lies strictly between zero and one in the domain.
\end{corollary}
\begin{proof}
Finitely many disjoint compact components have positive pairwise separation. Their indicator boundary functions are therefore continuous. Apply \cref{d12:14-06}; positivity and the maximum bound give the inequalities. The sum has boundary value one, so uniqueness makes it constant one. Linearity gives the data formula. For $m>1$ a function attaining zero or one inside would be constant by the strong principle, contradicting its distinct boundary values. No canonical slit-map existence is inferred from these functions.
\end{proof}

% END chapters/14

% BEGIN chapters/15
\chapter{Runge approximation and planar simple connectivity}\label{ch:d12-15}

\begin{lemma}[A compatible compact exhaustion]\family{D12.15.01}\label{d12:15-01}
Every open $\Omega\subset\C$ has compact sets $K_n\subset\operatorname{int}K_{n+1}$ exhausting its compact subsets, such that every component of $\Chat\setminus K_n$ meets $\Chat\setminus\Omega$. They may be chosen as finite unions of closed grid cells with selected bounded holes filled. Empty open sets and empty compact stages are allowed.
\end{lemma}
\begin{proof}
At stage $n$ cover the standard compact set $\{|z|\le n,\dist(z,\C\setminus\Omega)\ge1/n\}$ and the preceding compact $K_{n-1}$ by finitely many sufficiently small closed grid squares, including one neighboring layer, all contained in $\Omega$. Their union $L_n$ contains those sets in its interior. Its complement has finitely many bounded components: in a large box containing all its cells there are only finitely many remaining grid cells, and outside that box everything is connected. Fill precisely those bounded complementary components which do not meet $\C\setminus\Omega$, and call the resulting compact set $K_n$. Each filled component has boundary in $L_n$ and closure in $\Omega$, so the finite filling remains compact in $\Omega$. A remaining bounded component meets $\C\setminus\Omega$; the unbounded component meets the spherical complement at infinity. Nesting and compact exhaustion follow from the construction. The inclusion of a neighboring cell layer supplies the required interior containment, so no boundary-distance or hole-filling limit has been suppressed.
\end{proof}

\begin{lemma}[A finite polygon representation]\family{D12.15.02}\label{d12:15-02}
If $K$ is compact, $K\subset U$ open, and $f\in\Hol(U)$, there is a finite oriented polygonal cycle $\Gamma$ in $U\setminus K$ such that
\[
 f(z)=\frac1{2\pi\iu}\int_\Gamma\frac{f(\zeta)}{\zeta-z}\dd\zeta\qquad(z\in K).
\]
\end{lemma}
\begin{proof}
Take a grid mesh so small that all cells meeting $K$ and their neighboring layer lie in $U$. Sum their positive cell boundaries and cancel oppositely directed common edges. Every remaining edge misses $K$, with a positive distance by the neighboring layer; all lie in $U$. At each vertex the remaining directed edges have equal incoming and outgoing counts, so this finite chain decomposes into closed polygonal paths. Its index is one at interior points of the cell union not on grid lines and zero outside the union, by the explicit rectangle index formula. Add the cell Cauchy formulas at any off-grid point in the interior of this union, whether or not that point belongs to $K$. Their integrals cancel on common edges and give the displayed identity there. Both sides are continuous on the interior of the union, which contains $K$; every grid-line point there is a limit of off-grid interior points. The identity therefore holds throughout $K$, even if $K$ itself consists entirely of grid-line points. If $K$ is empty take the zero cycle.
\end{proof}
\source{CON VIII.1 pp.195--197 (PDF209--211); SS2.5.5 Lemma5.8 pp.61--62 (PDF80--81).}

\begin{lemma}[Rational Riemann sums]\family{D12.15.03}\label{d12:15-03}
The representation in \cref{d12:15-02} can be approximated uniformly on $K$ by rational functions whose poles lie on $\Gamma$, hence outside $K$.
\end{lemma}
\begin{proof}
For a parametrized segment $\gamma:[0,1]\to U\setminus K$ its integrand is $F(z,t)=f(\gamma(t))\gamma'(t)/(\gamma(t)-z)$. The denominator has a positive minimum on $K\times[0,1]$. Thus $F$ is uniformly continuous there. Riemann sums converge uniformly for $z\in K$: the error on each subinterval is bounded by its length times the common modulus of continuity, and the total length of the parameter interval is one. Every sum is a finite linear combination of $1/(\gamma(t_j)-z)$. Sum the approximations for the finitely many segments, retaining their orientations and coefficients.
\end{proof}

\begin{theorem}[Runge with specified allowed poles]\family{D12.15.04}\label{d12:15-04}
Let $K\subset\C$ be compact and let $A\subset\Chat\setminus K$ meet every component of that spherical complement. Every function holomorphic on a neighborhood of $K$ is uniformly approximable there by rational functions with all poles in $A$. A pole at infinity means a polynomial part; it is permitted only when infinity is in $A$.
\end{theorem}
\begin{proof}
Let $\mathcal R_A$ be the uniform closure on $K$ of the permitted rational functions. It is a closed algebra containing constants; addition and multiplication preserve the pole restriction, and products of uniformly converging bounded sequences converge uniformly. For finite $a\notin K$ write $r_a(z)=1/(z-a)$ and let $S=\{a:r_a\in\mathcal R_A\}$. It is closed relative to $\C\setminus K$, because resolvents converge uniformly when their poles approach a point still off $K$. It is open there: if $a\in S$ and $|b-a|<\dist(a,K)$, then
\[
 r_b=\frac{r_a}{1-(b-a)r_a}
 =\sum_{k\ge0}(b-a)^k r_a^{k+1}
\]
uniformly on $K$, putting $b\in S$. Every finite point of $A$ is in $S$. If infinity belongs to $A$, every sufficiently distant finite $b$ is in $S$ by the polynomial geometric expansion $r_b=-\sum_{k\ge0}z^k/b^{k+1}$. Since components of an open plane set are polygonally connected, open-and-closed propagation proves that $S$ contains every finite complementary component represented by a finite point of $A$, and also the unbounded component when it is represented by infinity. The unbounded component with infinity removed is connected: it is the unique plane component containing the exterior of a large disk. Thus all finite complementary points belong to $S$. The rational sums of \cref{d12:15-03} therefore lie in $\mathcal R_A$, and their limit is the desired $f$. If $K$ is empty the zero rational function suffices. This proves pole movement with an actual convergent algebra argument, not a deformation through poles on $K$.
\end{proof}
\source{CON VIII.1 pp.195--200 (PDF209--214); SS2.5.5 pp.60--64 (PDF79--83).}

\begin{corollary}[Approximation on an open set]\family{D12.15.05}\label{d12:15-05}
If $A\subset\Chat\setminus\Omega$ meets each component of that complement, every $f\in\Hol(\Omega)$ is a locally uniform limit of rational functions with poles in $A$. If $\Chat\setminus\Omega$ is connected, polynomials suffice. No connectedness of $\Omega$ is needed.
\end{corollary}
\begin{proof}
Take the exhaustion from \cref{d12:15-01}. Each component $V$ of $\Chat\setminus K_n$ contains a point of $\Chat\setminus\Omega$, and the entire connected component of that point in $\Chat\setminus\Omega$ lies in $V$. Thus $V$ contains a point of $A$. Runge approximates $f$ on $K_n$ with error $2^{-n}$. Compact exhaustion gives local uniform convergence. If the spherical complement of $\Omega$ is connected it contains infinity, so take $A=\{\infty\}$. The rational functions with only that possible pole are polynomials. In contrast, $1/z$ on an annulus is not approximable by polynomials on a circle surrounding zero: its integral there is $2\pi\iu$, whereas every polynomial integral is zero and uniform convergence preserves that integral.
\end{proof}

\begin{theorem}[The planar polynomial hull]\family{D12.15.06}\label{d12:15-06}
The polynomial hull $\widehat K=\{a:|p(a)|\le\sup_K|p|\text{ for every polynomial }p\}$ of a nonempty compact plane set is $K$ together with all bounded components of $\C\setminus K$.
\end{theorem}
\begin{proof}
On a bounded complementary component a polynomial has modulus at most its boundary supremum by the maximum principle; that boundary lies in $K$. Thus every filled hole lies in the hull. If $a$ lies in the unbounded component, the resolvent argument in the proof of \cref{d12:15-04}, using just $A=\{\infty\}$ on this one component, approximates $1/(z-a)$ on $K$ by a polynomial $p$. Choose the approximation so that $q(z)=1-(z-a)p(z)$ satisfies $\sup_K|q|<1/2$. But $q(a)=1$, separating $a$ from the hull. These two observations prove equality. The empty compact set has empty hull under the convention that constants already separate every point from the empty set.
\end{proof}

\begin{theorem}[Planar equivalences, with no circular prerequisite]\family{D12.15.07}\label{d12:15-07}
For a plane domain the following are equivalent: simple connectivity; connected spherical complement; zero closed-path periods for every holomorphic function; existence of primitives for every holomorphic function; existence of logarithms for every zero-free holomorphic function; existence of square roots for every such function; and existence of all positive-integer roots for every such function.
\end{theorem}
\begin{proof}
Simple connectivity implies zero periods by homotopy-Cauchy. Zero periods and primitives are equivalent by \cref{d12:03-04}. A primitive $H$ of $f'/f$ makes $fe^{-H}$ a nonzero constant, since its derivative vanishes. Choosing a logarithm of that one constant gives a global logarithm of $f$. Exponentiating fractions of it gives every integer root, hence square roots. Conversely, on a proper domain the universal square-root property is exactly the hypothesis of \cref{d12:12-02,d12:12-03}. Their extremal proof gives a biholomorphism onto the disk without any prior simple-connectivity equivalence. A loop contracts by transporting the straight-line disk contraction through that map. On $\Omega=\C$, straight-line contractions prove simple connectivity directly.
If the spherical complement is connected, \cref{d12:15-05} approximates every holomorphic function locally uniformly by polynomials, whose periods vanish. Thus all periods vanish. For the reverse implication suppose the compact spherical complement is disconnected. Partition it into disjoint nonempty compact sets $A,B$, with infinity in $B$ and thus $A$ bounded. Choose a finite union of small grid cells containing $A$ in its interior and disjoint from $B$, with one neighboring layer and still missing $B$. After internal-edge cancellation its boundary is a finite polygonal cycle in $\Omega$: it meets neither $A$ nor $B$. At a chosen $a\in A$ its index is one, by the same grid index calculation used in \cref{d12:15-02}. Each closed polygonal path in this cycle would have zero integral of the holomorphic function $1/(z-a)$ on $\Omega$. Their sum would therefore have zero integral, contradicting $2\pi\iu$. Hence zero periods force connected complement. This completes all implications. RMT was proved earlier from the homotopy hypothesis, and only its stronger square-root proof is reused here; no equivalence in this theorem was used to establish it.
\end{proof}
\source{CON VIII.2 pp.200--204 (PDF214--218); FB IV AppendixA pp.233--238 (PDF239--244).}

% END chapters/15

% BEGIN chapters/16
\chapter{Mittag--Leffler and Weierstrass constructions}\label{ch:d12-16}

\begin{theorem}[Infinite products and logarithms]\family{D12.16.01}\label{d12:16-01}
For nonzero complex factors $u_n$, convergence of their product to a nonzero value is equivalent to $u_n\to1$ and convergence of the series of principal logarithms of a sufficiently late tail. In particular $\sum|u_n-1|<\infty$ ensures nonzero product convergence. If $u_n$ are holomorphic on an open set and $\sum\sup_K|u_n-1|<\infty$ on every compact subset, their products converge locally uniformly to a holomorphic function. On a component where no factor is identically zero, the local zero order of the product is the sum of the zero orders of the factors; only finitely many terms of this sum are nonzero. If a factor is identically zero on a component, the product is identically zero there instead. A tail has no zeros on any fixed compact neighborhood once its factors are sufficiently close to one.
\end{theorem}
\begin{proof}
For $|w|\le1/2$, the power series for $\log(1+w)$ gives $|\log(1+w)|\le2|w|$. Thus absolute summability, or its uniform compact version, gives a convergent logarithm tail, whose exponential is nonzero; finitely many earlier factors are retained separately. Conversely if partial products $P_n$ converge to $P\ne0$, their ratios $u_n=P_n/P_{n-1}$ tend to one. A fixed disk about $P$ missing zero contains all late $P_n$ and admits one logarithm $L$. The difference $L(P_n)-L(P_{n-1})$ tends to zero and has exponential $u_n$. The principal $\log u_n$ also tends to zero, and the difference between these two logarithms is an integral multiple of $2\pi\iu$, hence is zero for all sufficiently large $n$. The series therefore telescopes and converges. The reverse implication follows by exponentiation. For holomorphic factors the tail logarithms converge uniformly on compact neighborhoods, so the exponential tail is holomorphic and nonzero by \cref{d12:11-02}; multiplication by the finite initial product proves both holomorphy and the exact zero-order assertion whenever the initial factors are not identically zero on the component. An identically zero factor makes every subsequent partial product zero there. Thus isolated zeros and identically zero factors are both treated without taking their logarithms.
\end{proof}

\begin{lemma}[Canonical factors]\family{D12.16.02}\label{d12:16-02}
Put $E_0(w)=1-w$ and $E_p(w)=(1-w)\exp(w+w^2/2+\cdots+w^p/p)$ for $p\ge1$. It has precisely one zero, a simple zero at one. For $|w|\le1/2$,
\[
 |\log E_p(w)|\le2|w|^{p+1},\qquad
 |E_p(w)-1|\le6|w|^{p+1},
\]
where the logarithm is the branch zero at zero.
\end{lemma}
\begin{proof}
The exponential never vanishes. In $|w|<1$ the logarithm is $-\sum_{k\ge p+1}w^k/k$, whose absolute value is at most $|w|^{p+1}/((p+1)(1-|w|))\le2|w|^{p+1}$. This is at most one on the half-disk. The elementary bound $|e^v-1|\le |v|e^{|v|}$, obtained by integrating $e^{tv}v$ for $0\le t\le1$, then gives the safe constant $2e<6$. The zero order follows from the linear factor at one.
\end{proof}

\begin{theorem}[Entire prescribed zeros]\family{D12.16.03}\label{d12:16-03}
Every locally finite effective divisor in $\C$ is the exact zero divisor of an entire function. If its nonzero points, repeated by positive multiplicity, are $a_n$ with $|a_n|\to\infty$, one may use
$z^m\prod_n E_n(z/a_n)$, where $m$ is the prescribed multiplicity at zero. Empty and finite divisors use the corresponding finite products.
\end{theorem}
\begin{proof}
A locally finite divisor has only finitely many points in each closed disk; it is countable and its repeated sequence can be ordered to tend to infinity. On $|z|\le R$, eventually $|a_n|>2R$, so the factor error is at most $6(1/2)^{n+1}$. The product theorem gives local uniform convergence and holomorphy. Near a designated point only finitely many initial factors can vanish; the tail is nonzero. Each occurrence gives one simple zero, so its total order is exactly the prescribed multiplicity, with no other zeros. If the divisor is empty the function one works.
\end{proof}

\begin{theorem}[Entire factorization]\family{D12.16.04}\label{d12:16-04}
Every entire $f\not\equiv0$ has a representation
$f(z)=z^m e^{g(z)}\prod_n E_{p_n}(z/a_n)$,
where $g$ is entire, the $a_n$ list its nonzero zeros with multiplicity, and suitable nonnegative integers $p_n$ can be chosen, for example $p_n=n$. This is not the stronger finite-order degree bound of Chapter~\ref{ch:d12-17}.
\end{theorem}
\begin{proof}
Construct the product $P$ with exactly those zeros by \cref{d12:16-03}. The quotient $f/P$, defined off their zero set, extends holomorphically and without zeros across every point by matching local orders. On $\C$ it has an entire logarithm by homotopy-Cauchy, since every plane loop contracts linearly. Calling that logarithm $g$ proves the formula. The zero function is excluded because it has no locally finite zero divisor; nonzero constants are obtained by a constant logarithm.
\end{proof}
\source{CON VII.5 pp.169--170 (PDF183--184); SS5.3--5.4 pp.140--147 (PDF159--166).}

\begin{theorem}[Prescribed zeros on a general open set]\family{D12.16.05}\label{d12:16-05}
Every locally finite effective divisor on an open plane set is the exact zero divisor of a holomorphic function there, chosen nonzero on every component outside its designated zeros.
\end{theorem}
\begin{proof}
Work first on one domain with infinitely many designated zeros. Choose a point $a$ and a disk $B(a,r)$ in it missing the divisor, possible by local finiteness. Under $T(z)=1/(z-a)$ the punctured domain becomes an open $G$ containing $\{|w|>1/r\}$; its transformed zeros are bounded. Repeat them according to multiplicity as $z_n$. They have no accumulation in $G$. The compact complement $G^c$ is nonempty in this infinite case; choose nearest $b_n\in G^c$ with $|z_n-b_n|=\dist(z_n,G^c)$. These distances tend to zero: otherwise a bounded subsequence would accumulate at a point of $G$, contradicting local finiteness. Both sequences lie in one fixed disk.
Consider
\[
 F(w)=\prod_{n\ge1}E_n\left(\frac{z_n-b_n}{w-b_n}\right),\qquad w\in G.
\]
Each factor is holomorphic there and has a single simple zero at $z_n$. On a compact $K\subset G$ its argument is at most $|z_n-b_n|/\dist(K,G^c)$, hence eventually at most $1/2$. The error estimate in \cref{d12:16-02} gives a summable geometric bound. Thus the product is holomorphic, with exactly the repeated zeros. Moreover $F(w)\to1$ as $w\to\infty$. To verify this last assertion uniformly for the whole product, if $|z_n|,|b_n|\le R$, every argument is at most $2R/(|w|-R)=\delta$. For large $|w|$, $\delta<1/2$ and the logarithm series is bounded by $2\sum_{n\ge1}\delta^{n+1}=2\delta^2/(1-\delta)$, which tends to zero. Hence its exponential tends to one. Consequently $F(T(z))$ extends over $a$ with value one by the removable singularity theorem and has precisely the original divisor. Finite divisors use a finite polynomial product in the original coordinate. Finally define these functions independently on each component of an arbitrary open set; components are open, so the assembled function is holomorphic. No primitive of a Mittag--Leffler solution, with uncontrolled periods, occurs in this proof.
\end{proof}
\source{CON VII.5 Theorem5.15 pp.170--173 (PDF184--187), nearest-complement factors and infinity normalization checked explicitly.}

\begin{corollary}[Holomorphic quotients]\family{D12.16.06}\label{d12:16-06}
Every scalar meromorphic function $f$ on an open plane set is $g/h$ with $g,h$ holomorphic and $h$ nonzero on each component away from its locally finite zero set.
\end{corollary}
\begin{proof}
Use \cref{d12:16-05} to choose $h$ whose zeros are exactly the poles of $f$ with their orders. The product $hf$ extends holomorphically across each pole by local Laurent factorization; set $g=hf$ with those extensions. The quotient agrees as a meromorphic function, including at canceled points by its local interpretation. The identically infinity map is excluded from the scalar meromorphic convention.
\end{proof}

\begin{theorem}[Mittag--Leffler on an open plane set]\family{D12.16.07}\label{d12:16-07}
Let $A$ be a locally finite subset of an open $\Omega$, and prescribe at each $a\in A$ a finite principal part $P_a(z)=\sum_{k=1}^{m_a}c_{ak}(z-a)^{-k}$. There is a meromorphic $f$ with these principal parts and no other poles. A designated point is an actual pole exactly when its prescribed principal part is nonzero. All solutions differ by a holomorphic function.
\end{theorem}
\begin{proof}
Choose the compatible exhaustion $K_n$ of \cref{d12:15-01}, with $K_0=\varnothing$. The finite sets $A_n=A\cap(K_n\setminus K_{n-1})$ partition $A$. Put $Q_n=\sum_{a\in A_n}P_a$. It is holomorphic on a neighborhood of $K_{n-1}$, since its finitely many poles lie outside that closed set. The complementary-component property and Runge give a rational $R_n$, holomorphic throughout $\Omega$, with $\sup_{K_{n-1}}|Q_n-R_n|<2^{-n}$; use $R_1=0$ when the compact set is empty. Then $f=\sum_n(Q_n-R_n)$ converges locally uniformly off $A$. Indeed on $K_m$, every term with $n>m+1$ has the displayed geometric bound, and only finitely many earlier terms remain.
Around $a\in A$ choose a closed isolating disk contained in the interior of some $K_m$. Subtract $P_a$ from the finite initial terms. They are then holomorphic near that point. The tail terms have no poles on that disk and converge uniformly there by the same estimate, so their limit is holomorphic. Thus $f-P_a$ is holomorphic near $a$, proving the exact principal part and all pole claims. Two solutions have removable differences at all designated points and no other poles, so their difference is holomorphic. Conversely adding any holomorphic function preserves the principal parts.
\end{proof}
\source{CON VIII.3 pp.204--208 (PDF218--222); FB IV.3 pp.219--222 (PDF225--228). The finite-polynomial principal parts here give meromorphic functions; FB's more general entire encodings may prescribe essential singularities and are not silently conflated with them.}

\begin{theorem}[Cotangent and sine]\family{D12.16.08}\label{d12:16-08}
Off the integer poles,
\[
 \pi\cot(\pi z)=\frac1z+\sum_{n\ge1}\frac{2z}{z^2-n^2},
 \qquad
 \sin(\pi z)=\pi z\prod_{n\ge1}\left(1-\frac{z^2}{n^2}\right).
\]
The paired series and product converge locally uniformly on their appropriate domains. Their pairing/order is part of the stated formulas.
\end{theorem}
\begin{proof}
Fix $z\notin\Z$, $z\ne0$, and integrate $\pi\cot(\pi w)/(w^2-z^2)$ around the positive square with $|\Ree w|,|\Imm w|\le N+1/2$. On its vertical sides $|\cot(\pi w)|\le1$ by the sine/cosine exponential formulas. On its horizontal sides these same formulas bound it by $(1+e^{-2\pi(N+1/2)})/(1-e^{-2\pi(N+1/2)})$. The denominator has modulus at least $(|w|-|z|)^2$, and perimeter is $O(N)$, so the integral tends to zero. Residues at integers $n$ are $1/(n^2-z^2)$; the residues at $z$ and $-z$ sum to $\pi\cot(\pi z)/z$. Therefore
$0=\lim_N\sum_{n=-N}^N1/(n^2-z^2)+\pi\cot(\pi z)/z$, which rearranges to the first formula. Its paired terms are $O_K(n^{-2})$ on every compact set off $\Z$, proving local uniform convergence.
The product $Q(z)=z\prod_{n\ge1}(1-z^2/n^2)$ converges locally uniformly by \cref{d12:16-01}, has exactly the simple integer zeros, and has $Q(z)/z\to1$ at zero. Off those zeros, termwise logarithmic differentiation, justified by the uniformly convergent logarithm tails and \cref{d12:11-02}, gives $Q'/Q=1/z+\sum2z/(z^2-n^2)=\pi\cot(\pi z)$. The entire quotient $\sin(\pi z)/Q(z)$ is zero-free, because both numerator and denominator have the same simple zeros. Its derivative is zero off the integers and hence everywhere by continuity. It is constant, and its value at zero is $\pi$. This proves the product without invoking finite-order Hadamard factorization or a future chapter.
\end{proof}
\source{FB III.7 cotangent partial fractions and IV.3 pp.220--223 (PDF226--229); SS5.3.2 pp.142--145 (PDF161--164).}

% END chapters/16


% BEGIN chapters/17
\chapter{Jensen theory and finite-order factorization}\label{ch:d12-17}

The identically zero function is excluded from zero-divisor and order
statements. Zeros and poles below are repeated according to multiplicity.

\begin{theorem}[Poisson--Jensen and Jensen]\label{d12:17-01}\family{D12.17.01}
Let $f$ be meromorphic on a neighborhood of $\{|z|\le R\}$, with no
zero or pole on its boundary, and let $a_j,b_j$ be its zeros and poles
inside. For $|z|<R$ away from those points,
\[
 \log|f(z)|=\frac1{2\pi}\int_0^{2\pi}
 \Ree\frac{Re^{\iu t}+z}{Re^{\iu t}-z}\log|f(Re^{\iu t})|\,dt
 -\sum_j G_R(z,a_j)+\sum_jG_R(z,b_j),
\]
where $G_R(z,a)=\log\left|\dfrac{R^2-\bar a z}{R(z-a)}\right|$.
If $f(0)$ is finite and nonzero, then
\[
 \frac1{2\pi}\int_0^{2\pi}\log|f(Re^{\iu t})|\,dt
 =\log|f(0)|+\sum_j\log\frac R{|a_j|}-\sum_j\log\frac R{|b_j|}.
\]
\end{theorem}
\begin{proof}
The divisor in the closed disk is finite, by isolated zeros and poles
and compactness. Put $B_a(z)=R(z-a)/(R^2-\bar a z)$, including
$B_0(z)=z/R$. It has one simple zero at $a$, no pole in the disk and
modulus one on its boundary. The function
$F=f\prod B_{a_j}^{-1}\prod B_{b_j}$ is holomorphic and nonvanishing
on the disk, with all apparent divisor singularities removed. Thus
$\log|F|$ is harmonic: a local logarithm exists by
\cref{d12:02-05}, and its real part is $\log|F|$.
Apply the Poisson formula \cref{d12:13-04}. On the boundary
$\log|F|=\log|f|$; in the interior substitute
$\log|B_a|=-G_R(\cdot,a)$ and rearrange. Evaluating at zero proves
Jensen's formula because $G_R(0,a)=\log(R/|a|)$.
\end{proof}

\begin{proposition}[The center and boundary zeros]\label{d12:17-02}\family{D12.17.02}
If $f(z)=z^m g(z)$ near zero, where $m\in\Z$ and $g(0)$ is finite and
nonzero, Jensen's formula instead has center term
$\log|g(0)|+m\log R$ and counts only nonzero interior zeros and poles.
For a function meromorphic on a neighborhood of the closed disk,
$\log|f(Re^{\iu t})|$ is integrable even when boundary zeros or poles
occur. Jensen's center formula remains valid, with boundary divisor
terms contributing zero. The displayed Poisson--Jensen theorem above
is used only on divisor-free boundary circles.
\end{proposition}
\begin{proof}
Apply the preceding theorem to $g=f/z^m$, then use
$\log|f(Re^{\iu t})|=m\log R+\log|g(Re^{\iu t})|$.
For the boundary assertion factor all the finitely many zeros and poles
in a neighborhood of the closed disk. The remaining factor is
nonvanishing there. Each elementary factor contributes an integral of
$\log|Re^{\iu t}-a|$, with a signed integer coefficient. This logarithm
is continuous except, when $|a|=R$, for a logarithmic singularity bounded
in absolute value by $C+|\log|t-t_0||$ near its one singular angle.
It is therefore integrable. Its mean is $\log\max(R,|a|)$: for
$|a|\ne R$ expand $\log(1-w)$ with $|w|<1$, whose nonconstant Fourier
terms have zero mean. At $|a|=R$ take the limit from $|a|>R$.
For radii ratios tending to one the negative logarithms are dominated
by $C+|\log|t-t_0||$, since
$|1-re^{\iu t}|^2=(1-r)^2+2r(1-\cos t)\ge c t^2$ for
$1/2\le r\le1$ and $|t|\le\pi$; the positive parts are bounded.
Dominated convergence gives the same mean at equality. Summing the
finite factorizations gives Jensen's formula, with $\log(R/|a|)=0$
at boundary points. No formal value of $\log f(0)$ at a zero is taken.
\end{proof}

\begin{corollary}[Counting zeros]\label{d12:17-03}\family{D12.17.03}
If $f$ is entire and $f(0)=1$, let $n(r)$ count its zeros in $|z|\le r$.
Then $n(r)\log2\le\log M(2r)$, where
$M(r)=\max_{|z|=r}|f(z)|$. For a general nonzero entire function first
factor its center zero and normalize the nonzero remaining value.
\end{corollary}
\begin{proof}
Each zero in $|z|\le r$ contributes at least $\log2$ to Jensen's sum
at radius $2r$; all the other contributions are nonnegative. The
boundary mean is at most $\log M(2r)$. The preceding boundary result
also covers a zero on either of the named circles.
\end{proof}

\begin{definition}[Order, exponent, rank and genus]\label{d12:17-04}\family{D12.17.04}
For a nonzero entire function set
\[
 \rho=\limsup_{r\to\infty}
       \frac{\log\log\max(e,M(r))}{\log r}\in[0,\infty].
\]
Nonzero constants and polynomials have order zero. For its nonzero
zeros $a_j$, let
$\lambda=\inf\{s>0:\sum_j|a_j|^{-s}<\infty\}$, with infimum
$+\infty$ for an empty set of admissible exponents and with $\lambda=0$
for a finite divisor. The canonical rank is the least integer $q\ge0$
such that $\sum|a_j|^{-q-1}<\infty$, or is infinite if none exists.
Write $E_q(w)=(1-w)\exp(w+\cdots+w^q/q)$, with $E_0(w)=1-w$.
If $f=z^m\prod E_q(z/a_j)e^{Q(z)}$ with $Q$ polynomial, its genus is
$\max(q,\deg Q)$; constants have degree convention zero here.
\end{definition}

\begin{proposition}[Growth of the divisor]\label{d12:17:exponent}
For a finite-order entire function, $\lambda\le\rho$ and
$q\le\lfloor\rho\rfloor$. If the nonzero zero sequence is infinite,
its counting function satisfies
$\lambda=\limsup_{r\to\infty}\log\max(1,n(r))/\log r$.
Whenever $q$ is finite, $q\le\lambda\le q+1$.
\end{proposition}
\begin{proof}
For every $s>\rho$, the definition of order and the preceding counting
bound give $n(r)\le C_s r^s$ for large $r$; factoring the center zero
only changes a fixed finite count and a logarithmic term. For $t>s$,
group the nonzero zeros into $2^k\le|a_j|<2^{k+1}$ to bound their
$t$-sum by a constant times $\sum_k2^{-k(t-s)}$. Thus every $t>\rho$
is an admissible exponent, and choosing $t=\lfloor\rho\rfloor+1$
proves the rank assertion. More generally, a convergent $s$-sum bounds
$n(r)$ by $Cr^s$. Conversely any eventual bound $n(r)\le Cr^u$ makes
every $s>u$ sum converge by the same dyadic argument. Taking infima
proves the limsup formula. Minimality of $q$ says that, when $q>0$,
the $q$-sum diverges, whereas the $(q+1)$-sum converges. These statements
give $q\le\lambda\le q+1$; the lower bound is automatic for $q=0$.
Whether the sum converges at an integral endpoint must not be discarded.
\end{proof}

\begin{lemma}[Circles avoiding a canonical divisor]\label{d12:17:lower-product}
Let $\rho<\infty$, $p=\lfloor\rho\rfloor$, and suppose
$n(r)\le C_s r^s$ for every $s>\rho$. For
$P(z)=\prod_jE_p(z/a_j)$ and each $\rho<t<p+1$, there are radii
$r_k\to\infty$ and a constant $C_t$ such that
\[
                   \log|P(z)|\ge-C_t r_k^t\quad (|z|=r_k).
\]
The empty or finite product is included.
\end{lemma}
\begin{proof}
Choose $\rho<s<t$. For a large $R$, there are at most $C R^s$ zeros
with $|a_j|\le4R$. Delete from $[R,2R]$ the intervals of radius
$R^{-s}$ about their moduli. The total deleted length is at most a
constant independent of $R$, so some $r\in[R,2R]$ remains for all
sufficiently large $R$. For $|z|=r$ and $|a_j|\le4R$,
\[
 \log|1-z/a_j|\ge-(s+1)\log R-\log4.
\]
Also, for $1\le l\le p<s$, dyadic grouping gives
$\sum_{|a_j|\le4R}|a_j|^{-l}\le C R^{s-l}$: each dyadic block is
bounded by $C2^{k(s-l)}$, a geometric sum. Consequently the logarithm
of the product of these near factors is at least $-C R^s\log R$,
including their exponential terms $z^l/(l a_j^l)$.

For $|a_j|>4R$, $|z/a_j|\le1/2$ and the convergent logarithm series
gives $|\log E_p(w)|\le2|w|^{p+1}$. A dyadic sum over the far zeros
gives
\[
 \sum_{|a_j|>4R}|a_j|^{-p-1}
 \le C R^{s-p-1}\sum_{k\ge0}2^{k(s-p-1)}
 \le C'R^{s-p-1},
\]
because $s<p+1$. Thus the far logarithms have total absolute value at
most $C R^s$. Combining both parts and using
$R^s\log R\le C_tR^t$ proves the bound. Choose $R=2^k$ and one such
radius for each sufficiently large $k$. Normal product convergence,
already ensured by the $(p+1)$-sum, justifies these logarithmic tail
estimates off the zeros.
\end{proof}

\begin{theorem}[Hadamard factorization]\label{d12:17-05}\family{D12.17.05}
Every nonzero entire function of finite order $\rho$ has the factorization
\[
 f(z)=z^m e^{Q(z)}\prod_jE_p(z/a_j),\qquad
 p=\lfloor\rho\rfloor,\qquad\deg Q\le p.
\]
It also has a factorization using its minimal canonical rank $q$, with
exponential polynomial of degree at most $p$; its genus is at most $\rho$.
\end{theorem}
\begin{proof}
The exponent estimate proves convergence of the $(p+1)$-sum, so
\cref{d12:16-01,d12:16-02} gives the entire product $P$ with exactly
the prescribed zeros. The quotient $f/(z^mP)$ is entire and zero-free.
The logarithm criterion \cref{d12:04-07} on $\C$ gives it the form
$e^h$, with $h$ entire. Choose $\rho<s<t<p+1$. On the circles of
the lower-product lemma, the growth definition and $r\ge1$ give
\[
                  \Ree h=\log|f|-m\log r-\log|P|\le C r^t.
\]
Write $h(z)=\sum_{l\ge0}c_lz^l$, $u=\Ree h$. For $l\ge1$ the
Taylor coefficient integral and its conjugate give
$c_l r^l=\pi^{-1}\int_0^{2\pi}u(re^{\iu\theta})e^{-\iu l\theta}d\theta$.
Subtract $B=\max_{|z|=r}u(z)$ inside this integral. The mean-value
property yields
$|c_l|r^l\le2(B-u(0))\le2(Cr^t-u(0))$.
Along the selected radii this tends to zero after division by $r^l$
whenever $l\ge p+1>t$. All such coefficients vanish, so $h=Q$ has
degree at most $p$.

For the minimal rank $q\le p$, absolute convergence of the relevant
coefficient sums gives
\[
 \prod_jE_p(z/a_j)=\prod_jE_q(z/a_j)
 \exp\left(\sum_{l=q+1}^p\frac{z^l}{l}\sum_j a_j^{-l}\right).
\]
Each inner sum converges absolutely since $l\ge q+1$. The added
exponent is a polynomial of degree at most $p$, proving the assertion
about minimal rank and genus. This is the additional growth argument
which ordinary Weierstrass factorization alone does not provide.
\end{proof}

\begin{theorem}[Blaschke zeros and products]\label{d12:17-06}\family{D12.17.06}
The nonzero zeros $a_j$ of a bounded, nonzero holomorphic function on
$\D$ satisfy $\sum_j(1-|a_j|)<\infty$. Conversely, a disk sequence
without interior accumulation satisfying this condition is the exact
zero sequence, with its listed multiplicities, of the bounded product
\[
             B(z)=z^m\prod_j\frac{|a_j|}{a_j}
                                      \frac{a_j-z}{1-\bar a_jz}.
\]
Here $m\ge0$ prescribes the center multiplicity; the product converges
normally and $|B|\le1$.
\end{theorem}
\begin{proof}
Factor the center zero as $f=z^m g$, $g(0)\ne0$. Jensen on radii
tending to one gives, for every finite subset of nonzero zeros,
$\sum\log(1/|a_j|)\le\log\|f\|_\infty-\log|g(0)|$; the term
$m\log r$ tends to zero. Taking increasing finite subsets proves
convergence, and $1-x\le-\log x$ gives the stated condition.
For sufficiency write the nonzero factor as $b_a$. Direct subtraction
gives, for $|z|\le r<1$,
\[
 |1-b_a(z)|\le(1-|a|)\frac{1+r}{1-r}.
\]
The normal product criterion therefore gives convergence and nonvanishing
away from its factors' zeros. Removing any finitely many local zero
factors leaves a nonvanishing limit, proving exact multiplicities. At
zero, $\prod|a_j|>0$: eventually $|a_j|\ge1/2$ and
$-\log|a_j|\le2(1-|a_j|)$, so this scalar product is nonzero.
Finally $|b_a(z)|\le1$, by
$|1-\bar az|^2-|a-z|^2=(1-|a|^2)(1-|z|^2)\ge0$.
Finite products and then their limit have modulus at most one.
\end{proof}

\begin{corollary}[Finite-order value consequences]\label{d12:17-07}\family{D12.17.07}
A nonconstant finite-order entire function omitting $a$ is
$a+e^{Q}$ for a nonconstant polynomial $Q$; it omits no other finite
value. A nonconstant entire function of finite nonintegral order assumes
every finite value infinitely often.
\end{corollary}
\begin{proof}
Subtracting a constant preserves the order of a nonconstant entire
function: the two maximum moduli differ by at most that constant's
modulus, and both tend to infinity by Liouville. Hadamard applied to
the zero-free $f-a$ gives the first representation. For $b\ne a$, solve
$Q(z)=\log(b-a)+2\pi\iu n$, using any one logarithm value and the
fundamental theorem of algebra. The infinitely many right sides give
distinct preimages, in particular excluding a second omitted value.
If instead $f-b$ has only finitely many zeros, its factorization is
$A(z)e^{Q(z)}$ with $A$ polynomial. For $\deg Q=l>0$, its order is
$l$: the upper bound follows from $|Q(z)|\le Cr^l$; on a ray making
the leading term of $Q$ positive real, the lower bound is
$\log|Ae^Q|\ge c r^l+\deg(A)\log r-C$ for large $r$.
If $Q$ is constant the function is polynomial and has order zero.
Thus finitely many preimages of $b$ would force integral order. This
contradicts nonintegrality, and proves the last assertion without using
the later Picard chapter.
\end{proof}
\source{CON XI.1--3, printed pp.280--290 (PDF pp.294--304), including
the center, meromorphic and Blaschke variants listed as exercises there;
SS V.5, printed pp.147--152 (PDF pp.166--171), for the avoiding-circle
growth route. The proof here supplies its own lower-product estimates
and real-part coefficient argument; an order infimum is never assumed
to be an attained exact growth bound.}

% END chapters/17

% BEGIN chapters/19
\chapter{Analytic continuation and monodromy}\label{ch:d12-19}

All base domains in this chapter are plane domains. Paths are continuous;
piecewise differentiability is needed only for the displayed period integrals.
An open connected subset of the plane is polygonally connected: the set of
points reachable from a fixed point by polygonal paths is open, and its
complement is open, because a small disk can be appended to any such path.
Connectedness therefore makes this set the whole domain.

\begin{definition}[Elements, germs and continuation]\label{d12:19-01}\family{D12.19.01}
A function element is a holomorphic function on a connected open set. Two
elements near $a$ have the same germ $[f]_a$ if they agree on a neighborhood
of $a$. Along $\gamma:[0,1]\to\Omega$, a continuation of $\varphi$ is a
family $\varphi_t$ of germs at $\gamma(t)$, with $\varphi_0=\varphi$, which
is locally represented by one holomorphic function: near each parameter $s$
there are an open set $U$, a holomorphic $f$ on $U$, and a parameter interval
$I$ containing $s$ such that $\gamma(I)\subset U$ and
$\varphi_t=[f]_{\gamma(t)}$ for $t\in I$.
\end{definition}

\begin{theorem}[Uniqueness along a fixed path]\label{d12:19:uniqueness}
Continuation, if it exists, is unique after its initial germ is fixed.
It can be reversed or concatenated, and is unchanged by an increasing
homeomorphic reparametrization of the parameter interval.
\end{theorem}
\begin{proof}
For two continuations consider the set $A$ of parameters where their germs
agree. At a parameter $s$, represent both families on a common connected
disk about $\gamma(s)$, after shortening the parameter interval. If they
agree at one parameter in that interval, the representing functions agree
on a neighborhood of the corresponding point and hence everywhere on the
disk by \cref{d12:05-04}. Thus either the entire shortened interval is in
$A$, or none of it is. Both $A$ and its complement are open. Since $0\in A$
and $[0,1]$ is connected, $A=[0,1]$. Reversal and reparametrization compose
the germ family with the corresponding parameter map. For concatenation,
the two representatives at the joining germ agree on a smaller disk; this
supplies the required local representative across the joining parameter.
\end{proof}

\begin{proposition}[Finite chains and continuous germ families]\label{d12:19-02}\family{D12.19.02}
Continuation is equivalent to the existence of a partition
$0=t_0<\cdots<t_N=1$, connected open sets $U_j$ containing
$\gamma([t_{j-1},t_j])$, and holomorphic functions $f_j$ on $U_j$, such
that $f_j$ and $f_{j+1}$ agree on the component of $U_j\cap U_{j+1}$
containing $\gamma(t_j)$. The initial and final germs are those of $f_1$
and $f_N$. It is also equivalent to a continuous lift into the germ space
whose basic open sets are $\{[f]_z:z\in U\}$.
\end{proposition}
\begin{proof}
The local parameter intervals in the definition cover $[0,1]$. A finite
subcover has a Lebesgue number: if none existed, intervals of lengths
tending to zero not lying in a member would have centers converging to a
point in one member, a contradiction. Choose a partition finer than this
number. Representatives on consecutive subintervals agree as germs at
the shared endpoint, hence agree on its overlap component by the identity
theorem. Components containing the compact subpaths may replace the
original representing sets. Conversely a chain represents the germ family
inside each subinterval. At a joint, choose a disk in the overlap component;
continuity of the path and agreement of the functions give a representative
across the joint. Basic germ neighborhoods intersect in open sets because
equality of germs persists in a neighborhood. On each such neighborhood
the projection to the base is a homeomorphism. Local representation is
therefore exactly continuity of the lift. This assertion concerns the germ
topology, and does not assert that its projection is a covering.
\end{proof}

\begin{proposition}[Permanence of relations]\label{d12:19-03}\family{D12.19.03}
Derivatives, sums and products continue by differentiating, adding and
multiplying the local representatives. A holomorphic identity between
finitely many simultaneously continued functions and their derivatives,
with single-valued holomorphic coefficients on the base, persists along
the path. Quotients are included wherever the denominator is nonzero.
\end{proposition}
\begin{proof}
Representatives which agree on an open overlap have equal derivatives
there, so the proposed operations define continuations. Substitute them
in the asserted identity. Its left side is holomorphic in each representing
neighborhood and vanishes in the first one. At a joint it vanishes on a
nonempty open overlap, and the identity theorem makes it vanish in the
next connected neighborhood. Finite induction through the chain proves
the claim. This propagates an existing relation; it supplies no existence
theorem for a differential equation or for a missing continuation.
\end{proof}

\begin{lemma}[Stability of a continuation chain]\label{d12:19-04}\family{D12.19.04}
If $\varphi$ continues along $\gamma$, there is $\delta>0$ such that every
path $\eta$ with the same endpoints and
$\sup_t|\eta(t)-\gamma(t)|<\delta$ has a continuation with the same
terminal germ. With a varying terminal point, the terminal germ is that
of the same final representative at the new point.
\end{lemma}
\begin{proof}
Choose a finite chain as above. At each joint choose a closed disk centered
at $\gamma(t_j)$ contained in the overlap component where the two functions
agree. Refine the partition if necessary so that each closed subpath is a
compact subset of its representing open set. The distances of these
finitely many compact sets to their respective closed complements are
positive; also the joint disks have positive radii. Take $\delta$ smaller
than all these distances and half all these radii. Then $\eta$ follows
the same representing sets, and its joint points remain in the chosen
agreement disks. Those same functions supply its continuation. The first
germ and, for a fixed endpoint, the last germ are unchanged.
\end{proof}

\begin{theorem}[Homotopy invariance]\label{d12:19-05}\family{D12.19.05}
Let $H:[0,1]^2\to\Omega$ be continuous, with $H(0,s)=a$ and $H(1,s)=b$.
Suppose the germ $\varphi$ at $a$ admits continuation along every path
$t\mapsto H(t,s)$. All the resulting terminal germs at $b$ agree.
\end{theorem}
\begin{proof}
Uniform continuity on the square says that $s\mapsto H(\cdot,s)$ is
continuous in the uniform norm. At each fixed $s$, the preceding stability
lemma gives a neighborhood of $s$ on which the terminal germ is constant.
A locally constant function on the connected interval $[0,1]$ is constant:
the inverse image of one value and its complement are disjoint open sets.
This proves the assertion, including the endpoints of the homotopy.
\end{proof}

\begin{theorem}[Monodromy]\label{d12:19-06}\family{D12.19.06}
Suppose $\Omega$ is simply connected and a holomorphic germ at $a\in\Omega$
admits continuation along every path starting at $a$. It is the germ of a
unique holomorphic function on $\Omega$.
\end{theorem}
\begin{proof}
Here simple connectivity is used in its fixed-endpoint path formulation:
two paths with the same endpoints are homotopic with endpoints fixed.
For the loop formulation this follows by concatenating one path with the
reverse of the other and contracting the resulting loop; backtracking
pieces contract by shortening their traversed parameter interval.
By homotopy invariance the terminal germ at $z$ is independent of the
chosen path from $a$ to $z$. Define $f(z)$ to be its value. Fix one such
path and a final representative $g$ on a disk about $z$. Appending the
straight segment to any sufficiently close point $w$ gives terminal germ
$[g]_w$, so $f=g$ throughout that disk. Hence $f$ is holomorphic. It has
the prescribed initial germ, and uniqueness follows from the identity
theorem on the connected domain.
\end{proof}

\begin{theorem}[Monodromy and obstructions]\label{d12:19-07}\family{D12.19.07}
A logarithm germ at a point of $\C^*$ continues along every path in
$\C^*$. Around a piecewise $C^1$ loop $\gamma$ its change is
$2\pi\iu\ind(\gamma,0)$. An $n$th-root germ changes by the factor
$\exp(2\pi\iu\ind(\gamma,0)/n)$. Simple connectivity alone does not
guarantee continuation. Moreover
\[
                         F(z)=\sum_{n=1}^\infty z^{n!}
\]
is holomorphic on $\D$ and has the unit circle as a natural boundary.
\end{theorem}
\begin{proof}
Cover the compact image of a path in $\C^*$ by disks avoiding zero. Local
logarithms exist by \cref{d12:02-05}. A finite subdivision lets us adjust
each consecutive logarithm by an integral multiple of $2\pi\iu$ to agree
at the joint, hence on its overlap disk. This constructs continuation.
Its derivative is $1/z$; the contour FTC gives the asserted change on a
loop. Exponentiating a continued logarithm divided by $n$ proves the root
formula. In particular the root returns precisely when the index is
divisible by $n$. The germ of $(1-z)^{-1}$ at zero cannot continue as a
holomorphic germ through $1$, although the proposed larger base $\C$ is
simply connected: an extension there would agree with $(1-z)^{-1}$ on a
punctured disk, and multiplication by $1-z$ would give $0=1$ at $1$.

For the series, $n!\ge n$ makes it normally convergent on each compact
subdisk, so its sum is holomorphic. If $\xi$ is a root of unity of order
$q$, then $\xi^{n!}=1$ for every $n\ge q$. The finitely many preceding
terms stay bounded as $r\uparrow1$, whereas
$\sum_{n\ge q}r^{n!}\to+\infty$: for every integer $L$, its first $L$
terms tend to $L$. Thus $F(r\xi)$ is unbounded near $\xi$, excluding a
holomorphic extension through $\xi$. Roots of unity are dense in the
circle. An extension across any point would agree with $F$ on the
connected intersection of a small disk with $\D$; a still smaller disk
contains a root of unity and would give a bounded extension there, which
we have just excluded. This proves the natural-boundary assertion.
\end{proof}
\source{CON IX.1--3, printed pp.210--221 (PDF pp.224--235),
especially printed pp.217--220 (PDF pp.231--234);
FOR \S7.1--7.8, printed pp.44--48 (PDF pp.53--57).
The all-intermediate-path hypothesis is retained explicitly; no surface
uniformization or future covering theorem is used.}

% END chapters/19

% BEGIN chapters/20
\chapter{Riemann surfaces and finite branched covers}\label{ch:d12-20}

\begin{definition}[Surfaces and atlases]\label{d12:20-01}\family{D12.20.01}
A Riemann surface is a nonempty connected, Hausdorff, second-countable topological
space with an atlas of homeomorphisms to plane domains whose transition
maps are holomorphic with holomorphic inverses. Compatible atlases define
the same complex structure. Domains in $\C$, with the identity chart,
are examples. Disconnected open subsets are unions of such surfaces,
not surfaces under the connected convention.
\end{definition}
The union of compatible atlases is again an atlas: the added transitions
are precisely the required compatible maps. Taking all compatible charts
therefore gives a uniquely determined maximal atlas. Connected surfaces
are path connected, by applying the open-and-closed reachability argument
in coordinate disks.

\begin{proposition}[The sphere and the projective line]\label{d12:20-02}\family{D12.20.02}
The quotient $\mathbb{CP}^1=(\C^2\setminus\{0\})/\C^*$ is a compact
Riemann surface. Its charts $[1:z]\mapsto z$ and $[w:1]\mapsto w$ identify
it with $\Chat$, with infinity coordinate $w=1/z$.
\end{proposition}
\begin{proof}
The chart domains are open because their inverse images in $\C^2\setminus0$
are given by a nonzero coordinate. Coordinate ratios are continuous and
constant on scalar orbits, so the quotient definition makes the chart
maps continuous; the displayed parametrizations give continuous inverses.
The transition is $w=1/z$. For Hausdorffness, normalize vectors to the
unit sphere $S^3$ and give phase orbits the distance
$\min_{|\alpha|=1}\|u-\alpha v\|$. The minimum exists, is zero exactly
on equal orbits, and the triangle inequality follows by multiplying the
two minimizing phases. The orbit map is continuous. Its compact quotient
maps bijectively onto this metric space and hence homeomorphically, since
images of closed subsets of a compact space are compact and closed in a
metric space. Thus the original quotient is compact and Hausdorff.
The two charts give a countable base, and their overlapping connected
images give connectedness. Finally the second chart says that a
neighborhood of infinity contains the complement of a sufficiently large
disk, exactly the topology of $\Chat$.
\end{proof}

\begin{theorem}[Maps, multiplicity and local form]\label{d12:20-03}\family{D12.20.03}
A continuous map between surfaces is holomorphic when its coordinate expressions are
holomorphic. This condition is chart independent. A nonconstant map has
at each point $a$ a positive integer multiplicity $m_a$ and coordinates
centered at $a$ and its image in which it is $t\mapsto t^{m_a}$.
It is locally biholomorphic exactly when $m_a=1$.
\end{theorem}
\begin{proof}
Changing charts composes with biholomorphic transition maps, so the chain
rule proves coordinate independence. A map constant on a neighborhood is
constant everywhere: the set of points having a neighborhood where it
equals that constant is open. At a point of its closure, continuity puts
the image in a coordinate disk about that constant; the local identity
theorem makes the map constant on a whole source disk there. This set is
therefore also closed, and connectedness makes it the entire surface.
For a nonconstant map its centered coordinate expression is not identically
zero. Its zero order is finite, and \cref{d12:05-05} gives the asserted
normal form. A target change multiplies the centered expression by a
holomorphic unit; a source change has a simple zero. Neither changes the
order. The power map is locally invertible for $m=1$ and is not injective
on any neighborhood of zero for $m>1$, proving the last assertion.
\end{proof}

\begin{theorem}[The continuation surface]\label{d12:20-04}\family{D12.20.04}
The germs obtainable by continuation of one holomorphic plane germ form
a Riemann surface $S$. The projection $p:S\to\Omega$ is locally
biholomorphic, and evaluation $[f]_z\mapsto f(z)$ is holomorphic.
The full space of all germs need not be second countable; the assertion
is about this single continuation component.
\end{theorem}
\begin{proof}
Use the basic germ neighborhoods of \cref{d12:19-02}. On one such
neighborhood projection is bijective to its representing domain, with
continuous inverse $z\mapsto[f]_z$. For germs over different points,
disjoint base disks separate them. For distinct germs over one point,
choose representatives $f,g$ on a common connected disk. If any germs
of $f$ and $g$ there agreed, the identity theorem would give $f=g$ on
the disk. Thus their basic neighborhoods are disjoint. The germ space
is Hausdorff. These projection charts have identity transitions on their
overlaps, and evaluation is represented by $f$.

The set $S$ is the path component of the initial germ: continuations give
paths in the germ topology, and a path in that topology gives a
continuation by local projection charts. It is open because each such
chart disk is path connected. To prove second countability, take the
countable collection of rational-center rational-radius disks compactly
contained in $\Omega$. Every finite continuation chain can be refined
to a chain of these disks. Given a finite word of disks starting at the
initial point, there is at most one sequence of representing functions:
the first is fixed by its initial germ, and each next one is fixed by
agreement on the nonempty connected intersection of two disks. There
are countably many finite words. The resulting final disk charts cover
$S$, and each chart has a countable base. Their bases together are a
countable base of $S$. No covering conclusion was needed for this proof.
\end{proof}

\begin{theorem}[Covering and lifting bridges]\label{d12:20-05}\family{D12.20.05}
A covering map is a map $p:E\to B$ for which every point of $B$ has a
neighborhood whose inverse image is a disjoint union of open sets mapped
homeomorphically onto that neighborhood. Paths have unique lifts from a
chosen initial point. A fixed-endpoint path homotopy has a continuous
lift from a chosen initial point, and its lifted endpoints agree.
The maps $\exp:\C\to\C^*$ and $z\mapsto z^n:\C^*\to\C^*$ are
holomorphic coverings, the latter having $n$ sheets.
\end{theorem}
\begin{proof}
Partition a path into finitely many subpaths lying in evenly covered
neighborhoods, using compactness and a Lebesgue number. Starting in its
chosen sheet, use that sheet's inverse on the first subpath; its endpoint
chooses the next sheet. This constructs a lift. Two lifts agreeing at a
parameter agree in a neighborhood by local inverses; their agreement set
and its complement are open, so connectedness gives uniqueness.
For a homotopy $H$, the inverse images $H^{-1}(U)$ of evenly covered
sets form an open cover of its compact parameter square. A finite
subcover has a Lebesgue number in that square's Euclidean metric.
Choose a finite rectangular grid with rectangles of smaller diameter;
each closed rectangle then maps into an evenly covered set. This uses
no metric or uniform-continuity assumption on the base space.
Prescribe the constant initial edge above the fixed initial endpoint.
Fill rectangles row by row, left to right, using the inverse sheet chosen
at the lower left corner. The pre-existing bottom and left edges stay
in that sheet because they are connected; uniqueness of path lifts makes
the formulas agree on shared edges. Finite gluing of continuous maps on
closed rectangles gives a continuous lift of the square. Above the fixed
terminal point its endpoint edge lies in a discrete fibre and is therefore
constant.

On a small disk $U\subset\C^*$ choose a logarithm $L$. The inverse
exponential sheets are $L(U)+2\pi\iu k$, $k\in\Z$; these are disjoint,
and all exponential preimages occur because equal exponentials differ by
$2\pi\iu k$. The inverse power sheets are
$\exp((L+2\pi\iu k)/n)(U)$, $0\le k<n$, with the same assertions.
Their local inverses are holomorphic. A local biholomorphism need not
be a covering: $\exp:\C\setminus\{0\}\to\C^*$ is surjective and locally
biholomorphic, but the path $e^t$, $-\eps\le t\le0$, starting at the
lift $-\eps$, cannot lift to its final point. Uniqueness for the full
exponential covering forces the missing final lift to be zero.
\end{proof}

\begin{theorem}[Compact maximum principle]\label{d12:20-06}\family{D12.20.06}
A scalar holomorphic function on a compact Riemann surface is constant.
\end{theorem}
\begin{proof}
Its modulus attains a maximum. In a coordinate disk through a maximizing
point, the strong maximum principle makes it constant. The identity
theorem, propagated through coordinate disks, makes it constant on the
connected surface. A sphere-valued holomorphic map is a different object;
nonconstant meromorphic functions on a compact surface are not excluded.
\end{proof}

\begin{theorem}[Finite degree]\label{d12:20-07}\family{D12.20.07}
Let $X$ be compact and $Y$ a connected Riemann surface. A nonconstant
holomorphic map $f:X\to Y$ is surjective, has finite fibres and finitely
many ramification points, and has a positive integer degree $d$ satisfying
\[
                         \sum_{a\in f^{-1}(b)}m_a=d\qquad(b\in Y).
\]
Away from its finitely many branch values it is a finite covering.
\end{theorem}
\begin{proof}
The local power form makes $f$ open. Its image is compact and thus closed
in the Hausdorff space $Y$, proving surjectivity. A fibre is closed and
has isolated points by the local power form, so compactness makes it
finite. The critical points are locally the zeros of the derivative of a
nonconstant coordinate expression; they are closed and isolated, and
are consequently finite on compact $X$.
Fix $b$ and its preimages $a_1,\ldots,a_k$. Choose pairwise disjoint
power-coordinate neighborhoods of the $a_j$ and smaller neighborhoods
whose closures lie inside them. The compact complement of the smaller
neighborhoods has image avoiding $b$, so a sufficiently small target
disk has its whole inverse image in the power neighborhoods. Shrink it
again so that each power chart contains all $m_j$ roots of every nonzero
target coordinate in that disk. Thus the sum of fibre multiplicities is
$\sum m_j$ throughout the disk, including its center. It is locally
constant on connected $Y$, hence is one integer $d$. At a non-branch
value all the local powers have exponent one; the same compact-complement
argument supplies an evenly covered disk with $d$ sheets.
\end{proof}

\begin{theorem}[Completion over a puncture]\label{d12:20-08}\family{D12.20.08}
A connected finite holomorphic covering of $0<|z|<r$ with $m$ sheets is
biholomorphically equivalent to $t\mapsto t^m$ on $0<|t|<r^{1/m}$.
It can therefore be completed by one point, with local coordinate $t$,
ramification multiplicity $m$ and projection $z=t^m$.
\end{theorem}
\begin{proof}
A loop in the punctured disk has a continuous argument by exponential
path lifting. Its change is $2\pi k$. Radially deform its modulus to the
base modulus, and interpolate its argument to the linear argument with
change $2\pi k$; these are fixed-endpoint homotopies. Thus its homotopy
class is precisely $k\in\Z$. The monodromy permutation $\tau$ of the
positive generator acts transitively on a fibre: join any two fibre
points by a path in the connected surface and project it. A transitive
permutation generated by one element on $m$ points is one $m$-cycle.

Choose corresponding initial points in the power cover and the given
cover. For $t$ in the punctured parameter disk, project any path from
the initial parameter to $t$ by $t\mapsto t^m$, and lift it in the given
cover. Its endpoint defines $\Phi(t)$. Changing the parameter path
adds a loop whose projected winding is a multiple of $m$, so $\tau^m=1$
and homotopy lifting make this endpoint independent of the choice.
The map is locally a composition of holomorphic inverse sheets. It is
onto because any path in the given surface projects to a path that also
lifts in the power cover. Two parameter points with the same image
differ by an $m$th root of unity; turning from one to the other projects
to the corresponding power of $\tau$. Such a power fixes a fibre point
only when the parameter points coincide. Hence $\Phi$ is a bijective
local biholomorphism, with holomorphic inverse. Adjoining the missing
center in the parameter disk now supplies the asserted completion.
For a disconnected finite covering, apply this argument to each
monodromy cycle separately, adding one point per cycle.
\end{proof}
\source{CON IX.6, printed pp.233--236 (PDF pp.247--250);
FOR \S1--2, printed pp.2--12 (PDF pp.11--21), and \S7, printed pp.44--48
(PDF pp.53--57), and \S8, printed pp.48--51 (PDF pp.57--60);
AHL VIII.1.2--1.7, printed pp.284--300 (PDF pp.300--316).
Hausdorffness and second countability are proved for the continuation
component, and finite-cover lifting and puncture completion are proved here.}

% END chapters/20

% BEGIN chapters/21
\chapter{Algebraic branches and their surface models}\label{ch:d12-21}

Polynomial coefficients are complex. A branch means an actual holomorphic
function on a stated open set, not an arbitrary pointwise choice of roots.

\begin{theorem}[A simple root varies holomorphically]\label{d12:21-01}\family{D12.21.01}
Let $P(z,w)$ be a polynomial in $w$ whose coefficients are holomorphic
near $a$, with $P(a,b)=0$ and
$P_w(a,b)\ne0$. There are neighborhoods on which its zeros near $b$
are precisely $w=f(z)$, with $f$ holomorphic and $f(a)=b$.
For a polynomial of fixed positive $w$-degree having distinct roots at
$a$ and nonzero leading coefficient there, all its roots are represented
by finitely many disjoint such branches on one disk about $a$.
\end{theorem}
\begin{proof}
Choose a small circle $C$ about $b$ on which $P(a,\cdot)$ is nonzero
and inside which $b$ is its only zero, of multiplicity one. On $C$ the
minimum modulus is positive. Uniform continuity permits a disk about
$a$ on which the perturbation is smaller than this minimum. Rouch\'e's
theorem gives exactly one zero inside $C$ for each parameter. Its value is
\[
                  f(z)=\frac1{2\pi\iu}\int_C
                       w\frac{P_w(z,w)}{P(z,w)}\,dw.
\]
The integrand and all its parameter derivatives are continuous and
uniformly bounded on smaller parameter disks times $C$, because the
denominator has a positive lower bound. Difference quotients, or Cauchy's
formula in the parameter, justify differentiation under this fixed
contour. Thus $f$ is holomorphic. There are no other zeros near $b$ by
the root count. For a polynomial take disjoint circles around all the
roots. Keep its leading coefficient nonzero after shrinking the parameter
disk. The fundamental theorem of algebra then counts its degree-many
roots, so the branches just constructed exhaust them.
\end{proof}

\begin{theorem}[The finite exceptional set]\label{d12:21-02}\family{D12.21.02}
Let $P\in\C[z,w]$ be irreducible and have $w$-degree $d>0$, with leading
coefficient $a_d(z)$. There is a nonzero polynomial $\Delta(z)$, the
resultant of $P$ and $P_w$, such that outside
\[
                    E=\{a_d\Delta=0\}\subset\C
\]
the polynomial has exactly $d$ distinct finite roots. In particular $E$
is finite. The point at infinity must be considered separately.
\end{theorem}
\begin{proof}
The coefficients of $P$ have no nonconstant common divisor in $\C[z]$,
since such a divisor would factor $P$. Gauss's lemma therefore makes
$P$ irreducible in $K[w]$, $K=\C(z)$. Here the needed form of that lemma
follows by clearing denominators in a factorization and canceling common
coefficient factors: for each irreducible $q\in\C[z]$, reduction modulo
$q$ shows that a product of two primitive polynomials is primitive,
because $(\C[z]/(q))[w]$ is an integral domain. Cancel the powers of
each $q$ in the cleared factorization to obtain polynomial factors.
Since the characteristic is zero, $P_w\ne0$ and has smaller degree;
thus $P$ and $P_w$ are coprime in $K[w]$.

Consider the linear map
\[
 K[w]_{<d-1}\oplus K[w]_{<d}\longrightarrow K[w]_{<2d-1},
                         \qquad(A,B)\longmapsto AP+BP_w.
\]
It is injective: a kernel relation and coprimality give $P\mid B$, so
$B=0$ by degree, and then $A=0$. Equal finite dimensions make its
determinant nonzero. Its matrix has entries in $\C[z]$, so that determinant
$\Delta$ is a nonzero polynomial. If at a parameter outside its zero set
$P$ and $P_w$ had a common root, evaluation at that root would annihilate
the whole image of the specialized map, contradicting its surjectivity
and the presence of the constant polynomial $1$ in its target. Outside
$a_d=0$ the degree stays $d$; absence of a common root is exactly
simplicity of every root. This proves the assertion without an unproved
resultant existence theorem. For $d=1$ the first summand is zero and
the same determinant argument applies.
\end{proof}

\begin{theorem}[Continuation and monodromy permutations]\label{d12:21-03}\family{D12.21.03}
Over $\Omega=\C\setminus E$, the set
$S=\{(z,w):z\in\Omega,\ P(z,w)=0\}$ is a $d$-sheeted holomorphic
covering of $\Omega$. Every branch continues along every path there.
A loop permutes the $d$ root germs, and homotopic loops give the same
permutation. Concatenation corresponds to successive application of the
permutations, with the first traversed loop acting first.
\end{theorem}
\begin{proof}
At a base point the preceding local root theorem writes the entire inverse
image of a small disk as the disjoint graphs of its $d$ branches. These
are covering charts and also give the topology and complex charts of
$S$; as a subspace of $\C^2$, $S$ is Hausdorff and second countable.
Covering path and homotopy lifting in \cref{d12:20-05} give the asserted
continuation and permutation rules. Reversing a loop reverses its
permutation, so every such correspondence is a permutation, not merely
a map of the fibre into itself.
\end{proof}

\begin{theorem}[Local cyclic parameters]\label{d12:21-04}\family{D12.21.04}
Choose a punctured disk about $c\in E$ containing no other exceptional
point. Each cycle of length $h$ in its local monodromy has a single-valued
parameter $t$ in which $z=c+t^h$ and its root is a holomorphic function
$w=W(t)$ for $t\ne0$. At infinity the corresponding formula is
$z=t^{-h}$.
\end{theorem}
\begin{proof}
Restrict the finite covering to the punctured disk. Its connected
components are exactly the monodromy cycles: a path between fibre points
projects to a loop, and lifted loops give all such paths. Apply the
punctured-disk classification \cref{d12:20-08} to each component.
The coordinate function $w$ pulled back by that biholomorphism is
holomorphic on the punctured parameter disk. At infinity apply the same
argument to the base coordinate $1/z$. This proves the formulas; it does
not yet assert removability of $W$ at the missing center.
\end{proof}

\begin{theorem}[Algebraic growth and finite Laurent tails]\label{d12:21-05}\family{D12.21.05}
In the preceding parameters, $W$ is meromorphic at $t=0$. In particular
it has a convergent Laurent expansion with only finitely many negative
powers. If $a_d(c)\ne0$, the finite-point branches are bounded and
$W$ is holomorphic at zero.
\end{theorem}
\begin{proof}
Every root of $w^d+b_{d-1}w^{d-1}+\cdots+b_0$ satisfies
$|w|\le1+\max_j|b_j|$. Indeed, if $|w|>1+B$, the sum of the moduli of
the lower terms is at most
$B(|w|^d-1)/(|w|-1)<|w|^d$, excluding a zero.
Near a finite point $c$, if $a_d$ has zero order $k$, the coefficients
$a_j/a_d$ are bounded by $C|z-c|^{-k}$ after reducing the disk. Thus
$|W(t)|\le C'|t|^{-hk}$. Multiplication by $t^{hk}$ gives a bounded
holomorphic function on a punctured disk, which extends by
\cref{d12:07-03}. If $k=0$, $W$ itself extends. At infinity the rational
coefficients $a_j/a_d$ grow at most as $C|z|^L$ for an integer $L\ge0$;
with $z=t^{-h}$ the same argument gives a pole of order at most $hL$.
The Laurent theorem now gives the claimed convergent finite negative
tail. This is an elementary fractional-power description; no formal
series convergence or general normalization theorem has been assumed.
\end{proof}

\begin{theorem}[The symmetric-coefficient converse]\label{d12:21-06}\family{D12.21.06}
Suppose a finite collection of $d$ analytic branches over the complement
of a finite set in $\Chat$ is invariant under continuation, and every
branch has at most power growth near each missing point in a finite
cyclic parameter. Then all branches satisfy one nonzero equation in
$\C[z,w]$.
\end{theorem}
\begin{proof}
On a small ordinary base disk form $\prod_{j=1}^d(w-f_j(z))$.
Continuation only permutes the factors, so its elementary symmetric
coefficients glue to single-valued holomorphic functions off the missing
points. In a punctured disk, finitely many cyclic parameters and their
power bounds give one common bound $C|z-c|^{-N}$ for every branch,
after increasing the integer $N$. Each coefficient, being a finite sum
of products of at most $d$ branches, has a power bound too. The removable
theorem after multiplication by a sufficiently high power makes it
meromorphic at $c$. The same reasoning in $1/z$ works at infinity.
Every coefficient is therefore meromorphic on $\Chat$ and is rational
by \cref{d12:07-06}. Multiplying the monic equation by a common
polynomial denominator yields the required polynomial equation, whose
leading coefficient is that nonzero denominator.
\end{proof}

\begin{corollary}[Irreducibility and connected branches]\label{d12:21:transitivity}
For the irreducible polynomial of \cref{d12:21-02}, the monodromy action
on its roots is transitive, so the covering $S$ is connected. Its finite
cyclic completions over $E$ and infinity form a compact Riemann surface
on which both $z$ and $w$ are meromorphic.
\end{corollary}
\begin{proof}
An orbit of the monodromy action is invariant under continuation. Its
root growth is bounded by the preceding theorem, so the product of its
linear factors has rational coefficients by the symmetric-coefficient
argument. The complementary orbit union has the same property. A proper
orbit would factor the monic polynomial $P/a_d$ in $\C(z)[w]$, contrary
to Gauss's lemma and irreducibility. Hence there is one orbit; path
lifting then joins any point of $S$ to the chosen initial fibre.
Attach one parameter center for each local cycle at each exceptional
point and at infinity. The finitely many attached disk charts have
holomorphic transitions to $S$, and the projection is the stated local
power. Distinct added points are separated by their disjoint punctured
cycle neighborhoods; an old point is separated from added points by
smaller base disks. This proves Hausdorffness. The old countable base
together with finitely many countable disk bases proves second
countability. The growth theorem extends $w$ meromorphically there.
Finally take small closed base disks about the exceptional values.
Their lifted closed parameter disks are compact. Above the remaining
compact base set, the root bound bounds $w$, so its preimage is a closed
bounded subset of $\C^2$ and is compact. These finitely many compact
sets cover the completed surface. This constructs this specific
finite-branch completion, not a universal normalization or algebraization
theorem for arbitrary analytic spaces.
\end{proof}

\begin{theorem}[Explicit smooth and singular examples]\label{d12:21-07}\family{D12.21.07}
The equation $w^n=z$ has the smooth parameter $t=w$, $z=t^n$; its compact
completion is $\mathbb{CP}^1$, branched with multiplicity $n$ at zero
and infinity. If $P(z)$ is square-free of degree $d\ge1$, the affine
curve $w^2=P(z)$ is smooth. Its compact surface has two unramified
points over infinity when $d$ is even and one point of multiplicity two
when $d$ is odd. These statements do not turn singular affine curves
into smooth embedded curves.
\end{theorem}
\begin{proof}
For $w^n=z$ the displayed parameter and inverse are polynomial. At
infinity $s=1/t$ gives $1/z=s^n$, proving the completion and multiplicity.
For $w^2=P(z)$, where $w\ne0$ the local root theorem gives $w$ as a
holomorphic function of $z$. At a root $a$ of $P$, square-freeness gives
$P'(a)\ne0$, so $z=P^{-1}(w^2)$ and $w$ is a local coordinate; the
projection has order two there.

For $d=2g+2$, put $t=1/z$, $v=wt^{g+1}$. The equation becomes
$v^2=t^{2g+2}P(1/t)=Q(t)$, with $Q(0)\ne0$. Its two local square
roots give two separate points at $t=0$, each with coordinate $t$ and
unramified projection. For $d=2g+1$, choose the even holomorphic root
$u(t)$ of $t^{4g+2}P(t^{-2})$, with $u(0)^2$ equal to the leading
coefficient. Then
\[
                  z=t^{-2},\qquad w=t^{-(2g+1)}u(t)
\]
parametrizes both sheets near infinity. Opposite parameters have the
same $z$ and opposite $w$, so exactly one parameter represents each
point on those sheets. It supplies one missing center and projection
$1/z=t^2$. The local formulas and the bounded-root compact-core
argument above prove compactness and the surface properties.

For comparison, the cusp $y^2=x^3$ has the parametrization
$(x,y)=(t^2,t^3)$. It is not a complex submanifold at the origin:
in any nonconstant local parametrization the identity
$2\ord y=3\ord x$ forces both derivatives to vanish, excluding the
nonzero differential required of an embedded manifold chart. The node
$y^2=x^2(x+1)$ has two distinct local branches
$y=\pm x\sqrt{1+x}$; deleting their common origin disconnects a
sufficiently small neighborhood, unlike a punctured disk. Separate
branch parameters give smooth surface models but do not make either
singular affine subset already a Riemann surface.
\end{proof}
\source{AHL VIII.2.1--2.3, printed pp.300--306 (PDF pp.316--322):
local root contour, resultant, continuation, power growth and symmetric
rational coefficients were compared directly. FOR \S8, printed
pp.48--51 (PDF pp.57--60), supplies the finite-branch comparison.
Leading-coefficient zeros, infinity and singular affine examples are
kept distinct.}

% END chapters/21


% BEGIN chapters/22
\chapter{Elliptic functions, tori and the associated cubic}\label{ch:d12-22}

\begin{theorem}[Lattice geometry and torus charts]\label{d12:22-01}\family{D12.22.01}
Let $\Lambda=\Z\omega_1+\Z\omega_2$, where $\omega_1\ne0$ and
$\Imm(\omega_2/\omega_1)>0$. Then $\Lambda$ is discrete,
$\sum_{\lambda\ne0}|\lambda|^{-s}<\infty$ for $s>2$, and
$T=\C/\Lambda$ is a compact Riemann surface. Its charts are induced by
small plane disks, with translation transition maps.
\end{theorem}
\begin{proof}
The real-linear map $(u,v)\mapsto u\omega_1+v\omega_2$ is invertible.
Its norm has a positive minimum on $|u|+|v|=1$, giving
$|m\omega_1+n\omega_2|\ge c(|m|+|n|)$. There are $4k$ integer pairs
with $|m|+|n|=k$, so the asserted sum is bounded by
$4c^{-s}\sum_{k\ge1}k^{1-s}$. The same inequality proves discreteness
and a positive minimum distance between distinct lattice points.
The formula
$d([z],[w])=\min_{\lambda\in\Lambda}|z-w-\lambda|$ defines a metric:
the minimum is attained because bounded disks contain finitely many
lattice points, and separation and the triangle inequality follow from
this formula. It induces the quotient topology, since the inverse image
of a metric ball is the union of its translated plane balls. For a
radius smaller than one third of the minimum lattice distance, the
quotient map is injective and open on every disk of that radius. These
give complex charts; two local lifts differ by one constant lattice
translation on each overlap component. Images of rational disks form
a countable base. The torus is Hausdorff by its metric and connected
as the image of $\C$. The closed fundamental parallelogram maps onto
it, proving compactness. Replacing a negatively oriented second basis
vector by its negative gives the stated orientation without changing
the underlying lattice.
\end{proof}

\begin{theorem}[Holomorphic elliptic functions and residues]\label{d12:22-02}\family{D12.22.02}
A $\Lambda$-periodic entire function is constant. For a meromorphic
$\Lambda$-periodic function the sum of its residues in a complete set
of classes is zero. In particular it cannot have a single simple pole.
\end{theorem}
\begin{proof}
The entire function is bounded on a closed fundamental parallelogram,
hence everywhere, so Liouville applies. A meromorphic function has
finitely many poles on the torus. Choose a translated parallelogram
with none on its boundary. Such a translate exists: in the real basis
coordinates avoid the countably many coordinate levels of the discrete
pole set in each direction. Opposite side integrals cancel by periodicity.
The residue theorem makes their sum $2\pi\iu$ times the residue sum,
which is therefore zero. A genuine simple pole has nonzero residue,
proving the last assertion.
\end{proof}

\begin{theorem}[Divisor constraints]\label{d12:22-03}\family{D12.22.03}
A nonconstant elliptic function has equally many zeros and poles on
$T$, with multiplicities. If their representatives are $a_1,\ldots,a_N$
and $b_1,\ldots,b_N$, then
$\sum a_j-\sum b_j\in\Lambda$. Every finite value has exactly $N$
preimages counted with multiplicity.
\end{theorem}
\begin{proof}
The elliptic logarithmic derivative $h=f'/f$ has residues equal to the
signed zero and pole orders. Its residue sum is zero by the preceding
theorem, proving equality of the counts. Applying this to $f-c$ proves
the value assertion, since its poles and their orders are unchanged.
Choose a translated positive parallelogram avoiding the entire divisor.
The residues of $zh(z)$ give
\[
 \sum a_j-\sum b_j
 =\frac1{2\pi\iu}\int_{\partial P}zh(z)\,dz
 =\omega_1 n_2-\omega_2 n_1,
\]
where $n_j=(2\pi\iu)^{-1}\int_{a}^{a+\omega_j}h(z)\,dz$.
Each $n_j$ is an integer: the endpoint values of $f$ agree, and that
integral is the winding number of its nonzero image path. The side-pair
calculation has minus $\omega_2$ times the bottom integral and plus
$\omega_1$ times the left upward integral. This establishes the formula
and its sign. Different divisor representatives change the sum only
by lattice vectors.
\end{proof}

\begin{theorem}[The regularized Weierstrass series]\label{d12:22-04}\family{D12.22.04}
The series
\[
 \wp(z)=\frac1{z^2}+\sum_{\lambda\in\Lambda\setminus\{0\}}
              \left(\frac1{(z-\lambda)^2}-\frac1{\lambda^2}\right)
\]
converges normally away from $\Lambda$ and defines a meromorphic
function with precisely a double pole of principal part
$(z-\lambda)^{-2}$ at every lattice point. Its derivative is
$\wp'(z)=-2\sum_{\lambda\in\Lambda}(z-\lambda)^{-3}$, normally and
absolutely convergent away from the poles.
\end{theorem}
\begin{proof}
For $|z|\le R$ and $|\lambda|>2R$, subtraction puts the summand over
the common denominator $\lambda^2(z-\lambda)^2$, and bounds its
modulus by $C_R|\lambda|^{-3}$. The lattice sum of this majorant
converges. The derivative terms are bounded by $C|\lambda|^{-3}$ on
the same compact set, so compact convergence and the derivative theorem
justify termwise differentiation. Near one lattice point remove its
one singular term; the remainder converges holomorphically on a small
disk, proving the exact principal part. The two separate inverse-square
sums have not been reindexed: they need not converge separately.
\end{proof}

\begin{theorem}[Parity and genuine periods]\label{d12:22-05}\family{D12.22.05}
The function $\wp$ is even, $\wp'$ is odd, and both are
$\Lambda$-periodic. The complete period lattice of $\wp$ is exactly
$\Lambda$.
\end{theorem}
\begin{proof}
Reindex the normally convergent regularized series by
$\lambda\mapsto-\lambda$ to obtain evenness. Its derivative is odd.
The absolutely convergent derivative series may be reindexed under
any lattice translation, proving periodicity of $\wp'$. Consequently
$\wp(z+\omega_j)-\wp(z)$ has zero derivative off the lattice.
Its poles cancel, so it extends to a constant entire function. At
$z=-\omega_j/2$ its value is zero by evenness; this point is not a
lattice point, by real independence of the basis. Thus each generator
is a period of $\wp$, hence so is every lattice element. Conversely,
if $\tau$ is a period, translation must preserve its pole set. The pole
at zero then forces $\tau\in\Lambda$.
\end{proof}

\begin{theorem}[Differential equation and branch values]\label{d12:22-06}\family{D12.22.06}
Set $g_2=60\sum_{\lambda\ne0}\lambda^{-4}$ and
$g_3=140\sum_{\lambda\ne0}\lambda^{-6}$. Then
\[
                \wp'^2=4\wp^3-g_2\wp-g_3,\qquad
                \wp''=6\wp^2-g_2/2.
\]
The three nonzero half-period classes $h_1=\omega_1/2$,
$h_2=\omega_2/2$, $h_3=(\omega_1+\omega_2)/2$ are exactly the
zeros of $\wp'$, each simple. Their values $e_j=\wp(h_j)$ are distinct,
and $4X^3-g_2X-g_3=4\prod_{j=1}^3(X-e_j)$. In particular
$g_2^3-27g_3^2\ne0$.
\end{theorem}
\begin{proof}
For $|z|$ smaller than the minimum nonzero lattice modulus, geometric
expansion of each summand, justified by the same compact majorants,
gives
\[
 \wp(z)=z^{-2}+A z^2+Bz^4+O(z^6),\qquad
       A=3\sum\lambda^{-4}=g_2/20,\quad B=5\sum\lambda^{-6}=g_3/28.
\]
Odd lattice sums vanish by absolute convergence and pairing opposite
points. Substitution shows that
$\wp'^2-4\wp^3+g_2\wp+g_3$ has neither a principal part nor a
constant term at zero: its possible $z^{-2}$ coefficient is
$-20A+g_2=0$, and its constant coefficient is $-28B+g_3=0$.
It is elliptic and has no poles, hence is constant, and its value at
zero after removal is zero. Differentiate the equation and cancel
$\wp'$ where it is nonzero. The identity theorem extends the resulting
second equation across its isolated zeros.

Oddness and periodicity give $\wp'(h_j)=0$. These are three distinct
classes. Since $\wp'$ has one triple pole on $T$, the divisor count
allows exactly three zeros with multiplicity. Each named zero is
therefore simple. The elliptic function $\wp-e_j$ has at least a
double zero at $h_j$, and has just one double pole. If $e_j=e_k$ for
$j\ne k$, its zero count would be at least four, a contradiction.
The differential equation now identifies the three distinct roots of
the cubic. A repeated root $r$ of that cubic would give
$g_2=12r^2$, $g_3=-8r^3$, hence $g_2^3=27g_3^2$.
Conversely that equality gives a repeated root: when $g_2=0$ take
$r=0$; otherwise take $r=-3g_3/(2g_2)$ and substitute in both the
cubic and its derivative. Thus distinct roots imply the asserted
nonzero discriminant.
\end{proof}

\begin{theorem}[The torus and its projective cubic]\label{d12:22-07}\family{D12.22.07}
The map $\wp:T\to\Chat$ has degree two and branch values
$e_1,e_2,e_3,\infty$. The projective curve
\[
 E:\quad Y^2Z=4X^3-g_2XZ^2-g_3Z^3\quad\text{in }[Z:X:Y]
\]
is nonsingular, and
\[
 \Phi([z])=[1:\wp(z):\wp'(z)]\ (z\notin\Lambda),\qquad
 \Phi([0])=O=[0:0:1]
\]
is a biholomorphism $T\to E$, including its infinity charts.
\end{theorem}
\begin{proof}
Every finite value of $\wp$ has two preimages with multiplicity by
\cref{d12:22-03}; its double pole supplies the fibre over infinity.
The only finite ramification points are the zeros of $\wp'$, and at
each $h_j$ the zero of $\wp-e_j$ has order exactly two. At zero the
coordinate $1/\wp=z^2+O(z^6)$ also has order two. This proves degree
and the branch-value list.

For clarity the projective topology is literal. Normalize nonzero
vectors of $\C^3$ to $S^5$ and identify unit phases. The distance
$\min_{|\alpha|=1}\|u-\alpha v\|$, as in
\cref{d12:20-02}, gives its Hausdorff quotient topology; compactness
comes from $S^5$. The three nonzero-coordinate charts have continuous
coordinate ratios and continuous inverses, giving a countable base.
The homogeneous equation defines a closed subset.
In the finite chart $Z=1$, write $x=X/Z$, $y=Y/Z$ and
$q(x)=4x^3-g_2x-g_3$. The equation is $y^2=q(x)$.
For $y\ne0$, $x$ is a holomorphic local coordinate, using the simple
root theorem \cref{d12:21-01} in $y$. At $y=0$, $x=e_j$ and
$q'(e_j)\ne0$, so $y$ is a holomorphic local coordinate, solving for
$x$. These charts are compatible because their ambient coordinates
are holomorphic functions of the respective parameters.

There is just one point at infinity: $Z=0$ forces $X=0$, giving $O$.
In its $Y\ne0$ chart put $u=X/Y$, $v=Z/Y$. Its equation is
\[
                     v=4u^3-g_2uv^2-g_3v^3.
\]
Its derivative with respect to $v$ at $(0,0)$ is one. The simple-root
contour theorem therefore gives $v=V(u)$ holomorphically, with
$V(u)=4u^3+O(u^7)$. This proves nonsingularity at infinity and gives
the missing complex chart, without relying on a set-theoretic curve
identification. The finite derivative checks prove nonsingularity
elsewhere.

The Laurent series give
\[
 \frac{\wp(z)}{\wp'(z)}=-\frac z2+O(z^5),\qquad
 \frac1{\wp'(z)}=-\frac{z^3}2+O(z^7).
\]
Thus $\Phi$ extends holomorphically at zero in the $(u,v)$ chart, and
its $u$-coordinate has derivative $-1/2$, giving a holomorphic local
inverse there. At a finite non-half-period its $x$-coordinate has
derivative $\wp'\ne0$. At a half-period its $y$-coordinate has
derivative $\wp''\ne0$, since the corresponding zero of $\wp'$ is
simple. Hence it is locally biholomorphic everywhere.

Finally $\wp(z)=\wp(w)$ implies $z\equiv w$ or $z\equiv-w$:
for a non-half-period the two displayed classes exhaust the degree-two
fibre; at a half-period its double zero exhausts that fibre. Equality
also of the odd derivative distinguishes the two classes unless that
derivative is zero, when the classes already agree. The only pole
class is zero. This proves injectivity of $\Phi$. For any finite
$(x,y)\in E$, choose $z$ with $\wp(z)=x$. The differential equation
gives $\wp'(z)=\pm y$, and replacing $z$ by $-z$ chooses the required
sign. Infinity is already attained. The bijective local biholomorphism
therefore has a holomorphic inverse on all of $E$. In particular $E$
is connected and compact; this proves precisely the lattice-associated
cubic assertion, not converse uniformization of arbitrary cubics.
\end{proof}

\begin{theorem}[Zeta, sigma and Legendre]\label{d12:22-08}\family{D12.22.08}
The normally convergent expressions
\[
 \zeta(z)=\frac1z+\sum_{\lambda\ne0}
       \left(\frac1{z-\lambda}+\frac1\lambda+\frac z{\lambda^2}\right),
 \qquad \sigma(z)=z\prod_{\lambda\ne0}E_2(z/\lambda)
\]
satisfy $\zeta'=-\wp$ and $\sigma'/\sigma=\zeta$.
Both are odd; $\sigma$ is entire with simple zeros precisely $\Lambda$.
There are constants $\eta_j$ such that
\[
 \zeta(z+\omega_j)=\zeta(z)+\eta_j,\qquad
 \sigma(z+\omega_j)=-e^{\eta_j(z+\omega_j/2)}\sigma(z),
 \qquad \eta_1\omega_2-\eta_2\omega_1=2\pi\iu.
\]
\end{theorem}
\begin{proof}
On a fixed compact set the regularized zeta summands are
$O(|\lambda|^{-3})$. The product estimate of \cref{d12:16-02} and
the same convergent lattice sum prove normal product convergence.
Its exact simple zeros and nonvanishing elsewhere follow from
\cref{d12:16-01}. Logarithmic differentiation on compact sets away
from the zeros gives the displayed zeta series; differentiating that
series gives $-\wp$. Reindexing the absolutely convergent product
by $\lambda\mapsto-\lambda$ proves oddness of $\sigma$, and its
logarithmic derivative is odd as well.
Since $\wp$ is periodic, the difference
$\zeta(z+\omega_j)-\zeta(z)$ has derivative zero. Its simple poles
and residues cancel, so the difference is a constant $\eta_j$.
The ratio
$\sigma(z+\omega_j)/(\sigma(z)e^{\eta_jz})$ has all its zeros
and poles removed and logarithmic derivative zero. It is a nonzero
constant. At $z=-\omega_j/2$, oddness makes that constant
$-e^{\eta_j\omega_j/2}$, proving sigma quasi-periodicity.
Integrate $\zeta$ around a positive translated fundamental
parallelogram avoiding its poles. It encloses one residue-one pole.
Opposite-side subtraction gives the integral
$\eta_1\omega_2-\eta_2\omega_1$, proving Legendre's relation.

These formulas determine arbitrary lattice shifts, without an
uncontrolled integration period. Explicitly, if $\lambda=m\omega_1+n\omega_2$
and $\eta_\lambda=m\eta_1+n\eta_2$, then
\[
 \sigma(z+\lambda)=(-1)^{m+n+mn}
             e^{\eta_\lambda(z+\lambda/2)}\sigma(z).
\]
Iterating a single generator for positive integers and then solving
the same formula backwards gives the factor
$(-1)^m e^{m\eta_jz+m^2\eta_j\omega_j/2}$ for every integer $m$.
Apply first the $\omega_1$ and then the $\omega_2$ shift. The difference
from the symmetric exponent above is
$mn(\eta_2\omega_1-\eta_1\omega_2)/2=-\pi\iu mn$;
its exponential supplies the extra sign $(-1)^{mn}$.
\end{proof}

\begin{theorem}[Addition and the chord--tangent law]\label{d12:22-09}\family{D12.22.09}
Transporting addition on $T$ through $\Phi$ gives exactly the
chord--tangent group law on $E$, with identity $O$ and inversion
$[Z:X:Y]\mapsto[Z:X:-Y]$. In particular, when the denominator is
nonzero,
\[
 \wp(z+w)=-\wp(z)-\wp(w)
       +\frac14\left(\frac{\wp'(z)-\wp'(w)}{\wp(z)-\wp(w)}\right)^2.
\]
The equality is a meromorphic identity; the group description includes
diagonal, vertical, half-period and infinity cases.
\end{theorem}
\begin{proof}
For distinct finite points $P=\Phi(z)$, $Q=\Phi(w)$ with different
$x$-coordinates, let $y=mx+b$ be their line. Its pullback
$\wp'(t)-m\wp(t)-b$ has a triple pole at zero and no other poles.
It has three zeros on $T$, counted with multiplicity, and their sum
is zero in $T$ by \cref{d12:22-03}. Two zeros are $z,w$; call the
third $v$. Thus $v=-z-w$, and reflection of the third intersection
$\Phi(v)$ is $\Phi(z+w)$. Substitution of $y=mx+b$ in the cubic
gives a cubic in $x$ whose leading coefficient is $4$ and whose
$x^2$ coefficient is $-m^2$. Therefore its three intersection
$x$-coordinates sum to $m^2/4$. Since $x(\Phi(v))=\wp(z+w)$,
this is the stated addition formula. The reflected ordinate is
$-y_P+m(x_P-x_{P+Q})$, giving the ordinate formula too.

For $P=Q$ with $y_P\ne0$, the tangent slope is
$m=(12x_P^2-g_2)/(2y_P)=\wp''(z)/\wp'(z)$.
The same pullback has a zero of order at least two at $z$; its total
zero order is three. Its remaining zero $v$ satisfies $2z+v=0$,
including order three at a flex. This proves the tangent rule and
the doubling formula with this slope.
If finite points have equal $x$-coordinates, they are equal or
opposites. The vertical line $X=x_PZ$ has intersection equation
$Z(Y^2-y_P^2Z^2)=0$, so its third intersection is $O$;
for $y_P=0$ it is the tangent with double intersection at $P$.
Thus opposites sum to $O$, and a half-period point doubles to $O$.
A line through $O$ and a finite $P$ is this vertical line, so its
third intersection is $-P$ and $O+P=P$. Finally the tangent at $O$
is $Z=0$; substitution gives intersection multiplicity three there,
so $O+O=O$. These cases exhaust all pairs of points. The operation
defined by chords and tangents is consequently the transported torus
addition; its associativity, commutativity and inverse laws follow
from the actual abelian group $\C/\Lambda$, rather than being
assumed from an unproved algebraic-curve group theorem.
\end{proof}
\source{AHL VII.2.4--3.3, printed pp.270--277 (PDF pp.286--293),
compared directly for residue/divisor constraints, regularized sums,
periodicity, zeta/sigma, Legendre and the differential equation.
FB V.3 Appendix A.6--A.7, printed pp.276--277 (PDF pp.281--282),
states the projective bijection; the nonsingularity, infinity chart,
analytic inverse and complete group-law bridge have been supplied here.}

% END chapters/22

% BEGIN chapters/23
\chapter{Bloch, Schottky and Picard}\label{ch:d12-23}

This chapter proves an omitted-values normality theorem, rather than
assuming it as the locally bounded Montel theorem. The numerical constant
below is a deliberately nonsharp proved constant.

\begin{theorem}[A quantitative Bloch theorem]\label{d12:23-01}\family{D12.23.01}
If $f$ is holomorphic on $\D$ and $f'(0)=1$, some disk $S\subset\D$
is a univalence disk for $f$, and $f(S)$ contains a disk of radius
$1/64$. More generally, on $B(a,R)$ its image contains a disk of
radius $R|f'(a)|/64$ when $f'(a)\ne0$, with a single-valued inverse
on that disk coming from a univalence disk.
\end{theorem}
\begin{proof}
On $|z|\le r=1/2$, maximize
$M=(r-|z|)|f'(z)|$. It is continuous, zero on the boundary and at
least $1/2$ at zero. A maximizing point $b$ is interior and
$\delta=(r-|b|)/2>0$. For $|z-b|<\delta$,
$|f'(z)|\le M/(r-|z|)<M/\delta=2|f'(b)|$.
The function $h(z)=f'(z)/f'(b)-1$ vanishes at $b$ and has modulus
less than three on that disk. Schwarz's lemma after rescaling gives
$|h(z)|\le3|z-b|/\delta$. On the smaller disk $S=B(b,\delta/6)$
it is at most $1/2$. For distinct $z,w\in S$, integration along
their segment gives
\[
 |f(z)-f(w)-f'(b)(z-w)|\le\tfrac12|f'(b)||z-w|,
\]
so $f$ is injective there. On its boundary, the same estimate from
$b$ gives a remainder of modulus at most
$|f'(b)|\delta/12$ after subtraction of the linear term.
For $|v|<|f'(b)|\delta/12$, Rouch\'e applied to
$f(z)-f(b)-v$ and $f'(b)(z-b)$ gives one zero inside $S$.
Thus its image contains the disk with this radius, which equals
$M/24\ge1/48>1/64$. All disks used lie strictly inside $\D$;
no extension to its boundary was assumed. Apply this to
$(f(a+Rz)-f(a))/(Rf'(a))$ to obtain the general assertion.
\end{proof}

\begin{lemma}[An omitted lattice of auxiliary values]\label{d12:23:auxiliary}
If a holomorphic $f$ on a simply connected domain omits zero and one,
there is a holomorphic $g$ there such that
\[
                      f=-\exp(\pi\iu\cosh(2g)).
\]
The image of $g$ contains no disk of radius two.
\end{lemma}
\begin{proof}
Take a holomorphic logarithm $t$ of $f$ and put $F=t/(2\pi\iu)$.
It omits every integer, since an integral value would give $f=1$.
Choose holomorphic square roots $A^2=F$, $B^2=F-1$ and let $H=A-B$.
Then $H(A+B)=1$, so $H$ is nonzero and has a holomorphic logarithm
$g$. These logarithms and roots exist by \cref{d12:04-07}.
The identity
$\cosh(2g)=(H^2+H^{-2})/2=2F-1$ proves the formula for $f$.

Put $a_n=\log(\sqrt n+\sqrt{n-1})$, $n\ge1$. The function $g$
omits each $\pm a_n+\pi\iu k/2$, $k\in\Z$: at such a point
$\cosh(2g)=(-1)^k(2n-1)$, making $F$ an integer and $f=1$.
The numbers $a_n$ increase from zero to infinity and their successive
gaps are at most $a_2<1$. Indeed the derivative of
$\log(\sqrt x+\sqrt{x-1})$ is
$1/(2\sqrt{x(x-1)})$, a decreasing positive function for $x>1$;
integrating it over consecutive unit intervals proves the gap bound.
Also $1+\sqrt2<5/2<e$, proving $a_2<1$.
The omitted points are the vertices of rectangles of width less than
one and height $\pi/2<2$. Every point of a rectangle is at distance
at most half its diagonal, hence less than $\sqrt5/2<2$, from a
vertex. Thus every disk of radius two contains an omitted vertex.
The elementary bound $\pi<4$ used here follows, for example, from
$\pi/4=\int_0^1(1+x^2)^{-1}dx<1$. This proves the stated, deliberately
weaker-than-optimal, excluded-disk bound.
\end{proof}

\begin{theorem}[Little Picard]\label{d12:23-02}\family{D12.23.02}
A nonconstant entire function omits at most one finite complex value.
\end{theorem}
\begin{proof}
An affine change of the range sends two putatively omitted values
to zero and one. The auxiliary lemma produces an entire $g$. If $f$
is nonconstant, $g$ is nonconstant and $g'(a)\ne0$ at some point.
Bloch on $B(a,R)$ gives a disk of radius $R|g'(a)|/64$ in its image.
For sufficiently large $R$ this radius exceeds two, contradicting
the auxiliary lemma. Therefore a function omitting two values is
constant. This argument has not used finite order or normality of
omitted-value families.
\end{proof}

\begin{theorem}[Schottky control]\label{d12:23-03}\family{D12.23.03}
For $A>0$ and $0\le r<1$ there is an explicit finite constant
$C(A,r)$ such that every holomorphic $f:\D\to\C\setminus\{0,1\}$
with $|f(0)|\le A$ satisfies $|f(z)|\le C(A,r)$ on $|z|\le r$.
One choice is, with $A_*=\max(2,A)$,
\[
 C_0=1+\frac{\log A_*}{2\pi},\quad K=\sqrt{C_0}+\sqrt{C_0+1},
 \quad C(A,r)=1+\exp\left(\pi\exp\left[
             2\left(\log K+\pi+\frac{128r}{1-r}\right)\right]\right).
\]
Scaling gives the corresponding bound on every proper subdisk of any
disk where the same omitted values and center bound hold.
\end{theorem}
\begin{proof}
First assume $1/2\le|f(0)|\le A_*$. In the auxiliary construction
choose $|\Imm t(0)|\le\pi$ and $|\Imm g(0)|\le\pi$.
Then $|F(0)|\le C_0$. Both
$H(0)=A(0)-B(0)$ and $H(0)^{-1}=A(0)+B(0)$ have modulus at most
$K$, whence $|g(0)|\le\log K+\pi$.
At $a\in\D$, apply Bloch to $g$ on $B(a,1-|a|)$. Because its image
has no disk of radius two, the resulting radius cannot exceed two:
$|g'(a)|\le128/(1-|a|)$. This also holds where the derivative is zero.
Integrating along the segment from zero to $z$, $|z|\le r$, gives
$|g(z)|\le\log K+\pi+128r/(1-r)$.
Finally $|\cosh(2g)|\le e^{2|g|}$ and
$|\exp(\pi\iu\cosh(2g))|\le\exp(\pi e^{2|g|})$, proving the
bound without its harmless extra one. If $0<|f(0)|<1/2$, apply this
case to $1-f$, which also omits zero and one and has center modulus
between $1/2$ and $3/2\le A_*$. Then
$|f|\le1+|1-f|$ gives the displayed constant. The change of variables
$z\mapsto a+Rz$ proves scaling. The restriction $r<1$ is essential;
no uniform bound at the disk boundary has been asserted.
\end{proof}

\begin{theorem}[Montel--Carath\'eodory]\label{d12:23-04}\family{D12.23.04}
On a plane domain $\Omega$, holomorphic functions omitting zero and
one form a spherically normal family. Every sequence has a subsequence
converging compactly in the spherical metric to a holomorphic function
or to the constant infinity. The finite limit may be the constant zero
or one.
\end{theorem}
\begin{proof}
Fix $a\in\Omega$. The subfamily with $|f(a)|\le1$ is locally bounded
throughout $\Omega$. To see this, join $a$ to any point by a polygonal
path and choose a finite chain of disks compactly contained in $\Omega$,
with each next center in a proper concentric subdisk of the preceding
one. Schottky bounds the first subdisk uniformly, hence the value at
the next center; repeating its bound through this finite chain gives
a uniform bound on a neighborhood of the final point. A compact set
has a finite cover by such neighborhoods, giving its uniform bound.
The locally bounded Montel theorem \cref{d12:11-03} therefore applies.

For an arbitrary sequence, infinitely many center values are at most
one, or infinitely many exceed one. In the first case use the preceding
finite-limit normality. In the second apply it to $1/f$, which is
holomorphic and omits zero and one, with center values less than one.
A subsequence converges ordinarily to a holomorphic $h$. Hurwitz's
zero-free alternative \cref{d12:09-05} says $h$ is nonvanishing or
identically zero. If it is nonvanishing, its minimum modulus on every
compact set is positive, so inversion gives ordinary compact convergence
of $f$ to $1/h$. If $h=0$, the spherical distance of $f$ to infinity
is $|1/f|/\sqrt{1+|1/f|^2}$ and tends to zero uniformly on compact
sets. This proves spherical normality. Constants approaching infinity
show why finite local boundedness is a different conclusion.
\end{proof}

\begin{theorem}[Great Picard]\label{d12:23-05}\family{D12.23.05}
If a holomorphic function has an isolated essential singularity at
$a$, then in every punctured neighborhood of $a$ it takes every finite
value, with at most one exception, infinitely many times. More precisely,
there is at most one finite value whose preimage is finite in some
punctured neighborhood of $a$.
\end{theorem}
\begin{proof}
Translate $a$ to zero. Suppose first that $f$ omits two values on
$0<|z|<R$. Normalize them to zero and one, and put $f_n(z)=f(z/n)$
on that punctured disk. Omitted-value normality gives a subsequence
converging spherically to a holomorphic $h$ or to infinity. On the
fixed circle $|z|=R/2$, convergence to a finite $h$ is ordinary:
the compact range of $h$ stays a positive spherical distance from
infinity, and in this finite chart the spherical and ordinary metrics
are uniformly equivalent on a fixed bounded set. Thus $|f|\le M$
on circles of radii $r_k=R/(2n_k)\downarrow0$, with one fixed finite
$M$. The maximum principle on each closed annulus
$r_{k+1}\le|z|\le r_k$ gives the same bound throughout it. Hence
$f$ is bounded on a punctured disk and its singularity is removable.

If the limit is infinity, apply the same circle and annulus argument
to the holomorphic function $1/f$. Its values tend uniformly to zero
on the selected circles. The maximum principle makes it bounded
near zero, so it extends holomorphically there. It takes value zero
at zero, because the circle suprema tend to zero. It is not identically
zero, since it was nonzero on the punctured disk. Its zero has finite
positive order, so $f$ has a pole. Both alternatives exclude an
essential singularity. The reciprocal argument supplies the step
sometimes left as an exercise in this proof route.

Finally, if two distinct finite values each have only finitely many
preimages in some punctured neighborhood, choose a common smaller
radius. Remove the finitely many remaining preimage points by reducing
the radius below their positive distances from zero. The two values
are now omitted, contradicting the preceding conclusion. Thus the
set of values with finitely many preimages near the singularity has
at most one member. Every other value has infinitely many preimages
in every punctured neighborhood, exactly the stated quantifiers.
\end{proof}

\begin{theorem}[The sphere-valued versions]\label{d12:23-06}\family{D12.23.06}
Meromorphic functions on a domain omitting three fixed distinct sphere
values form a spherically normal family. A meromorphic function on
$\C$ omitting three sphere values is constant. At an isolated essential
singularity of a meromorphic map to $\Chat$, all sphere values with
at most two exceptions are taken infinitely often in every punctured
neighborhood.
\end{theorem}
\begin{proof}
Choose one fixed M\"obius map sending the omitted values to
$0,1,\infty$. Its composition with each function has no poles, so
is holomorphic and omits zero and one. The preceding normality theorem
gives a compact spherical subsequence. The inverse M\"obius map is
uniformly continuous on the compact sphere; composing the limit gives
spherical compact convergence to a meromorphic map or a sphere constant.
On the whole plane, Little Picard makes the transformed entire function
constant and hence makes the original constant.
If a meromorphic map omitted three values near its puncture, the same
transformation would give the holomorphic omitted-two-values function
treated in the Great Picard proof. That function extends holomorphically
or with a pole, hence as a holomorphic map to the sphere. Composing
with the inverse M\"obius map extends the original sphere map, which
contradicts essentiality. Three values each having only finitely many
preimages could be made omitted in a smaller punctured disk by deleting
their finitely many preimages. Thus at most two sphere values have
that finite-preimage property. Poles are counted as preimages of
infinity. When the original function is holomorphic, infinity is
already one omitted sphere value, leaving at most one finite exception.
\end{proof}
\source{CON XII.1--4, printed pp.292--301 (PDF pp.306--315), directly
checked for the derivative-disk Bloch argument, the auxiliary cosh
transformation, Schottky propagation and the two normal-limit cases.
The constants here are independently proved nonsharp constants. The
infinite-limit and finite-preimage steps are fully supplied; neither
bounded Montel nor an unproved modular uniformization stands in for
the omitted-value theorem.}

% END chapters/23

\backmatter
% BEGIN source-register
\chapter*{Source register}\addcontentsline{toc}{chapter}{Source register}
The source files, editions, sizes and SHA256 identities are retained in
\texttt{TEXTBOOK\_INVENTORY.json}. The original PDFs remain separate from the
MathGenome release; possession of a user-supplied file is not a redistribution
license. Source comparison abbreviations are:
\begin{description}
\item[CON] J. B. Conway, \emph{Functions of One Complex Variable I}, second edition.
\item[AHL] L. V. Ahlfors, \emph{Complex Analysis}, third edition, 1979.
\item[SS] E. M. Stein and R. Shakarchi, \emph{Complex Analysis}, 2003.
\item[FB] E. Freitag and R. Busam, \emph{Complex Analysis}, second English edition, 2009.
\item[FOR] O. Forster, \emph{Lectures on Riemann Surfaces}, English edition 1981,
corrected fourth printing 1999; only the explicitly selected introductory scope.
\end{description}
Printed page numbers and physical PDF page numbers are different. The source
crosswalk records both when inspected; no universal offset is assumed across
all books. Differences in analytic/holomorphic definitions, spherical limits,
and boundary hypotheses are resolved by explicit statements and bridges.

% END source-register

% BEGIN verification-register
\chapter*{Verification register}\addcontentsline{toc}{chapter}{Verification register}
The adopted scope comprises 22 mathematical chapters and 151 stable content
families. The three author self-checks enumerate all constituent statements,
source comparisons, hypotheses, proof dependencies and boundary cases.
These author reports are not themselves independent acceptance.

Three non-author cross-reviews read every assigned statement and proof:
\begin{center}
\begin{tabular}{lll}\toprule
Stable chapter IDs & Adopted families & Independent reviewer\\\midrule
01--08 & 52 & Author of 09--16\\
09--16 & 55 & Author of 17, 19--23\\
17, 19--23 & 44 & Author of 01--08\\\bottomrule
\end{tabular}
\end{center}
All three completed their full assigned scope, including necessary proof
supports, and recorded no unresolved blocker after repair. The exact final
chapter bytes are bound by SHA256 in the audit reports. This is independent
cross-review, not a claim of a new blind review or of infallibility.

The resolved findings include the punctured-limit sequence restriction,
the explicit finite-variation premise for Stieltjes sums, the general-covering
homotopy subdivision and locally constant-germ propagation arguments, and the
explicit C1-rounding bridge needed before using the prior Green identity.
A source-page locator was also clarified. None of these repairs removes an
adopted principal conclusion. Additional checks during authoring distinguished
given regular regions from arbitrary Jordan traces and supplied the
pre-limit barrier comparison for canonical Green functions.

The coordinating reader independently inspected the difficult Goursat/full-cycle
Cauchy closure, BV/AC contour bridges, finite versus spherical normality,
extremal Riemann mapping, Perron/Green arguments, Runge and general-domain
products, finite-order factorization, the torus/cubic infinity chart and all
group-law cases, and the complete Bloch--Schottky--Picard route. A separate
integration audit reconciles the ledger, scope, prior interfaces, cross-chapter
proof architecture and exact review coverage.

The PDF is rebuilt from the reviewed modular source. Its standalone LaTeX
export, source files, principal-result catalog, bibliography identities,
review reports and PDF hashes are recorded in the accompanying
\texttt{FORMALIZATION\_INPUTS.json}. Compilation and visual inspection verify
document integrity and layout; they are not mathematical proof certificates.
The prior MathGenome v0.6.0 interfaces are reuse candidates whose actual
hypotheses must still match each future Lean use. No future D22, D23 or D24
theorem is an implicit assumption.

All adopted principal results retain their hypotheses and conclusions.
Valid equivalent proof routes are allowed; unused intermediate steps from a
particular source proof are not additional obligations. The expressly deferred
graduate and application branches are not represented as proved here.
No D12 Lean proof or formalization completion is claimed by this preparation.

% END verification-register

\end{document}
